% Numerical Methods — Pure Mathematics with Statistics Upper Sixth — Cours signé OnBuch+
% Official MINESEC syllabus. Compiler avec : tectonic PMS06-numerical-methods.tex
\newcommand{\DOCMATIERE}{Pure Mathematics with Statistics}
\newcommand{\DOCNIVEAU}{Upper Sixth}
\newcommand{\DOCMODULE}{Upper Sixth — Pure mathematics with statistics}
\newcommand{\DOCLECON}{6}
\newcommand{\DOCTITRE}{Numerical Methods}
% OnBuch+ — préambule « cours signés » ENRICHI (compilé avec tectonic / XeLaTeX)
% Système visuel « Pop » d'OnBuch+ : fond crème (jamais blanc), bordures encre
% épaisses, ombres « dures » décalées, titres Archivo Black en capitales.
% Version MATHÉMATIQUES : graphes (pgfplots/tikz), boîtes pédagogiques
% (définition, propriété, méthode, attention, le savais-tu, exercice résolu,
% exercice, TP, bilan), page de garde et signature OnBuch+ ancrée en bas.
%
% Macros attendues dans main.tex AVANT \input{preamble} :
%   \DOCMATIERE  \DOCNIVEAU  \DOCMODULE  \DOCLECON (numéro)  \DOCTITRE
\documentclass[11pt]{article}

\usepackage{fontspec}
\usepackage[english]{babel}
\usepackage[a4paper,margin=2cm,top=3.4cm,bottom=2.8cm,headheight=1.4cm,footskip=1.3cm]{geometry}
\usepackage{amsmath,amssymb,amsfonts}
\usepackage{mathtools} % \xmapsto, \coloneqq... souvent utilisés par les modèles
\usepackage{cancel} % simplifications barrées (bilans de forces)
% \abs{z} (module, valeur absolue) et \norm{u} (norme), utilisés par les modèles
\providecommand{\abs}{}\let\abs\relax\DeclarePairedDelimiter{\abs}{\lvert}{\rvert}
\providecommand{\norm}{}\let\norm\relax\DeclarePairedDelimiter{\norm}{\lVert}{\rVert}
\usepackage[table]{xcolor} % [table] : fournit aussi \rowcolors (alternance de lignes)
\usepackage{pagecolor}
\usepackage{tikz}
\usetikzlibrary{arrows.meta,positioning,calc,shapes.geometric,shapes.misc,shapes.arrows,decorations.pathmorphing,decorations.pathreplacing,decorations.markings,patterns,shadows,mindmap,trees,fit,backgrounds,matrix,intersections,angles,quotes}
\usepackage{pgfplots}
\usepackage[version=4]{mhchem} % équations nucléaires \ce{^{235}_{92}U}
\usepackage[european,siunitx]{circuitikz} % schémas électriques
\pgfplotsset{compat=1.18}
\usepgfplotslibrary{fillbetween,groupplots} % aires entre courbes (calcul intégral)
\usepackage{enumitem}
\usepackage{fancyhdr}
% [most] charge toutes les bibliothèques tcolorbox (listings, minted, raster,
% documentation...) et gonfle inutilement la mémoire TeX sur un document long
% (~300 boîtes) : on ne charge que ce qui est réellement utilisé.
\usepackage[skins,breakable]{tcolorbox}
\usepackage{graphicx}
\usepackage{colortbl}
\usepackage{booktabs}
\usepackage{tabularx}
\usepackage{diagbox} % case d'en-tête diagonale des tableaux à double entrée
% Colonnes extensibles centrée / gauche / droite (C, L, R), souvent utilisées par les modèles
\newcolumntype{C}{>{\centering\arraybackslash}X}
\newcolumntype{L}{>{\raggedright\arraybackslash}X}
\newcolumntype{R}{>{\raggedleft\arraybackslash}X}
\usepackage{array}
\usepackage{multicol}
\usepackage{siunitx}
% parse-numbers=false : accepte aussi \SI{2,5 \times 10^{-3}}{...}
\sisetup{locale=UK,output-decimal-marker={.},group-minimum-digits=4,parse-numbers=false}
\DeclareSIUnit\ohm{\ensuremath{\symup{\Omega}}} % U+2126 absent de la police
\DeclareSIUnit\division{div} % réglages d'oscilloscope (ms/div, V/div)
\DeclareSIUnit\tr{tr} % tours (vitesse de rotation, ex. \SI{1500}{\tr\per\minute})
\usepackage{qrcode}
% \degree : commande fréquemment utilisée par les modèles pour un angle ou une
% température en dehors de \SI{}{} (non fournie par les paquets chargés).
\providecommand{\degree}{^{\circ}}
% \qed (fin de démonstration), utilisé par les modèles sans amsthm
\providecommand{\qed}{\hfill$\square$}
% \barre : double barre des valeurs interdites dans un tableau de variations (tkz-tab)
\providecommand{\barre}{\ensuremath{\|}}
% environnement proof (amsthm non chargé) utilisé par les modèles
\ifdefined\proof\else\newenvironment{proof}[1][Proof]{\par\noindent\textit{#1.}\ }{\qed\par}\fi
\usepackage{lastpage}
\usepackage{hyperref}
\hypersetup{hidelinks}

% ============================================================
% Palette « Pop » OnBuch+ — mêmes hex que la classe P (plus_ui.dart)
% ============================================================
\definecolor{poporange}{HTML}{F59321}
\definecolor{popdark}{HTML}{E07A0C}
\definecolor{popcream}{HTML}{FFF6E8}
\definecolor{popink}{HTML}{1C1714}
\definecolor{poppurple}{HTML}{7A5AE0}
\definecolor{popgreen}{HTML}{2E9E6A}
\definecolor{poppink}{HTML}{E0457B}
\definecolor{popblue}{HTML}{2F7DE1}
\definecolor{popgold}{HTML}{C98A12}
\definecolor{popmuted}{HTML}{8A7F76}
\definecolor{popline}{HTML}{EADFD2}
\definecolor{popgray}{HTML}{8A7F76}
\colorlet{popgrey}{popgray}
% Alias tolérants (le modèle nomme parfois les couleurs différemment)
\colorlet{popwhite}{white}
\colorlet{popblack}{popink}
\colorlet{popred}{poppink}
\colorlet{popyellow}{popgold}
% Teintes claires pour fonds de tableaux / zones
\colorlet{popblueL}{popblue!10!white}
\colorlet{poporangeL}{poporange!14!white}
\colorlet{popgreenL}{popgreen!12!white}
\colorlet{poppurpleL}{poppurple!12!white}
\colorlet{poppinkL}{poppink!12!white}
\colorlet{popgoldL}{popgold!14!white}

\pagecolor{popcream}

% ============================================================
% Typographie
% ============================================================
\defaultfontfeatures{Ligatures=TeX}
\setmainfont{PlusJakartaSans-Medium.ttf}[Path=../../../fonts/,BoldFont=PlusJakartaSans-Bold.ttf]
% Maths en sans-serif (Fira Math) avec chiffres et lettres droites de Plus
% Jakarta, pour que les formules se fondent dans le texte.
\usepackage[math-style=ISO,bold-style=ISO]{unicode-math}
\setmathfont{FiraMath-Regular.otf}[Path=../../../fonts/]
\setmathfont{PlusJakartaSans-Medium.ttf}[Path=../../../fonts/,range={up/num,up/{Latin,latin},\mathup}]
% (icomma retiré : décimales avec point en anglais)
% Caractères mathématiques Unicode absents des polices de texte (Plus Jakarta,
% Archivo Black) : rendus par leur équivalent mathématique, en texte comme en
% formule — sinon ils s'affichent en carré vide (ex. « Arithmétique dans ℤ »).
\usepackage{newunicodechar}
\newunicodechar{ℝ}{\ensuremath{\mathbb{R}}}
\newunicodechar{ℤ}{\ensuremath{\mathbb{Z}}}
\newunicodechar{ℂ}{\ensuremath{\mathbb{C}}}
\newunicodechar{ℕ}{\ensuremath{\mathbb{N}}}
\newunicodechar{ℚ}{\ensuremath{\mathbb{Q}}}
\newunicodechar{𝔻}{\ensuremath{\mathbb{D}}}
\newunicodechar{α}{\ensuremath{\alpha}}
\newunicodechar{β}{\ensuremath{\beta}}
\newunicodechar{γ}{\ensuremath{\gamma}}
\newunicodechar{δ}{\ensuremath{\delta}}
\newunicodechar{ε}{\ensuremath{\varepsilon}}
\newunicodechar{θ}{\ensuremath{\theta}}
\newunicodechar{λ}{\ensuremath{\lambda}}
\newunicodechar{μ}{\ensuremath{\mu}}
\newunicodechar{σ}{\ensuremath{\sigma}}
\newunicodechar{τ}{\ensuremath{\tau}}
\newunicodechar{φ}{\ensuremath{\varphi}}
\newunicodechar{ψ}{\ensuremath{\psi}}
\newunicodechar{ω}{\ensuremath{\omega}}
\newunicodechar{Σ}{\ensuremath{\Sigma}}
% majuscules produites par \MakeUppercase dans les titres (α-aminés -> Α)
\newunicodechar{Α}{\ensuremath{\alpha}}
\newunicodechar{Β}{\ensuremath{\beta}}
\newunicodechar{Γ}{\ensuremath{\Gamma}}
\newunicodechar{Δ}{\ensuremath{\Delta}}
\newunicodechar{Ε}{\ensuremath{\varepsilon}}
\newunicodechar{Θ}{\ensuremath{\Theta}}
\newunicodechar{Λ}{\ensuremath{\Lambda}}
\newunicodechar{Μ}{\ensuremath{\mu}}
\newunicodechar{Τ}{\ensuremath{\tau}}
\newunicodechar{Φ}{\ensuremath{\Phi}}
\newunicodechar{Ψ}{\ensuremath{\Psi}}
\newunicodechar{Ω}{\ensuremath{\Omega}}
\newunicodechar{↦}{\ensuremath{\mapsto}}
\newunicodechar{∈}{\ensuremath{\in}}
\newunicodechar{∉}{\ensuremath{\notin}}
\newunicodechar{∧}{\ensuremath{\wedge}}
\newunicodechar{⇔}{\ensuremath{\Leftrightarrow}}
\newunicodechar{⟺}{\ensuremath{\Longleftrightarrow}}
\newunicodechar{⇒}{\ensuremath{\Rightarrow}}
\newunicodechar{⟹}{\ensuremath{\Longrightarrow}}
\newunicodechar{∘}{\ensuremath{\circ}}
\newunicodechar{∪}{\ensuremath{\cup}}
\newunicodechar{∩}{\ensuremath{\cap}}
\newunicodechar{≡}{\ensuremath{\equiv}}
\newunicodechar{⊂}{\ensuremath{\subset}}
\newunicodechar{⊥}{\ensuremath{\perp}}
\newunicodechar{∃}{\ensuremath{\exists}}
\newunicodechar{∀}{\ensuremath{\forall}}
\newunicodechar{∛}{\ensuremath{\sqrt[3]{\,}}}
\newunicodechar{∜}{\ensuremath{\sqrt[4]{\,}}}
\newunicodechar{⃗}{\ensuremath{{}^{\to}}}
\newunicodechar{ⁿ}{\ifmmode^{n}\else\textsuperscript{\ensuremath{n}}\fi}
\newunicodechar{ˣ}{\ifmmode^{x}\else\textsuperscript{\ensuremath{x}}\fi}
\newunicodechar{ᵇ}{\ifmmode^{b}\else\textsuperscript{\ensuremath{b}}\fi}
\newunicodechar{ʳ}{\ifmmode^{r}\else\textsuperscript{\ensuremath{r}}\fi}
\newunicodechar{ᵘ}{\ifmmode^{u}\else\textsuperscript{\ensuremath{u}}\fi}
\newunicodechar{ʸ}{\ifmmode^{y}\else\textsuperscript{\ensuremath{y}}\fi}
\newunicodechar{ᵃ}{\ifmmode^{a}\else\textsuperscript{\ensuremath{a}}\fi}
\newunicodechar{ᵉ}{\ifmmode^{e}\else\textsuperscript{\ensuremath{e}}\fi}
\newunicodechar{ᵏ}{\ifmmode^{k}\else\textsuperscript{\ensuremath{k}}\fi}
\newunicodechar{ᵖ}{\ifmmode^{p}\else\textsuperscript{\ensuremath{p}}\fi}
\newunicodechar{ᴺ}{\ifmmode^{N}\else\textsuperscript{\ensuremath{N}}\fi}
\newunicodechar{ᶜ}{\ifmmode^{c}\else\textsuperscript{\ensuremath{c}}\fi}
\newunicodechar{ʰ}{\ifmmode^{h}\else\textsuperscript{\ensuremath{h}}\fi}
\newunicodechar{ᵠ}{\ifmmode^{q}\else\textsuperscript{\ensuremath{q}}\fi}
\newunicodechar{ₙ}{\ifmmode_{n}\else\textsubscript{\ensuremath{n}}\fi}
\newunicodechar{ₐ}{\ifmmode_{a}\else\textsubscript{\ensuremath{a}}\fi}
\newunicodechar{ᵢ}{\ifmmode_{i}\else\textsubscript{\ensuremath{i}}\fi}
\newunicodechar{ᵣ}{\ifmmode_{r}\else\textsubscript{\ensuremath{r}}\fi}
\newunicodechar{ₖ}{\ifmmode_{k}\else\textsubscript{\ensuremath{k}}\fi}
\newunicodechar{ᵦ}{\ifmmode_{\beta}\else\textsubscript{\ensuremath{\beta}}\fi}
\newunicodechar{ₓ}{\ifmmode_{x}\else\textsubscript{\ensuremath{x}}\fi}
\newfontfamily\popdisplay{ArchivoBlack-Regular.ttf}[Path=../../../fonts/]
\newfontfamily\poplogo{SpaceGrotesk-Bold.ttf}[Path=../../../fonts/]
\newfontfamily\popsemi{PlusJakartaSans-SemiBold.ttf}[Path=../../../fonts/]
\newfontfamily\popxbold{PlusJakartaSans-ExtraBold.ttf}[Path=../../../fonts/]
\newfontfamily\popxboldls{PlusJakartaSans-ExtraBold.ttf}[Path=../../../fonts/,LetterSpace=3.0]

\setlist[enumerate]{leftmargin=1.7em,itemsep=5pt,topsep=4pt}
\setlist[itemize]{leftmargin=1.7em,itemsep=4pt,topsep=4pt,label={\color{poporange}\textbullet}}
\setlength{\parindent}{0pt}
\setlength{\parskip}{5pt}
\renewcommand{\arraystretch}{1.25}

% pgfplots : style Pop par défaut
\pgfplotsset{
  every axis/.append style={axis line style={popink,line width=1pt},tick style={popink},
    grid style={popline,line width=0.5pt},label style={font=\small},tick label style={font=\footnotesize}},
  popcurve/.style={poporange,line width=1.8pt,no markers,smooth},
}
% Même style disponible dans un \draw TikZ simple (hors pgfplots)
\tikzset{popcurve/.style={poporange,line width=1.8pt}}
\tikzset{popgrid/.style={popline,very thin}}
\colorlet{popcurve}{poporange} % le nom du style est parfois utilisé comme couleur

% ============================================================
% Pill / squiggle
% ============================================================
\newcommand{\poppill}[3]{%
  \tikz[baseline=-0.55ex]\node[draw=popink,line width=1.2pt,rounded corners=2.6pt,%
    fill=#2,inner xsep=6.5pt,inner ysep=3pt]{%
    \popxboldls\fontsize{7}{7}\selectfont\color{#3}\MakeUppercase{#1}};%
}
\newcommand{\popsquiggle}[1]{%
  \tikz[baseline=-0.4ex]{
    \draw[#1,line width=2.6pt,line cap=round]
      plot [smooth] coordinates {(0,0.09) (0.34,0.01) (0.68,0.21) (1.02,0.01) (1.36,0.10)};
  }%
}
% Mot-clé surligné (marqueur orange sous le texte)
\newcommand{\cle}[1]{{\popxbold\color{popdark}#1}}

% ============================================================
% En-tête / pied
% ============================================================
\pagestyle{fancy}
\fancyhf{}
\renewcommand{\headrulewidth}{0pt}
\renewcommand{\footrulewidth}{0pt}
\lhead{\raisebox{-0.15ex}{\poplogo\fontsize{13}{13}\selectfont\color{popink}OnBuch\color{poporange}+}}
\rhead{\poppill{\DOCMATIERE\ \textbullet\ \DOCNIVEAU\ \textbullet\ Chapter \DOCLECON}{white}{popink}}
\lfoot{\popxbold\fontsize{7.6}{7.6}\selectfont\color{popmuted}\MakeUppercase{Signed course OnBuch+}\ \textbullet\ {\popsemi\DOCTITRE}}
\rfoot{\tikz[baseline=-0.6ex]\node[rounded corners=8.5pt,fill=popink,minimum height=17pt,minimum width=17pt,inner xsep=5pt,inner sep=0]{\popxbold\fontsize{8}{8}\selectfont\color{popcream}\thepage\,/\,\pageref*{LastPage}};}

% ============================================================
% Titres
% ============================================================
\newcommand{\chaptitre}[2]{%
  \begin{tcolorbox}[enhanced,colback=poporange,colframe=popink,boxrule=2.6pt,arc=11pt,
      drop shadow=popink,left=15pt,right=15pt,top=12pt,bottom=11pt,before skip=3pt,after skip=18pt]
    \poppill{#2}{popink}{popcream}\par\vspace{9pt}
    {\raggedright\hyphenpenalty=10000\exhyphenpenalty=10000\popdisplay\fontsize{22}{24}\selectfont\color{popcream}\MakeUppercase{#1}\par}
  \end{tcolorbox}
}
% Section numérotée : pastille orange avec le numéro + titre Archivo
\newcounter{popsec}
\newcommand{\coursec}[1]{%
  \stepcounter{popsec}\setcounter{popsub}{0}%
  \par\vspace{16pt}\noindent
  \begin{minipage}{\linewidth}
  \tikz[baseline=-0.7ex]\node[draw=popink,line width=1.6pt,fill=poporange,rounded corners=4pt,minimum size=22pt,inner sep=0]{\popdisplay\fontsize{12}{12}\selectfont\color{popcream}\thepopsec};%
  \hspace{8pt}{\popdisplay\fontsize{15}{17}\selectfont\color{popink}\MakeUppercase{#1}}\par
  \vspace{4pt}\noindent{\color{poporange}\rule{36pt}{3.6pt}}
  \end{minipage}\par\nopagebreak\vspace{8pt}
}
\newcounter{popsub}
\newcommand{\courssub}[1]{%
  \stepcounter{popsub}%
  \par\vspace{9pt}\noindent{\popxbold\fontsize{11.5}{13}\selectfont\color{popink}{\color{poporange}\thepopsec.\thepopsub}\hspace{6pt}#1}\par\nopagebreak\vspace{4pt}
}

% ============================================================
% Boîtes pédagogiques — même recette : fond blanc, bordure épaisse colorée,
% ombre dure encre, badge plein. Titre optionnel [#1].
% ============================================================
\tcbset{popbase/.style={breakable,enhanced,colback=white,boxrule=2.3pt,arc=9pt,
  shadow={3pt}{-3pt}{0pt}{popink},left=14pt,right=14pt,top=11pt,bottom=11pt,before skip=11pt,after skip=15pt,
  fontupper=\popsemi\fontsize{10.6}{14.5}\selectfont\color{popink},
  fontlower=\popsemi\fontsize{10.6}{14.5}\selectfont\color{popink}}}
% Coche dessinée (le glyphe ✓ n'existe pas dans les polices du document)
\newcommand{\popcheck}[1]{\tikz[baseline=-0.5ex]\draw[#1,line width=1.6pt,line cap=round,line join=round](-0.13,0.0)--(-0.04,-0.09)--(0.13,0.1);}
\providecommand{\coche}{\popcheck{popgreen}}
\providecommand{\resultat}[1]{\par\smallskip\noindent\fcolorbox{popgreen}{popgreenL}{\parbox{\dimexpr\linewidth-2\fboxsep-2\fboxrule\relax}{\textbf{Result.} #1}}\par\smallskip}
\newcommand{\popboxhead}[3]{\poppill{#1}{#2}{white}\ifx\relax#3\relax\else\hspace{6pt}{\popxbold\fontsize{10.6}{12}\selectfont\color{popink}#3}\fi\par\vspace{7pt}}

\newtcolorbox{aretenir}[1][]{popbase,colframe=popgreen,before upper={\popboxhead{Key points}{popgreen}{#1}}}
\newtcolorbox{exemplebox}[1][]{popbase,colframe=poporange,before upper={\popboxhead{Exemple}{poporange}{#1}}}
\newtcolorbox{definition}[1][]{popbase,colframe=popblue,before upper={\popboxhead{Definition}{popblue}{#1}}}
\newtcolorbox{propriete}[1][]{popbase,colframe=poppurple,before upper={\popboxhead{Property}{poppurple}{#1}}}
\newtcolorbox{methode}[1][]{popbase,colframe=popgold,before upper={\popboxhead{Method}{popgold}{#1}}}
\newtcolorbox{attention}[1][]{popbase,colframe=poppink,before upper={\popboxhead{Warning}{poppink}{#1}}}
\newtcolorbox{experience}[1][]{popbase,colframe=popblue,colback=popblueL,before upper={\popboxhead{Experiment}{popblue}{#1}}}
% « Le savais-tu ? » : carte violette pleine (contexte, culture, Cameroun)
\newtcolorbox{savaistu}[1][]{popbase,colback=poppurple,colframe=popink,
  fontupper=\popsemi\fontsize{10.4}{14.2}\selectfont\color{white},
  before upper={\poppill{Did you know?}{white}{poppurple}\ifx\relax#1\relax\else\hspace{6pt}{\popxbold\color{white}#1}\fi\par\vspace{7pt}}}
% Exercice résolu : énoncé en haut, \tcblower puis la solution sur fond crème
\newcounter{popexo}\newcounter{popexr}
\newtcolorbox{exoresolu}[1][]{popbase,colframe=poporange,colbacklower=popcream,
  segmentation style={popink,line width=1.2pt,dashed},
  before upper={\stepcounter{popexr}\popboxhead{Worked example \thepopexr}{poporange}{#1}},
  before lower={\poppill{Solution}{popink}{popcream}\par\vspace{6pt}}}
% Exercice (non résolu en place) — la correction est dans la partie Corrigés
\newtcolorbox{exercice}[1][]{popbase,colframe=popink,
  before upper={\stepcounter{popexo}\popboxhead{Exercise \thepopexo}{popink}{#1}}}
% Corrigé
\newtcolorbox{corrige}[1][]{popbase,colframe=popgreen,colback=popgreenL,
  before upper={\popboxhead{Solution}{popgreen}{#1}}}
% Objectifs / prérequis / bilan
\newtcolorbox{objectifs}{popbase,colframe=popink,colback=poporangeL,
  before upper={\popboxhead{Chapter objectives}{popink}{}}}
\newtcolorbox{prerequis}{popbase,colframe=popink,colback=popblueL,
  before upper={\popboxhead{Prerequisites}{popblue}{}}}
\newtcolorbox{bilan}{popbase,colframe=popink,colback=white,boxrule=2.8pt,
  before upper={\popboxhead{Fiche bilan}{poporange}{L'essentiel à connaître pour l'examen}}}
% Figure encadrée avec légende
\newtcolorbox{popfigure}[1][]{enhanced,breakable=false,colback=white,colframe=popline,boxrule=1.4pt,arc=8pt,
  left=10pt,right=10pt,top=10pt,bottom=8pt,before skip=10pt,after skip=12pt,halign=center,
  lower separated=false,halign lower=center,
  fontlower=\popsemi\fontsize{9}{11.5}\selectfont\color{popmuted},
  #1}
\newcommand{\legende}[1]{\par\vspace{6pt}{\popsemi\fontsize{9}{11.5}\selectfont\color{popmuted}#1}}

% ============================================================
% Signature OnBuch+
% ============================================================
\newcommand{\poplogoinline}[1]{{\poplogo\fontsize{#1}{#1}\selectfont\color{popink}OnBuch\color{poporange}+}}
\newcommand{\signatureonbuch}{%
  \begin{tcolorbox}[enhanced,colback=popink,colframe=popink,boxrule=2.2pt,arc=10pt,
      drop shadow=poporange,left=14pt,right=14pt,top=10pt,bottom=10pt,before skip=0pt,after skip=0pt]
    \tikz[baseline=-0.7ex]\node[circle,fill=popgreen,draw=popcream,line width=1.3pt,minimum size=22pt,inner sep=0]{\popcheck{white}};%
    \hspace{9pt}\begin{minipage}[c]{0.62\linewidth}
      {\popxbold\fontsize{12}{13}\selectfont\color{popcream}Signed\ \textbullet\ {\poplogo OnBuch\color{poporange}+}}\par\vspace{2pt}
      {\raggedright\popsemi\fontsize{8}{10.5}\selectfont\color{popline}\MakeUppercase{OnBuch+ teaching team}\\\MakeUppercase{Follows the official MINESEC syllabus}\par}
    \end{minipage}\hfill
    {\poplogo\fontsize{22}{22}\selectfont\color{popcream}OnBuch\color{poporange}+}
  \end{tcolorbox}%
}

% ============================================================
% Page de garde
% #1 titre · #2 matière · #3 niveau · #4 module · #5 n° leçon · #6 durée ·
% #7 description / objectifs (texte ou itemize)
% ============================================================
\newcommand{\pagegarde}[7]{%
  \thispagestyle{empty}
  \newgeometry{a4paper,margin=1.7cm,top=1.6cm,bottom=1.4cm}%
  \begin{tikzpicture}[remember picture,overlay]
    \fill[poporange!18] ([xshift=-3.2cm,yshift=-3.6cm]current page.north east) circle (4.6cm);
    \fill[poppurple!14] ([xshift=2.4cm,yshift=7.6cm]current page.south west) circle (3.8cm);
  \end{tikzpicture}%
  {\poplogo\fontsize{30}{30}\selectfont\color{popink}OnBuch\color{poporange}+}\hfill
  \raisebox{0.6ex}{\poppill{Signed course \textbullet\ Premium}{popink}{popcream}}\par
  \vspace{1pt}\popsquiggle{poppurple}\par
  \vspace{8pt}
  \begin{tcolorbox}[enhanced,colback=poporange,colframe=popink,boxrule=2.8pt,arc=13pt,
      drop shadow=popink,left=18pt,right=18pt,top=12pt,bottom=13pt,after skip=10pt]
    \poppill{#2 \textbullet\ #3}{popink}{popcream}\hspace{5pt}\poppill{#4}{white}{popink}\par\vspace{8pt}
    {\popdisplay\fontsize{14}{14}\selectfont\color{popink}CHAPTER #5}\par\vspace{4pt}
    {\raggedright\hyphenpenalty=10000\exhyphenpenalty=10000\popdisplay\fontsize{25}{27}\selectfont\color{popcream}\MakeUppercase{#1}\par}
    \vspace{8pt}
    {\popxbold\fontsize{9.5}{9.5}\selectfont\color{popink}\MakeUppercase{Suggested time: #6}}
  \end{tcolorbox}
  \begin{tcolorbox}[enhanced,colback=white,colframe=popink,boxrule=2.2pt,arc=10pt,
      drop shadow=popink,left=16pt,right=16pt,top=10pt,bottom=10pt,after skip=8pt]
    \poppill{In this chapter}{poppurple}{white}\par\vspace{5pt}
    \popsemi\fontsize{9.3}{12.6}\selectfont\color{popink}#7
  \end{tcolorbox}
  \noindent\begin{minipage}[t]{0.487\linewidth}
    \begin{tcolorbox}[enhanced,colback=popgreen,colframe=popink,boxrule=2.2pt,arc=10pt,
        drop shadow=popink,left=11pt,right=11pt,top=10pt,bottom=10pt,height=3.1cm,valign=center]
      \tcbox[colback=white,colframe=popink,boxrule=1.3pt,arc=5pt,boxsep=0pt,nobeforeafter,
        left=3pt,right=3pt,top=3pt,bottom=3pt,valign=center]{\qrcode[height=1.9cm]{https://play.google.com/store/apps/details?id=com.luvvix.onbuch&hl=en}}%
      \hspace{8pt}\begin{minipage}[c]{0.5\linewidth}
        {\popdisplay\fontsize{11}{12.5}\selectfont\color{white}\MakeUppercase{Download OnBuch}}\par\vspace{4pt}
        {\raggedright\popsemi\fontsize{8.4}{11}\selectfont\color{white}Free \textbullet\ Google Play\\AI tutor, past papers, results\par}
      \end{minipage}
    \end{tcolorbox}
  \end{minipage}\hfill
  \begin{minipage}[t]{0.487\linewidth}
    \begin{tcolorbox}[enhanced,colback=poppurple,colframe=popink,boxrule=2.2pt,arc=10pt,
        drop shadow=popink,left=11pt,right=11pt,top=10pt,bottom=10pt,height=3.1cm,valign=center]
      \tikz[baseline=-0.7ex]\node[circle,fill=white,draw=popink,line width=1.3pt,minimum size=1.6cm,inner sep=0]{\popdisplay\fontsize{24}{24}\selectfont\color{poppurple}+};%
      \hspace{8pt}\begin{minipage}[c]{0.6\linewidth}
        {\popdisplay\fontsize{11}{12.5}\selectfont\color{white}\MakeUppercase{Go OnBuch+}}\par\vspace{4pt}
        {\raggedright\popsemi\fontsize{8.4}{11}\selectfont\color{white}All signed courses, unlimited AI tutor, PDF sheets — OnBuch+ tab in the app\par}
      \end{minipage}
    \end{tcolorbox}
  \end{minipage}
  \vspace{8pt}
  \popcontenu
  \vfill
  \signatureonbuch
  \restoregeometry
  \newpage
}

% Rangée « Ce cours contient » (page de garde)
\newcommand{\poptile}[3]{% #1 couleur #2 gros texte #3 légende
  \begin{tcolorbox}[enhanced,colback=#1,colframe=popink,boxrule=2pt,arc=9pt,shadow={2.5pt}{-2.5pt}{0pt}{popink},
      left=6pt,right=6pt,top=9pt,bottom=9pt,halign=center,valign=center,height=1.75cm,width=\linewidth,nobeforeafter]
    {\popdisplay\fontsize{12.5}{13}\selectfont\color{white}\MakeUppercase{#2}}\par\vspace{3pt}
    {\centering\popsemi\fontsize{7.8}{9.5}\selectfont\color{white}#3\par}
  \end{tcolorbox}}
\newcommand{\popcontenu}{%
  \par\noindent{\popxboldls\fontsize{7.5}{7.5}\selectfont\color{popmuted}THIS SIGNED COURSE CONTAINS}\par\vspace{7pt}
  \noindent\begin{minipage}[t]{0.235\linewidth}\poptile{popblue}{Course}{complete, illustrated, step by step}\end{minipage}\hfill
  \begin{minipage}[t]{0.235\linewidth}\poptile{poporange}{Methods}{and worked examples}\end{minipage}\hfill
  \begin{minipage}[t]{0.235\linewidth}\poptile{poppink}{Exercises}{exam style + solutions}\end{minipage}\hfill
  \begin{minipage}[t]{0.235\linewidth}\poptile{popgreen}{Summary sheet}{the essentials to remember}\end{minipage}\par}

% Sommaire
\newtcolorbox{sommaire}{enhanced,colback=white,colframe=popink,boxrule=2.2pt,arc=10pt,
  drop shadow=popink,left=18pt,right=18pt,top=13pt,bottom=13pt,before skip=6pt,after skip=20pt,
  before upper={\poppill{Contents}{poppurple}{white}\par\vspace{9pt}\popsemi\fontsize{10.6}{14.5}\selectfont\color{popink}}}

% Signature de fin de cours — poussée tout en bas de la dernière page
\newcommand{\findecours}{%
  \par\vfill\nopagebreak
  \begin{center}\popsquiggle{poporange}\end{center}\vspace{4pt}
  \signatureonbuch
}
\AtBeginDocument{\def\llbracket{\mathopen{[\mkern-3.5mu[}}\def\rrbracket{\mathclose{]\mkern-3.5mu]}}} % intervalles d'entiers (glyphe ⟦ absent de la police)
% Glyphes absents de Fira Math : redéfinis à partir de glyphes disponibles
\AtBeginDocument{%
  \def\setminus{\mathbin{\backslash}}\let\smallsetminus\setminus
  \def\vdots{\mathord{\vbox{\baselineskip3.2pt\lineskiplimit0pt\kern4pt\hbox{.}\hbox{.}\hbox{.}}}}%
  \def\ddots{\mathinner{\mkern1mu\raise6.5pt\hbox{.}\mkern2mu\raise3.6pt\hbox{.}\mkern2mu\raise0.7pt\hbox{.}\mkern1mu}}%
  \def\triangle{\mathord{\scalebox{0.9}{$\symup{\Delta}$}}}%
  \def\nabla{\mathord{\raisebox{\depth}{\scalebox{1}[-1]{$\symup{\Delta}$}}}}%
  \def\checkmark{\popcheck{popdark}}%
  \def\star{\mathbin{\ast}}\def\bigstar{\mathord{\ast}}%
  \def\diamond{\mathbin{\scalebox{0.6}{$\Diamond$}}}%
  \def\bigcup{\mathop{\mathchoice{\raisebox{-0.3em}{\scalebox{1.7}{$\cup$}}}{\scalebox{1.25}{$\cup$}}{\cup}{\cup}}\displaylimits}%
  \def\bigcap{\mathop{\mathchoice{\raisebox{-0.3em}{\scalebox{1.7}{$\cap$}}}{\scalebox{1.25}{$\cap$}}{\cap}{\cap}}\displaylimits}%
  \def\lesseqqgtr{\mathrel{\lessgtr}}\def\bigoplus{\mathop{\oplus}\displaylimits}%
}
\providecommand{\ssi}{if and only if} % abréviation fréquente
\AtBeginDocument{\providecommand{\wideparen}{\overparen}} % arc de cercle

% Fonctions usuelles que les modèles utilisent sans que amsmath les fournisse
\providecommand{\arsinh}{\operatorname{arsinh}}\providecommand{\arcosh}{\operatorname{arcosh}}
\providecommand{\artanh}{\operatorname{artanh}}\providecommand{\arsech}{\operatorname{arsech}}
\providecommand{\arcsch}{\operatorname{arcsch}}\providecommand{\arcoth}{\operatorname{arcoth}}
\providecommand{\arcsinh}{\operatorname{arsinh}}\providecommand{\arccosh}{\operatorname{arcosh}}
\providecommand{\arctanh}{\operatorname{artanh}}\providecommand{\sech}{\operatorname{sech}}
\providecommand{\csch}{\operatorname{csch}}\providecommand{\arcsec}{\operatorname{arcsec}}
\providecommand{\arccsc}{\operatorname{arccsc}}\providecommand{\arccot}{\operatorname{arccot}}
\providecommand{\sgn}{\operatorname{sgn}}\providecommand{\lcm}{\operatorname{lcm}}
\providecommand{\Var}{\operatorname{Var}}\providecommand{\Cov}{\operatorname{Cov}}

%% FIGFIX-BEGIN v3 — correctifs globaux des figures (docs/onbuch-generation/tools/figfix)
\usepackage{adjustbox}\usepackage{etoolbox}
% toute figure TikZ/pgfplots plus large que la ligne est réduite à la largeur disponible
\BeforeBeginEnvironment{tikzpicture}{\begin{adjustbox}{max width=\linewidth}}
\AfterEndEnvironment{tikzpicture}{\end{adjustbox}}
% légendes pgfplots lisibles par-dessus les courbes
\pgfplotsset{every axis/.append style={legend style={fill=white,fill opacity=0.92,text opacity=1,draw=black!15,inner sep=3pt}}}
% étiquettes d'arêtes (syntaxe edge["..."]) avec fond blanc
\tikzset{every edge quotes/.append style={fill=white,inner sep=1.2pt}}
% bulles de cartes mentales : texte à la ligne dans la bulle, bulle un peu plus grande
\tikzset{every concept/.append style={text width=2.0cm,align=center,minimum size=2.3cm,inner sep=2pt}}
%% FIGFIX-END
\begin{document}

\pagegarde{Numerical Methods}{Pure Mathematics with Statistics}{Upper Sixth}{Upper Sixth — Pure mathematics with statistics}{6}{5 periods}{This chapter introduces numerical methods for solving equations that cannot be solved exactly: locating roots by sign changes, then refining them with the bisection method and the Newton–Raphson iteration. Students learn to control accuracy, justify convergence, and apply these techniques to GCE-style modelling problems.
\vspace{4pt}
\begin{multicols}{2}\small
\begin{itemize}
\item By the end of this chapter I can explain why numerical methods are needed and what an approximate root is.
\item By the end of this chapter I can locate a root of f(x) = 0 in an interval using a change of sign of a continuous function.
\item By the end of this chapter I can use tables of values and graphs to bracket roots and count the roots of an equation.
\item By the end of this chapter I can rearrange an equation into the form x = g(x) and use simple iteration x\_{n+1} = g(x\_n) where appropriate.
\item By the end of this chapter I can carry out the bisection method step by step and state the interval containing the root after n steps.
\item By the end of this chapter I can apply the Newton–Raphson formula x\_{n+1} = x\_n − f(x\_n)/f′(x\_n) to find a root to a given accuracy.
\end{itemize}\end{multicols}}

\begin{sommaire}
\begin{multicols}{2}
\begin{enumerate}
\item Introduction and Location of Roots
\item The Bisection Method
\item The Newton–Raphson Method
\item Accuracy, Comparison and Applications
\item Integration activity
\item Methods and worked examples
\item Exercises
\item Solutions to the exercises
\item Summary sheet
\end{enumerate}
\end{multicols}
\end{sommaire}

\begin{objectifs}
\begin{itemize}
\item By the end of this chapter I can explain why numerical methods are needed and what an approximate root is.
\item By the end of this chapter I can locate a root of f(x) = 0 in an interval using a change of sign of a continuous function.
\item By the end of this chapter I can use tables of values and graphs to bracket roots and count the roots of an equation.
\item By the end of this chapter I can rearrange an equation into the form x = g(x) and use simple iteration x\_{n+1} = g(x\_n) where appropriate.
\item By the end of this chapter I can carry out the bisection method step by step and state the interval containing the root after n steps.
\item By the end of this chapter I can apply the Newton–Raphson formula x\_{n+1} = x\_n − f(x\_n)/f′(x\_n) to find a root to a given accuracy.
\item By the end of this chapter I can interpret iteration geometrically using tangent lines and staircase/cobweb diagrams.
\item By the end of this chapter I can decide when a numerical method fails or converges slowly and choose a suitable starting value.
\item By the end of this chapter I can determine the accuracy of an approximation (decimal places, error bounds) and apply numerical methods to real-life problems.
\end{itemize}
\end{objectifs}

\begin{prerequis}
\begin{itemize}
\item Differentiation of polynomials, exponentials, logarithms and trigonometric functions (Lower Sixth calculus)
\item Continuity and the behaviour of graphs of functions, including sketching curves
\item Solving equations algebraically and manipulating functions (factorisation, rearrangement)
\item Use of a scientific calculator in radians and degrees, and correct rounding to given decimal places
\item Sequences and the idea of a limit (recurrence relations)
\end{itemize}
\end{prerequis}

\newpage

\coursec{Introduction and Location of Roots}

A water-supply engineer in Bafoussam must size a pipe so that the flow equation $x^3 - 7x + 3 = 0$ is satisfied, yet no algebraic formula gives the answer. A mobile-money agent in Douala wants to know the exact interest rate hidden in a repayment schedule — again, the equation cannot be solved exactly. In real life, most equations that matter have no neat closed-form solution. Numerical methods give us powerful, systematic ways to find roots to any desired accuracy, using only arithmetic. This chapter shows how a simple calculator and a clever algorithm can solve problems that defeat pure algebra.

\courssub{Why Numerical Methods?}

From Lower Sixth onwards you have solved equations such as $x^2 - 5x + 6 = 0$ (factorisation or the quadratic formula) and $\sin x = \tfrac{1}{2}$ (standard values). But try solving any of these:

\begin{itemize}
\item $x^3 - x - 2 = 0$ (a cubic with no rational root);
\item $e^x = 3x$ (mixing an exponential with a linear term);
\item $\cos x = x$ (mixing a trigonometric function with a polynomial).
\end{itemize}

No algebraic manipulation will ever produce an exact answer in terms of surds, fractions or standard constants. Such equations are called \cle{transcendental} or simply \emph{non-algebraic} equations. Yet each of them \emph{does} have solutions, and in engineering, finance, physics and statistics we genuinely need those solutions — often to 4, 6 or 10 decimal places.

\begin{definition}[Root of an equation]
A \cle{root} of the equation $f(x) = 0$ (also called a \cle{zero} of the function $f$) is a number $\alpha$ such that $f(\alpha) = 0$. Geometrically, the roots of $f(x) = 0$ are the $x$-coordinates of the points where the curve $y = f(x)$ meets the $x$-axis.
\end{definition}

\begin{definition}[Approximate root and tolerance]
An \cle{approximate root} of $f(x) = 0$ is a number $x_0$ which is close to the true root $\alpha$. We say $x_0$ approximates $\alpha$ \emph{with tolerance} $\varepsilon$ (a small positive number) if
\[ |x_0 - \alpha| < \varepsilon. \]
For example, stating a root ``to 2 decimal places'' means the error is at most $0.005$; the tolerance is $\varepsilon = 0.005$.
\end{definition}

\begin{attention}[Exact versus approximate]
Writing $x = 1.521$ as ``the solution'' of $x^3 - x - 2 = 0$ is sloppy: it is an \emph{approximation}. The correct statement is ``the root is $1.521$ correct to 3 decimal places'', meaning the true root $\alpha$ satisfies $1.5205 \leq \alpha < 1.5215$. In the GCE examination, always state the accuracy of your answer.
\end{attention}

The strategy of this whole chapter has two phases:
\begin{enumerate}
\item \textbf{Locate} the root: find an interval $[a, b]$ that is known to contain it (this section).
\item \textbf{Refine} the root: use an algorithm — bisection (Section 2) or Newton--Raphson (Section 3) — to shrink the interval or improve the estimate until the desired tolerance is reached.
\end{enumerate}

\courssub{The Intermediate Value Theorem: the Sign-Change Criterion}

Suppose a function $f$ is \cle{continuous} on an interval $[a, b]$ — informally, its graph can be drawn from $(a, f(a))$ to $(b, f(b))$ without lifting the pen. If $f(a)$ is negative and $f(b)$ is positive, the curve starts below the $x$-axis and ends above it. Somewhere in between, it \emph{must} cross the axis. That crossing point is a root.

\begin{propriete}[Intermediate Value Theorem — sign-change criterion]
Let $f$ be continuous on the closed interval $[a, b]$. If $f(a)$ and $f(b)$ have \textbf{opposite signs} (i.e. $f(a)\,f(b) < 0$), then the equation $f(x) = 0$ has \textbf{at least one root} in the open interval $(a, b)$. (The full proof belongs to analysis and is \emph{not} examinable; the statement and its use \emph{are}.)
\end{propriete}

\begin{popfigure}
\begin{center}
\begin{tikzpicture}[scale=1.1]
\draw[->, thick, popmuted] (-0.5,0) -- (5.5,0) node[right] {$x$};
\draw[->, thick, popmuted] (0,-1.8) -- (0,2.6) node[above] {$y$};
\draw[very thick, popblue] (0.8,-1.2) .. controls (1.8,-1.6) and (2.2,1.8) .. (3.0,0.9) .. controls (3.6,0.2) and (4.0,1.2) .. (4.6,2.0);
\fill[poporange] (0.8,-1.2) circle (2pt) node[below left, popink] {$(a,\ f(a))<0$};
\fill[poporange] (4.6,2.0) circle (2pt) node[above right, popink] {$(b,\ f(b))>0$};
\draw[dashed, popmuted] (0.8,-1.2) -- (0.8,0) node[above, popink] {$a$};
\draw[dashed, popmuted] (4.6,2.0) -- (4.6,0) node[below, popink] {$b$};
\fill[popgreen] (2.42,0) circle (2.5pt) node[below, popink] {$\alpha$};
\end{tikzpicture}
\legende{A continuous curve with $f(a)<0$ and $f(b)>0$ must cross the $x$-axis at least once between $a$ and $b$.}
\end{center}
\end{popfigure}

\begin{attention}[What the sign change does and does not say]
\begin{itemize}
\item A sign change \emph{proves} at least one root in $(a,b)$, but says \textbf{nothing about uniqueness}: there could be 3, 5, \ldots\ roots (any odd number, counting multiplicity).
\item \textbf{No} sign change does \textbf{not} prove there is no root: the interval could contain an \emph{even} number of roots, or a \emph{double root} where the curve is tangent to the axis.
\item If $f$ is \emph{not continuous} on $[a,b]$, the test is unreliable: $f(x) = \dfrac{1}{x}$ gives $f(-1) = -1 < 0$ and $f(1) = 1 > 0$, yet $1/x = 0$ has \emph{no} solution — the ``sign change'' comes from the discontinuity at $x = 0$, a \emph{false positive}.
\end{itemize}
\end{attention}

\begin{popfigure}
\begin{center}
\begin{tikzpicture}[scale=1.3]
\draw[->, thick, popmuted] (-2.6,0) -- (2.6,0) node[right] {$x$};
\draw[->, thick, popmuted] (0,-0.5) -- (0,2.8) node[above] {$y$};
\draw[very thick, poppurple, domain=-2.2:2.2, samples=80] plot (\x, {0.25*\x*\x});
\fill[poporange] (0,0) circle (2.5pt) node[below right, popink] {double root};
\node[popink, align=center] at (0,2.4) {$f(x)>0$ on both sides: \\ no sign change, yet a root exists};
\end{tikzpicture}
\legende{Counter-example: a curve tangent to the $x$-axis (double root) produces no sign change.}
\end{center}
\end{popfigure}

\courssub{Locating Roots by a Table of Values}

The most direct way to exploit the sign-change criterion is to evaluate $f(x)$ at integer points and look for a sign change between consecutive values.

\begin{methode}[Locating roots with a table of values]
\begin{enumerate}
\item Write the equation in the form $f(x) = 0$.
\item Compute $f(x)$ at consecutive integers $x = \ldots, -2, -1, 0, 1, 2, \ldots$ over the required range.
\item Each sign change between $n$ and $n+1$ brackets a root in $(n, n+1)$ — provided $f$ is continuous there.
\item To locate the root more finely, tabulate $f$ at steps of $0.1$ inside that unit interval, then $0.01$, and so on.
\item Check the graph (or further values) to make sure no root has been missed (even numbers of roots hide from the test).
\end{enumerate}
\end{methode}

\begin{exoresolu}[Locating the root of $x^3 - x - 2 = 0$]
Show that the equation $x^3 - x - 2 = 0$ has a root between $1$ and $2$, and locate it between two consecutive tenths.
\tcblower
Let $f(x) = x^3 - x - 2$; $f$ is a polynomial, hence continuous everywhere. Compute:
\begin{center}
\begin{tabular}{|c|cccc|}
\hline
\rowcolor{popblueL} $x$ & $0$ & $1$ & $2$ & $3$ \\
\hline
$f(x)$ & $-2$ & $-2$ & $4$ & $22$ \\
\hline
\end{tabular}
\end{center}
Since $f(1) = -2 < 0$ and $f(2) = 4 > 0$, by the Intermediate Value Theorem there is a root $\alpha \in (1, 2)$. Now refine with steps of $0.1$:
\begin{center}
\begin{tabular}{|c|ccccc|}
\hline
\rowcolor{popblueL} $x$ & $1.3$ & $1.4$ & $1.5$ & $1.6$ & $1.7$ \\
\hline
$f(x)$ & $-1.103$ & $-0.656$ & $-0.125$ & $0.496$ & $1.213$ \\
\hline
\end{tabular}
\end{center}
The sign change occurs between $1.5$ and $1.6$: hence $1.5 < \alpha < 1.6$. (In fact $\alpha \approx 1.5214$.)
\end{exoresolu}

\begin{popfigure}
\begin{center}
\begin{tikzpicture}
\begin{axis}[
  width=11cm, height=7cm,
  axis lines=middle,
  xlabel={$x$}, ylabel={$y$},
  xmin=-2.2, xmax=2.6, ymin=-4.5, ymax=5.5,
  xtick={-2,-1,1,1.5214,2}, xticklabels={$-2$,$-1$,$1$,$\alpha$,$2$},
  ytick={-2,4},
  domain=-2.1:2.3, samples=120,
  every axis x label/.style={at={(ticklabel* cs:1.02)}, anchor=west},
]
\addplot[very thick, popblue] {x^3 - x - 2};
\addplot[only marks, mark=*, mark size=2pt, poporange] coordinates {(1,-2) (2,4)};
\addplot[only marks, mark=*, mark size=2.5pt, popgreen] coordinates {(1.5214,0)};
\node[popink, anchor=west] at (axis cs:1.05,-2.6) {$f(1)=-2<0$};
\node[popink, anchor=south west] at (axis cs:1.55,4.1) {$f(2)=4>0$};
\draw[dashed, popmuted] (axis cs:1,-2) -- (axis cs:1,0);
\draw[dashed, popmuted] (axis cs:2,4) -- (axis cs:2,0);
\end{axis}
\end{tikzpicture}
\legende{Graph of $f(x)=x^3-x-2$: the sign change between $x=1$ and $x=2$ brackets the root $\alpha \approx 1.5214$.}
\end{center}
\end{popfigure}

\courssub{Graphical Location of Roots}

A quick sketch often reveals \emph{how many} real roots an equation has and roughly \emph{where} they are — essential information before starting any numerical algorithm, since the table method can miss roots.

\textbf{Method 1: sketch $y = f(x)$ directly.} The roots are the $x$-intercepts. This works well when $f$ is easy to sketch.

\textbf{Method 2: split into two familiar curves.} Rearrange $f(x) = 0$ as $g(x) = h(x)$ where $g$ and $h$ are standard graphs (lines, parabolas, exponentials, trigonometric curves). The roots of $f(x) = 0$ are the $x$-coordinates of the \emph{intersection points} of $y = g(x)$ and $y = h(x)$.

\begin{methode}[Locating roots graphically by splitting]
\begin{enumerate}
\item Rearrange the equation into the form $g(x) = h(x)$ with $g, h$ easy to sketch.
\item Sketch $y = g(x)$ and $y = h(x)$ on the same axes (a careful freehand sketch is enough).
\item Count the intersection points: this gives the \emph{number} of real roots.
\item Read off approximate positions and choose \emph{integer bounds} $n, n+1$ bracketing each root.
\item Confirm each bracket by a sign change of $f(x) = g(x) - h(x)$.
\end{enumerate}
\end{methode}

\begin{exemplebox}[The equation $e^x = 3x$]
Take $g(x) = e^x$ and $h(x) = 3x$. The exponential curve and the straight line meet \emph{twice}, so $e^x = 3x$ has exactly two real roots. Reading the graph, one root lies between $0$ and $1$, the other between $1$ and $2$. Check with $f(x) = e^x - 3x$:
\[ f(0) = 1 > 0,\quad f(1) = e - 3 \approx -0.28 < 0,\quad f(2) = e^2 - 6 \approx 1.39 > 0. \]
Both sign changes are confirmed: roots in $(0,1)$ and $(1,2)$.
\end{exemplebox}

\begin{popfigure}
\begin{center}
\begin{tikzpicture}
\begin{axis}[
  width=11cm, height=7.5cm,
  axis lines=middle,
  xlabel={$x$}, ylabel={$y$},
  xmin=-1.2, xmax=2.4, ymin=-1, ymax=8,
  xtick={-1,1,2}, ytick={2,4,6},
  legend style={at={(0.03,0.97)}, anchor=north west, draw=popline},
]
\addplot[very thick, popblue, domain=-1.2:2.1, samples=100] {exp(x)};
\addlegendentry{$y=e^x$}
\addplot[very thick, poporange, domain=-1.2:2.4, samples=3] {3*x};
\addlegendentry{$y=3x$}
\addplot[only marks, mark=*, mark size=2.5pt, popgreen] coordinates {(0.619,1.857) (1.512,4.537)};
\node[popink, anchor=north] at (axis cs:0.619,1.3) {$\alpha_1$};
\node[popink, anchor=south east] at (axis cs:1.512,4.9) {$\alpha_2$};
\end{axis}
\end{tikzpicture}
\legende{The curves $y=e^x$ and $y=3x$ intersect twice: the equation $e^x=3x$ has two real roots, $\alpha_1 \in (0,1)$ and $\alpha_2 \in (1,2)$.}
\end{center}
\end{popfigure}

\begin{exoresolu}[Counting and bracketing all roots of $\cos x = x$ and of $x^3 - 7x + 3 = 0$]
(a) Show that $\cos x = x$ has exactly one real root and bracket it between consecutive integers. (b) Show that $x^3 - 7x + 3 = 0$ has three real roots and give integer bounds for each.
\tcblower
\textbf{(a)} Sketch $y = \cos x$ (oscillating between $-1$ and $1$) and $y = x$ (a straight line through the origin). For $x > 1$ the line is above the cosine's maximum; for $x < -1$ it is below its minimum. They cross exactly once, in $(0, 1)$. Check with $f(x) = \cos x - x$ ($x$ in radians):
\[ f(0) = 1 > 0, \qquad f(1) = \cos 1 - 1 \approx 0.5403 - 1 = -0.4597 < 0. \]
Sign change confirmed: the unique root lies in $(0, 1)$.

\textbf{(b)} Let $f(x) = x^3 - 7x + 3$, continuous everywhere. Tabulate at integers:
\begin{center}
\begin{tabular}{|c|cccccccc|}
\hline
\rowcolor{popblueL} $x$ & $-3$ & $-2$ & $-1$ & $0$ & $1$ & $2$ & $3$ & $4$ \\
\hline
$f(x)$ & $-3$ & $9$ & $9$ & $3$ & $-3$ & $-3$ & $9$ & $39$ \\
\hline
\end{tabular}
\end{center}
Sign changes occur between $-3$ and $-2$, between $0$ and $1$, and between $2$ and $3$. Since a cubic has \emph{at most} three real roots, these are all of them:
\[ \alpha_1 \in (-3,-2), \qquad \alpha_2 \in (0,1), \qquad \alpha_3 \in (2,3). \]
This is exactly the engineer's equation from Bafoussam — and we now know it has three real solutions, so a physical constraint will be needed to select the meaningful one.
\end{exoresolu}

\begin{attention}[Traps in locating roots]
\begin{itemize}
\item \textbf{Radians!} When mixing trigonometric functions with polynomials ($\cos x = x$), your calculator must be in \emph{radian} mode.
\item A single table at unit steps can \emph{miss} two roots lying in the same unit interval (e.g. roots at $1.2$ and $1.8$). Always combine the table with a sketch or with knowledge of the maximum number of roots (a polynomial of degree $n$ has at most $n$ real roots).
\item Never apply the sign-change test across an asymptote: $f(x) = \dfrac{1}{x}$ ``changes sign'' between $-1$ and $1$ but has no zero there.
\end{itemize}
\end{attention}

\courssub{First Ideas of Iteration}

Locating a root between two integers is only the beginning. To \emph{refine} it, numerical methods use \cle{iteration}: starting from an estimate $x_0$, we apply a rule that produces a (hopefully better) estimate $x_1$, then $x_2$, and so on, generating a sequence $(x_n)$.

One natural idea is to rearrange $f(x) = 0$ into the form
\[ x = g(x) \]
and then compute
\[ x_{n+1} = g(x_n), \qquad n = 0, 1, 2, \ldots \]
If the sequence $(x_n)$ settles down to a limit $\ell$, then (provided $g$ is continuous) $\ell = g(\ell)$, so $\ell$ is a root of the rearranged equation. We say the iteration \cle{converges}. If the terms drift away or oscillate wildly, it \cle{diverges} — and the choice of rearrangement matters enormously.

\begin{exemplebox}[A first taste of iteration: $x^3 - x - 2 = 0$]
Rearrange as $x = \sqrt[3]{x + 2}$, i.e. $g(x) = (x+2)^{1/3}$. Starting from $x_0 = 1.5$:
\begin{align*}
x_1 &= (1.5 + 2)^{1/3} \approx 1.5183,\\
x_2 &= (1.5183 + 2)^{1/3} \approx 1.5211,\\
x_3 &\approx 1.5214, \qquad x_4 \approx 1.5214.
\end{align*}
The sequence converges rapidly to the root $\alpha \approx 1.5214$ found earlier. But the equally valid rearrangement $x = x^3 - 2$ gives $x_1 = 1.375$, $x_2 \approx 0.5996$, $x_3 \approx -1.7846$, $x_4 \approx -7.684$, \ldots — divergence! Same equation, same starting point, different rearrangement, opposite behaviour.
\end{exemplebox}

\begin{aretenir}[Key points of this section]
\begin{itemize}
\item A \cle{root} of $f(x) = 0$ is a value $\alpha$ with $f(\alpha) = 0$; most real-world equations can only be solved \emph{approximately}, to a stated \cle{tolerance}.
\item \textbf{Intermediate Value Theorem:} if $f$ is continuous on $[a,b]$ and $f(a)f(b) < 0$, then $f(x) = 0$ has at least one root in $(a,b)$.
\item A sign change proves \emph{existence}, not \emph{uniqueness}; absence of a sign change does \emph{not} prove absence of roots (double roots, even numbers of roots); discontinuities can create false sign changes.
\item Roots can be located by a \emph{table of values} or \emph{graphically}, either from $y = f(x)$ or by splitting into two intersecting curves $y = g(x)$, $y = h(x)$.
\item Iteration $x_{n+1} = g(x_n)$ generates a sequence of estimates that may converge to a root — the engine behind the bisection and Newton--Raphson methods of the next sections.
\end{itemize}
\end{aretenir}

\begin{exercice}[Locating roots]
\begin{enumerate}
\item Show that the equation $x^3 + 2x - 5 = 0$ has exactly one real root and that it lies between $1$ and $2$. Locate it between two consecutive tenths.
\item By sketching $y = \ln x$ and $y = 2 - x$, show that $\ln x = 2 - x$ has exactly one root, and find integer bounds for it.
\item Explain why the sign change of $f(x) = \dfrac{x+1}{x-2}$ between $x = 1$ and $x = 3$ does not guarantee a root of $f(x) = 0$ in $(1, 3)$. Where is the actual root?
\item Rearrange $x^2 - 5x + 3 = 0$ in the form $x = g(x)$ in two different ways and test, with $x_0 = 4$, which rearrangement seems to converge.
\end{enumerate}
\end{exercice}

\begin{corrige}[Exercise 1]
\textbf{1.} $f(x) = x^3 + 2x - 5$ is continuous; $f(1) = -2 < 0$, $f(2) = 7 > 0$, so a root lies in $(1,2)$. Since $f'(x) = 3x^2 + 2 > 0$ for all $x$, $f$ is strictly increasing, so the root is unique. Refining: $f(1.3) \approx -0.203 < 0$, $f(1.4) \approx 0.544 > 0$, so the root lies in $(1.3, 1.4)$.

\textbf{2.} $y = \ln x$ is increasing from $-\infty$, $y = 2 - x$ is a decreasing line; they cross exactly once, for $x \in (1, 2)$ since at $x=1$: $\ln 1 = 0 < 1$; at $x = 2$: $\ln 2 \approx 0.693 > 0$. With $f(x) = \ln x + x - 2$: $f(1) = -1 < 0$, $f(2) \approx 0.693 > 0$. Root in $(1,2)$.

\textbf{3.} $f$ is \emph{not continuous} on $[1,3]$: it has a vertical asymptote at $x = 2$. The sign change is a false positive caused by the discontinuity. The actual root satisfies $x + 1 = 0$, i.e. $x = -1$, which is outside $(1,3)$.

\textbf{4.} (i) $x = \dfrac{x^2 + 3}{5}$: $x_0 = 4 \to x_1 = 3.8 \to x_2 \approx 3.488 \to \ldots$ converging towards $4.303$? Checking: the roots are $\dfrac{5 \pm \sqrt{13}}{2} \approx 4.303$ and $0.697$; the values decrease towards $4.303$ — convergence. (ii) $x = 5 - \dfrac{3}{x}$: $x_0 = 4 \to x_1 = 4.25 \to x_2 \approx 4.294 \to x_3 \approx 4.301$ — also convergent to $4.303$. (Other rearrangements such as $x = \sqrt{5x - 3}$ behave differently; the point is that convergence depends on the rearrangement.)
\end{corrige}

\coursec{The Bisection Method}

Having learned how to locate a root of a continuous function $f(x)$ by finding an interval $[a,b]$ where $f(a)f(b)<0$, we now need a systematic procedure to \emph{approximate} that root to any desired accuracy. The simplest and most robust such procedure is the \textbf{bisection method} (also called the \emph{interval-halving method}). It repeatedly halves the interval, always keeping the subinterval where the sign change occurs, thereby trapping the root in an ever-shrinking interval.

\courssub{Principle of the Bisection Method}

Suppose $f$ is continuous on $[a,b]$ and $f(a)f(b)<0$. By the Intermediate Value Theorem, there is at least one root $\alpha \in (a,b)$. The idea is straightforward:
\begin{enumerate}
    \item Compute the midpoint $m = \dfrac{a+b}{2}$.
    \item Evaluate $f(m)$.
    \item If $f(m)=0$, we have found the exact root (rare in practice).
    \item If $f(a)f(m)<0$, the root lies in $[a,m]$; set $b=m$.
    \item If $f(m)f(b)<0$, the root lies in $[m,b]$; set $a=m$.
\end{enumerate}
In either case, the new interval is exactly half the length of the previous one and still brackets the root. Repeating this process produces a sequence of nested intervals $[a_n,b_n]$ whose lengths tend to zero, forcing both endpoints to converge to the root $\alpha$.

\begin{definition}[Bisection Algorithm]
Given a continuous function $f$ on $[a,b]$ with $f(a)f(b)<0$ and a tolerance $\varepsilon>0$:
\begin{enumerate}
    \item Set $n=0$, $a_0=a$, $b_0=b$.
    \item While $b_n - a_n > \varepsilon$ (or until endpoints agree to required decimal places):
        \begin{itemize}
            \item Compute $m_n = \dfrac{a_n+b_n}{2}$.
            \item If $f(a_n)f(m_n) < 0$, set $a_{n+1}=a_n$, $b_{n+1}=m_n$.
            \item Else, set $a_{n+1}=m_n$, $b_{n+1}=b_n$.
        \end{itemize}
    \item The approximate root is $m_n$ (or $\frac{a_n+b_n}{2}$) with error $\le \frac{b_n-a_n}{2}$.
\end{enumerate}
\end{definition}

\begin{methode}[Step-by-Step Bisection Procedure]
\begin{enumerate}
    \item \textbf{Verify the bracket:} Ensure $f$ is continuous on $[a,b]$ and $f(a)f(b)<0$.
    \item \textbf{Draw a table} with columns: $n$, $a_n$, $b_n$, $m_n$, $f(m_n)$, Sign of $f(m_n)$, New Interval.
    \item \textbf{Iterate:} Compute $m_n$, evaluate $f(m_n)$, check its sign against $f(a_n)$ (or $f(b_n)$), and write the new interval.
    \item \textbf{Stop:} When the interval width $b_n-a_n$ is smaller than the required tolerance, or when $a_n$ and $b_n$ agree to the required number of decimal places.
    \item \textbf{State the root:} The best estimate is the midpoint of the final interval. Give the answer to the requested accuracy.
\end{enumerate}
\end{methode}

\courssub{Worked Example: Solving $x^3 - 7x + 3 = 0$ in $[0,1]$}

Let $f(x)=x^3-7x+3$. We have $f(0)=3>0$ and $f(1)=-3<0$, so a root lies in $[0,1]$. We perform bisection until the root is determined correct to \textbf{two decimal places} (i.e., interval width $<0.005$ or endpoints agree to 2 d.p.).

\begin{exemplebox}[Bisection Iterations for $x^3-7x+3=0$]
\begin{center}
\begin{tabular}{|c|c|c|c|c|c|c|}
\hline
$n$ & $a_n$ & $b_n$ & $m_n=\frac{a_n+b_n}{2}$ & $f(m_n)$ & Sign & New Interval $[a_{n+1},b_{n+1}]$ \\ \hline
0 & 0 & 1 & 0.5 & $0.125-3.5+3=-0.375$ & $-$ & $[0,0.5]$ \\ \hline
1 & 0 & 0.5 & 0.25 & $0.015625-1.75+3=1.265625$ & $+$ & $[0.25,0.5]$ \\ \hline
2 & 0.25 & 0.5 & 0.375 & $0.052734-2.625+3=0.427734$ & $+$ & $[0.375,0.5]$ \\ \hline
3 & 0.375 & 0.5 & 0.4375 & $0.083740-3.0625+3=0.021240$ & $+$ & $[0.4375,0.5]$ \\ \hline
4 & 0.4375 & 0.5 & 0.46875 & $0.103027-3.28125+3=-0.178223$ & $-$ & $[0.4375,0.46875]$ \\ \hline
5 & 0.4375 & 0.46875 & 0.453125 & $0.093018-3.171875+3=-0.078857$ & $-$ & $[0.4375,0.453125]$ \\ \hline
6 & 0.4375 & 0.453125 & 0.4453125 & $0.088287-3.117188+3=-0.028901$ & $-$ & $[0.4375,0.4453125]$ \\ \hline
7 & 0.4375 & 0.4453125 & 0.44140625 & $0.085999-3.089844+3=-0.003845$ & $-$ & $[0.4375,0.44140625]$ \\ \hline
8 & 0.4375 & 0.44140625 & 0.439453125 & $0.084869-3.076172+3=0.008697$ & $+$ & $[0.439453125,0.44140625]$ \\ \hline
\end{tabular}
\end{center}
\end{exemplebox}

After 8 iterations, the interval is $[0.439453125,\;0.44140625]$. Its width is $0.001953125 < 0.005$. Both endpoints round to \textbf{0.44} to two decimal places. The midpoint of the final interval is $0.4404296875 \approx 0.44$.

\begin{aretenir}[Final Answer]
The root of $x^3-7x+3=0$ in $[0,1]$ correct to two decimal places is $\boldsymbol{\alpha = 0.44}$.
\end{aretenir}

\begin{popfigure}
\begin{tikzpicture}[scale=1.2]
\draw[thick,->] (-0.2,0) -- (1.2,0) node[right] {$x$};
\foreach \x in {0,0.25,0.5,0.75,1} \draw (\x,0.05) -- (\x,-0.05) node[below] {\small $\x$};
% Intervals
\draw[popblue, very thick] (0,0.3) -- (1,0.3) node[left] {$[a_0,b_0]$};
\draw[poporange, very thick] (0,0.6) -- (0.5,0.6) node[left] {$[a_1,b_1]$};
\draw[popgreen, very thick] (0.25,0.9) -- (0.5,0.9) node[left] {$[a_2,b_2]$};
\draw[poppurple, very thick] (0.375,1.2) -- (0.5,1.2) node[left] {$[a_3,b_3]$};
\draw[poppink, very thick] (0.4375,1.5) -- (0.5,1.5) node[left] {$[a_4,b_4]$};
% Midpoints
\foreach \m/\c in {0.5/popblue, 0.25/poporange, 0.375/popgreen, 0.4375/poppurple, 0.46875/poppink} {
    \draw[\c, dashed] (\m,0) -- (\m,1.6);
    \fill[\c] (\m,1.6) circle (1.5pt);
}
\node at (0.5,1.8) {$m_0$}; \node at (0.25,1.8) {$m_1$}; \node at (0.375,1.8) {$m_2$}; \node at (0.4375,1.8) {$m_3$}; \node at (0.46875,1.8) {$m_4$};
\end{tikzpicture}
\legende{Successive nested intervals in the bisection method. Each midpoint $m_n$ splits the current interval; the half containing the root is kept.}
\end{popfigure}

\courssub{Stopping Criteria and Error Bound}

At each step the interval length is halved: $b_{n+1}-a_{n+1} = \frac{1}{2}(b_n-a_n)$. After $n$ bisections, the length is
\[
b_n - a_n = \frac{b_0 - a_0}{2^n}.
\]
The true root $\alpha$ lies somewhere in $[a_n,b_n]$. If we take the midpoint $m_n$ as our approximation, the maximum possible error is half the interval length.

\begin{propriete}[Error Bound after $n$ Bisections]
If the initial bracket is $[a,b]$ and $n$ bisections have been performed, then for the midpoint $m_n = \frac{a_n+b_n}{2}$,
\[
|\alpha - m_n| \le \frac{b_n - a_n}{2} = \frac{b-a}{2^{n+1}}.
\]
\end{propriete}

This formula allows us to predict \textbf{how many iterations} are needed to guarantee a given accuracy \emph{before} starting the computation.

\begin{methode}[Choosing the Number of Iterations for a Given Accuracy]
To achieve an error $< \varepsilon$ (or interval width $< \delta$):
\begin{enumerate}
    \item For error bound: solve $\dfrac{b-a}{2^{n+1}} < \varepsilon$ for the smallest integer $n$.
    \item For interval width: solve $\dfrac{b-a}{2^n} < \delta$ for the smallest integer $n$.
    \item Equivalently: $n > \dfrac{\ln(b-a) - \ln(\delta)}{\ln 2}$ (width) or $n > \dfrac{\ln(b-a) - \ln(2\varepsilon)}{\ln 2}$ (error).
\end{enumerate}
\end{methode}

\begin{exemplebox}[Predicting Iterations]
For $f(x)=x^3-7x+3$ on $[0,1]$, we want the root correct to \textbf{3 decimal places} (error $< 0.0005$).
Initial width $b-a=1$. We need $\frac{1}{2^{n+1}} < 0.0005 \Rightarrow 2^{n+1} > 2000 \Rightarrow n+1 > \log_2 2000 \approx 10.97$.
So $n \ge 10$ iterations guarantee the required accuracy. (In practice, we stop as soon as the endpoints agree to 3 d.p., which may happen slightly earlier.)
\end{exemplebox}

\begin{popfigure}
\begin{tikzpicture}
\begin{axis}[
    width=12cm, height=7cm,
    axis lines=middle,
    xlabel=$x$, ylabel=$y$,
    xmin=-0.1, xmax=1.1,
    ymin=-3.5, ymax=3.5,
    xtick={0,0.25,0.5,0.75,1},
    ytick={-3,-2,-1,0,1,2,3},
    grid=major, grid style={popline},
    clip=false
]
\addplot[popblue, thick, domain=0:1, samples=100] {x^3 - 7*x + 3};
\addplot[only marks, mark=*, mark size=2, poporange] coordinates {(0,3) (1,-3)};
\node[poporange, anchor=south west] at (axis cs:0,3) {$f(0)=3$};
\node[poporange, anchor=north west] at (axis cs:1,-3) {$f(1)=-3$};
% First midpoint m0=0.5
\addplot[only marks, mark=*, mark size=2.5, popgreen] coordinates {(0.5,-0.375)};
\node[popgreen, anchor=south] at (axis cs:0.5,-0.375) {$m_0$};
% Second midpoint m1=0.25
\addplot[only marks, mark=*, mark size=2.5, poppurple] coordinates {(0.25,1.265625)};
\node[poppurple, anchor=south] at (axis cs:0.25,1.265625) {$m_1$};
% Shade retained interval after 2 steps: [0.25,0.5]
\draw[popgreenL, opacity=0.3, fill=popgreenL] (axis cs:0.25,0) rectangle (axis cs:0.5,0.15);
\end{axis}
\end{tikzpicture}
\legende{Graph of $f(x)=x^3-7x+3$ showing the first two midpoints $m_0=0.5$ and $m_1=0.25$. The shaded region $[0.25,0.5]$ is the interval retained after two bisections.}
\end{popfigure}

\courssub{Advantages and Disadvantages}

\begin{center}
\begin{tabularx}{\linewidth}{|l|X|X|}
\hline
\rowcolor{popblueL}
\textbf{Aspect} & \textbf{Advantages} & \textbf{Disadvantages} \\ \hline
\textbf{Convergence} & \textbf{Guaranteed} for any continuous $f$ with a sign change. No divergence possible. & \textbf{Slow} (linear convergence). Error halves each step; $\sim 3.3$ iterations per decimal digit. \\ \hline
\textbf{Simplicity} & Very easy to program/implement; only requires function evaluations (no derivative). & Requires a \textbf{valid starting bracket} $[a,b]$ with $f(a)f(b)<0$. \\ \hline
\textbf{Robustness} & Unaffected by function shape (flat regions, multiple roots in interval — finds one). & Cannot find roots of even multiplicity (no sign change, e.g., $f(x)=x^2$ at $0$). \\ \hline
\textbf{Error Control} & Exact error bound known at every step: $\frac{b_n-a_n}{2}$. & No way to accelerate; always one bit per iteration. \\ \hline
\end{tabularx}
\end{center}

\courssub{GCE Examination Technique}

In the GCE Advanced Level examination, you will often be asked to:
\begin{itemize}
    \item Perform a fixed number of bisections (e.g., "Carry out 4 iterations").
    \item Find the root to a stated accuracy (e.g., "correct to 3 decimal places").
    \item State the number of iterations needed to achieve a given tolerance.
\end{itemize}

\begin{attention}[Common Mistakes to Avoid]
\begin{enumerate}
    \item \textbf{Sign errors in $f(m)$:} One wrong sign sends you to the wrong half-interval, and all subsequent rows are meaningless. \textbf{Double-check every $f(m)$ calculation.}
    \item \textbf{Using an endpoint as the final answer:} The best estimate is the \textbf{midpoint of the final interval}, not $a_n$ or $b_n$. The midpoint halves the maximum error.
    \item \textbf{Stopping too early:} "Correct to 2 d.p." means the interval width must be $<0.005$ \textbf{or} the endpoints must round to the same 2 d.p. value. Do not stop just because $f(m)$ is small.
    \item \textbf{Messy tables:} Draw a neat table with a ruler. Show $a_n$, $b_n$, $m_n$, $f(m_n)$, and the new interval clearly. Examiners award marks for the table itself.
    \item \textbf{Forgetting to verify the initial bracket:} Always state $f(a)$ and $f(b)$ and note $f(a)f(b)<0$ before starting.
\end{enumerate}
\end{attention}

\begin{attention}[The "Midpoint vs Endpoint" Trap]
Suppose after several iterations you have $[0.439, 0.441]$. The root is $0.44$ to 2 d.p.
\begin{itemize}
    \item Midpoint $= 0.440$ $\to$ rounds to $0.44$. \textbf{Correct.}
    \item Left endpoint $= 0.439$ $\to$ rounds to $0.44$. (Works here but not guaranteed.)
    \item Right endpoint $= 0.441$ $\to$ rounds to $0.44$. (Works here but not guaranteed.)
\end{itemize}
If the interval were $[0.439, 0.440]$, midpoint $=0.4395 \to 0.44$, but right endpoint $=0.440 \to 0.44$, left $=0.439 \to 0.44$. Still okay.
But if interval were $[0.434, 0.436]$ (width 0.002), midpoint $=0.435 \to 0.44$, endpoints give $0.43$ and $0.44$ — \textbf{disagree}. Always use the midpoint.
\end{attention}

\begin{exoresolu}[GCE Style Question]
\textbf{Statement:} The function $f(x)=e^x - 3x$ has a root $\alpha$ in the interval $[0,1]$.
\begin{enumerate}
    \item Verify that $\alpha$ lies in $[0,1]$.
    \item Perform \textbf{three} iterations of the bisection method to approximate $\alpha$.
    \item Using your results, state the value of $\alpha$ correct to \textbf{one decimal place}, giving a reason for your answer.
\end{enumerate}

\tcblower
\textbf{Detailed Solution:}

\textbf{1. Bracket verification:}
$f(0)=e^0-0=1 > 0$.
$f(1)=e-3 \approx 2.71828-3 = -0.28172 < 0$.
$f$ is continuous (exponential and polynomial), and $f(0)f(1)<0$. Hence a root $\alpha \in (0,1)$.

\textbf{2. Three iterations:}
We need $e^x$ values: $e^{0.5}\approx1.64872$, $e^{0.75}\approx2.11700$, $e^{0.625}\approx1.86825$.

\begin{center}
\begin{tabular}{|c|c|c|c|c|c|c|}
\hline
$n$ & $a_n$ & $b_n$ & $m_n$ & $f(m_n)$ & Sign & New Interval \\ \hline
0 & 0 & 1 & 0.5 & $1.64872 - 1.5 = 0.14872$ & $+$ & $[0.5, 1]$ \\ \hline
1 & 0.5 & 1 & 0.75 & $2.11700 - 2.25 = -0.13300$ & $-$ & $[0.5, 0.75]$ \\ \hline
2 & 0.5 & 0.75 & 0.625 & $1.86825 - 1.875 = -0.00675$ & $-$ & $[0.5, 0.625]$ \\ \hline
\end{tabular}
\end{center}

\textbf{3. Root to 1 d.p.:}
After 3 iterations, the interval is $[0.5, 0.625]$. Width $=0.125$.
The midpoint is $m_3 = 0.5625$.
To 1 d.p., the endpoints are $0.5$ and $0.6$ (since $0.625$ rounds to $0.6$). They do \textbf{not} agree.
The required interval width for 1 d.p. certainty is $<0.05$. Here $0.125 > 0.05$.
\textbf{Conclusion:} With only three iterations, the root \textbf{cannot} be stated reliably to one decimal place. More iterations are needed until the interval width is less than $0.05$ (or endpoints agree to 1 d.p.).
\end{exoresolu}

\begin{savaistu}[Historical Note]
The bisection method is essentially a binary search on a continuous domain. It was known to ancient mathematicians in geometric form (e.g., Archimedes' method of exhaustion for $\pi$). In the 19th century, Bernard Bolzano formalised the Intermediate Value Theorem, giving the method a rigorous foundation. Today, it is the "fail-safe" fallback in many numerical libraries (like Brent's method) because it never diverges.
\end{savaistu}

\begin{exercice}[Practice Problems]
\begin{enumerate}
    \item Use the bisection method to find the root of $x^3 - 2x - 5 = 0$ in $[2,3]$ correct to 2 decimal places. Show your table.
    \item How many bisections are needed to guarantee an error $< 10^{-4}$ for a root initially bracketed in $[1,4]$?
    \item The function $f(x)=\sin x - x/2$ has a root in $[\pi/2, \pi]$. Perform 4 iterations of bisection. What is the best estimate and its maximum error?
    \item Explain why the bisection method cannot find the root of $f(x)=x^2$ at $x=0$, even though $f(0)=0$.
\end{enumerate}
\end{exercice}

\begin{corrige}[Exercise 1]
$f(x)=x^3-2x-5$. $f(2)=8-4-5=-1 < 0$, $f(3)=27-6-5=16 > 0$. Root in $[2,3]$.
Target: 2 d.p. $\Rightarrow$ width $<0.005$ or endpoints agree to 2 d.p.
Iterations (values to 6 d.p.):
\begin{center}
\begin{tabular}{|c|c|c|c|c|c|c|}
\hline
$n$ & $a_n$ & $b_n$ & $m_n$ & $f(m_n)$ & Sign & New Interval \\ \hline
0 & 2 & 3 & 2.5 & 5.625 & $+$ & $[2, 2.5]$ \\ \hline
1 & 2 & 2.5 & 2.25 & 1.890625 & $+$ & $[2, 2.25]$ \\ \hline
2 & 2 & 2.25 & 2.125 & 0.341797 & $+$ & $[2, 2.125]$ \\ \hline
3 & 2 & 2.125 & 2.0625 & -0.345703 & $-$ & $[2.0625, 2.125]$ \\ \hline
4 & 2.0625 & 2.125 & 2.09375 & -0.009277 & $-$ & $[2.09375, 2.125]$ \\ \hline
5 & 2.09375 & 2.125 & 2.109375 & 0.165253 & $+$ & $[2.09375, 2.109375]$ \\ \hline
6 & 2.09375 & 2.109375 & 2.1015625 & 0.077408 & $+$ & $[2.09375, 2.1015625]$ \\ \hline
7 & 2.09375 & 2.1015625 & 2.09765625 & 0.033926 & $+$ & $[2.09375, 2.09765625]$ \\ \hline
8 & 2.09375 & 2.09765625 & 2.095703125 & 0.012278 & $+$ & $[2.09375, 2.095703125]$ \\ \hline
9 & 2.09375 & 2.095703125 & 2.0947265625 & 0.001494 & $+$ & $[2.09375, 2.0947265625]$ \\ \hline
10 & 2.09375 & 2.0947265625 & 2.09423828125 & -0.003894 & $-$ & $[2.09423828125, 2.0947265625]$ \\ \hline
\end{tabular}
\end{center}
After 10 iterations, interval width $= 0.000488 < 0.005$. Endpoints $2.0942$ and $2.0947$ both round to $2.09$.
\textbf{Root $\approx 2.09$ to 2 d.p.}
\end{corrige}

\begin{corrige}[Exercise 2]
Initial width $b-a = 3$. Error bound after $n$ iterations: $\frac{3}{2^{n+1}} < 10^{-4}$.
$2^{n+1} > 30000$. $2^{14}=16384$, $2^{15}=32768$. So $n+1 \ge 15 \Rightarrow n \ge 14$.
\textbf{14 iterations} guarantee error $< 10^{-4}$.
\end{corrige}

\begin{corrige}[Exercise 3]
$f(x)=\sin x - x/2$. $[\pi/2, \pi] \approx [1.5708, 3.1416]$.
$f(\pi/2)=1 - \pi/4 \approx 0.2146 >0$. $f(\pi)=0 - \pi/2 \approx -1.5708 <0$.
\begin{center}
\begin{tabular}{|c|c|c|c|c|c|c|}
\hline
$n$ & $a_n$ & $b_n$ & $m_n$ & $f(m_n)$ & Sign & New Interval \\ \hline
0 & $\pi/2$ & $\pi$ & $3\pi/4\approx2.3562$ & $-0.4710$ & $-$ & $[\pi/2, 3\pi/4]$ \\ \hline
1 & $\pi/2$ & $3\pi/4$ & $5\pi/8\approx1.9635$ & $-0.0578$ & $-$ & $[\pi/2, 5\pi/8]$ \\ \hline
2 & $\pi/2$ & $5\pi/8$ & $9\pi/16\approx1.7671$ & $0.0972$ & $+$ & $[9\pi/16, 5\pi/8]$ \\ \hline
3 & $9\pi/16$ & $5\pi/8$ & $19\pi/32\approx1.8653$ & $0.0242$ & $+$ & $[19\pi/32, 5\pi/8]$ \\ \hline
\end{tabular}
\end{center}
Best estimate: midpoint of final interval $= \frac{19\pi/32 + 5\pi/8}{2} = \frac{39\pi}{64} \approx 1.9144$.
Maximum error $= \frac{5\pi/8 - 19\pi/32}{2} = \frac{\pi/32}{2} = \frac{\pi}{64} \approx 0.0491$.
\end{corrige}

\begin{corrige}[Exercise 4]
$f(x)=x^2$. $f(0)=0$. For any interval $[a,b]$ containing $0$ but not having $0$ as an endpoint, $f(a)>0$ and $f(b)>0$ (since $x^2>0$ for $x\neq0$). Thus $f(a)f(b)>0$, no sign change. The Intermediate Value Theorem condition fails. The method relies on a sign change to bracket the root; a root of even multiplicity (touching the axis) does not produce a sign change, so bisection cannot "see" it.
\end{corrige}

\coursec{The Newton–Raphson Method}

The bisection method of the previous section is safe and reliable: it never fails when a sign change has been located. But it is \emph{slow} — each step gains only one binary digit, and reaching four decimal places typically requires more than a dozen iterations. We now study a much faster method, due to Isaac Newton and Joseph Raphson, which uses not only the \emph{values} of $f$ but also its \emph{slope}. When it works, it is spectacularly efficient; when it fails, it fails badly — so we must also learn when and why it can go wrong.

\courssub{The Geometric Idea}

Suppose we want to solve $f(x)=0$ and we have an estimate $x_0$ of the root $\alpha$. Instead of working with the curve $y=f(x)$ itself, which may be complicated, we replace it near $x_0$ by the simplest curve we fully understand: the \cle{tangent line} to the graph at the point $(x_0, f(x_0))$. The tangent is a good local approximation of the curve, so the point where the tangent crosses the $x$-axis should be close to the point where the curve crosses the $x$-axis — that is, close to the root. We call this intercept $x_1$, then repeat the process: draw the tangent at $(x_1, f(x_1))$, take its intercept $x_2$, and so on.

\begin{popfigure}
\begin{center}
\begin{tikzpicture}
\begin{axis}[
  width=0.86\linewidth, height=7cm,
  axis lines=middle, xlabel=$x$, ylabel=$y$,
  xmin=1.8, xmax=4.4, ymin=-4, ymax=10,
  xtick={2.6458,2.875,4}, xticklabels={$\alpha$,$x_1$,$x_0$},
  ytick=\empty, samples=100, domain=1.8:4.4,
  legend style={at={(0.03,0.97)}, anchor=north west, font=\small}
]
\addplot[popblue, thick] {x^2-7};
\addlegendentry{$y=x^2-7$}
\addplot[poporange, thick, domain=2.2:4.3] {8*x-23};
\addlegendentry{tangent at $x_0=4$}
\addplot[poppurple, thick, domain=2.4:3.4] {5.75*x-15.2656};
\addlegendentry{tangent at $x_1=2.875$}
\addplot[only marks, mark=*, popdark] coordinates {(4,9) (2.875,1.2656)};
\draw[dashed, gray] (axis cs:4,0) -- (axis cs:4,9);
\draw[dashed, gray] (axis cs:2.875,0) -- (axis cs:2.875,1.2656);
\node[popdark, font=\small] at (axis cs:4.25,9.3) {$(x_0,f(x_0))$};
\node[popdark, font=\small] at (axis cs:3.35,2.6) {$(x_1,f(x_1))$};
\end{axis}
\end{tikzpicture}
\end{center}
\legende{Newton--Raphson applied to $f(x)=x^2-7$. The tangent at $x_0=4$ meets the $x$-axis at $x_1=2.875$; the tangent there meets it at $x_2\approx 2.6549$, already very close to $\alpha=\sqrt{7}\approx 2.6458$.}
\end{popfigure}

\courssub{Derivation of the Formula}

The tangent line to $y=f(x)$ at the point $(x_n, f(x_n))$ has gradient $f'(x_n)$, so its equation is
\[ y - f(x_n) = f'(x_n)\,(x - x_n). \]
The next estimate $x_{n+1}$ is where this tangent crosses the $x$-axis, i.e. where $y=0$:
\[ -f(x_n) = f'(x_n)\,(x_{n+1} - x_n)
\quad\Longrightarrow\quad
x_{n+1} - x_n = -\frac{f(x_n)}{f'(x_n)}. \]

\begin{propriete}[Newton--Raphson iteration]
To solve $f(x)=0$, choose a starting value $x_0$ and iterate
\[ \boxed{\;x_{n+1} = x_n - \frac{f(x_n)}{f'(x_n)}\;} \qquad (n = 0, 1, 2, \dots) \]
provided $f'(x_n)\neq 0$. The derivation from the tangent line is examinable and takes only two lines, as shown above.
\end{propriete}

\begin{methode}[Applying the Newton--Raphson method]
\begin{enumerate}
\item Write the equation in the form $f(x)=0$ and compute $f'(x)$ \textbf{carefully}.
\item Choose $x_0$ sensibly: the midpoint (or an endpoint) of a sign-change interval, or a value read from a graph of $y=f(x)$.
\item Compute $x_1, x_2, \dots$ with the formula, keeping \textbf{at least two more decimal places} than the accuracy required.
\item Stop when two successive iterates agree to the required accuracy, and \textbf{verify} the answer (e.g. by a sign change of $f$ across the claimed rounded value).
\end{enumerate}
\end{methode}

\begin{attention}[Differentiate first, differentiate right]
Every iterate depends on $f'$. If your derivative is wrong, \emph{every single value} you compute afterwards is wrong, even if the arithmetic is perfect. Before iterating, check $f'$ on a simple value of $x$.
\end{attention}

\begin{attention}[Radians!]
Whenever $f$ involves trigonometric functions, your calculator \textbf{must be in radian mode}. The derivative of $\cos x$ is $-\sin x$ only when $x$ is measured in radians.
\end{attention}

\courssub{Worked Examples}

\begin{exoresolu}[Computing $\sqrt{7}$ to 4 decimal places]
Solve $x^2-7=0$ by Newton--Raphson and hence find $\sqrt{7}$ to 4 d.p.
\tcblower
$f(x)=x^2-7$, so $f'(x)=2x$ and the iteration simplifies to
\[ x_{n+1}=x_n-\frac{x_n^2-7}{2x_n}=\frac{1}{2}\!\left(x_n+\frac{7}{x_n}\right). \]
Since $f(2)=-3<0$ and $f(3)=2>0$, the root lies in $(2,3)$; take $x_0=3$.
\begin{center}
\begin{tabular}{|c|c|c|c|}
\hline
\rowcolor{popblueL} $n$ & $x_n$ & $f(x_n)$ & $x_{n+1}$ \\ \hline
0 & 3.000000 & 2.000000 & 2.666667 \\ \hline
1 & 2.666667 & 0.111111 & 2.645833 \\ \hline
2 & 2.645833 & 0.000434 & 2.645751 \\ \hline
3 & 2.645751 & $-0.0000001$ & 2.645751 \\ \hline
\end{tabular}
\end{center}
$x_3$ and $x_4$ (both $2.645751$) agree to 4 d.p., so $\sqrt{7}\approx \cle{2.6458}$. Note the speed: four iterations suffice, where bisection would need about thirteen.
\end{exoresolu}

\begin{exoresolu}[Solving $e^x-3x=0$]
Show that $e^x-3x=0$ has a root near $1.5$ and find it to 4 d.p.
\tcblower
$f(x)=e^x-3x$; $f(1.5)=4.4817-4.5=-0.0183<0$ and $f(1.6)=4.9530-4.8=0.1530>0$: sign change, so a root lies in $(1.5,1.6)$. Take $x_0=1.5$. With $f'(x)=e^x-3$:
\begin{center}
\begin{tabular}{|c|c|c|c|c|}
\hline
\rowcolor{popblueL} $n$ & $x_n$ & $f(x_n)$ & $f'(x_n)$ & $x_{n+1}$ \\ \hline
0 & 1.500000 & $-0.018311$ & 1.481689 & 1.512358 \\ \hline
1 & 1.512358 & 0.000336 & 1.537410 & 1.512139 \\ \hline
2 & 1.512139 & $-0.000001$ & 1.537135 & 1.512139 \\ \hline
\end{tabular}
\end{center}
The iterates agree to 4 d.p.: the root is $\cle{1.5121}$. (The equation has a second root near $0.6190$; a different $x_0$ finds it.)
\end{exoresolu}

\begin{exoresolu}[Solving $\cos x = x$]
Find, to 4 d.p., the root of $\cos x - x = 0$ (calculator in \textbf{radians}).
\tcblower
$f(x)=\cos x - x$, $f'(x)=-\sin x - 1$. Since $f(0)=1>0$ and $f(1)=-0.4597<0$, take $x_0=0.75$.
\begin{center}
\begin{tabular}{|c|c|c|c|c|}
\hline
\rowcolor{popblueL} $n$ & $x_n$ & $f(x_n)$ & $f'(x_n)$ & $x_{n+1}$ \\ \hline
0 & 0.750000 & $-0.018311$ & $-1.681639$ & 0.739111 \\ \hline
1 & 0.739111 & $-0.000026$ & $-1.673612$ & 0.739095 \\ \hline
2 & 0.739085 & $0.000000$ & $-1.673577$ & 0.739085 \\ \hline
\end{tabular}
\end{center}
The root is $\cle{0.7391}$ (4 d.p.). This number is the unique fixed point of $\cos$: pressing \textsf{cos} repeatedly on a radian-mode calculator converges to it!
\end{exoresolu}

\courssub{Speed of Convergence}

\begin{propriete}[Quadratic convergence]
Near a \cle{simple root} $\alpha$ (i.e. $f(\alpha)=0$, $f'(\alpha)\neq 0$), with a good starting value, Newton--Raphson converges \emph{quadratically}: the error roughly squares at each step, so the number of correct decimal places approximately \textbf{doubles} at each iteration. (Proof not required at GCE.)
\end{propriete}

Look at the $\sqrt{7}$ table: the errors are roughly $3.5\times10^{-1}$, $2.1\times10^{-2}$, $8\times10^{-5}$, $10^{-9}$ — each error is proportional to the \emph{square} of the previous one. This is why Newton--Raphson is the standard method built into calculators and computers.

\begin{savaistu}[Historical note]
Isaac Newton described a version of this method in \emph{De analysi per aequationes numero terminorum infinitas} (1669), but only for polynomials and without the general derivative formula. Joseph Raphson published the general iterative form $x_{n+1}=x_n-f(x_n)/f'(x_n)$ in \emph{Analysis aequationum universalis} (1690). The method is therefore rightly named after both.
\end{savaistu}

\courssub{When Newton--Raphson Fails}

\begin{attention}[Failure cases — always check convergence]
\begin{itemize}
\item \textbf{Horizontal tangent:} if $f'(x_n)=0$ (or is tiny), the formula divides by (nearly) zero and $x_{n+1}$ is undefined or thrown far away. This happens near stationary points.
\item \textbf{Poor $x_0$:} starting far from the root can send the iteration to a \emph{different} root, into a cycle, or off to infinity.
\item \textbf{Oscillation:} for $f(x)=x^3-2x+2$ with $x_0=0$, one gets $x_1=1$, then $x_2=0$, $x_3=1,\dots$ — an endless loop that never converges.
\item \textbf{Roots at inflection points} (e.g. $f(x)=x^3$ at $0$, where $f'(0)=0$): convergence, when it occurs, is no longer quadratic.
\end{itemize}
\textbf{Habit:} watch the sequence of iterates. If they do not settle down, change $x_0$ or fall back on bisection.
\end{attention}

\begin{popfigure}
\begin{center}
\begin{tikzpicture}
\begin{axis}[
  width=0.8\linewidth, height=6.5cm,
  axis lines=middle, xlabel=$x$, ylabel=$y$,
  xmin=-2.4, xmax=2.4, ymin=-2, ymax=4,
  xtick={-2,-1,0,1,2}, ytick=\empty,
  samples=120, domain=-2.2:2.2,
  legend style={at={(0.03,0.97)}, anchor=north west, font=\small}
]
\addplot[popblue, thick] {x^3-2*x+2};
\addlegendentry{$y=x^3-2x+2$}
\addplot[poporange, thick, domain=-1.2:1.6] {2-2*x};
\addlegendentry{tangent at $x_0=0$}
\addplot[poppurple, thick, domain=-0.6:2.2] {x};
\addlegendentry{tangent at $x_1=1$}
\addplot[only marks, mark=*, popdark] coordinates {(0,2) (1,1)};
\draw[dashed, gray] (axis cs:0,0) -- (axis cs:0,2);
\draw[dashed, gray] (axis cs:1,0) -- (axis cs:1,1);
\node[popdark, font=\small] at (axis cs:0.15,2.4) {$x_0=0$};
\node[popdark, font=\small] at (axis cs:1.3,0.7) {$x_1=1$};
\end{axis}
\end{tikzpicture}
\end{center}
\legende{Failure by oscillation: for $f(x)=x^3-2x+2$, the tangent at $x_0=0$ gives $x_1=1$, and the tangent at $x_1=1$ sends us back to $x_2=0$. The iteration cycles forever and never reaches the real root (near $-1.77$).}
\end{popfigure}

\courssub{Link with Fixed-Point Iteration (Enrichment)}

Newton--Raphson is a special case of \cle{fixed-point iteration} $x_{n+1}=g(x_n)$, which solves $x=g(x)$. Indeed, the Newton formula is $x_{n+1}=g(x_n)$ with
\[ g(x)=x-\frac{f(x)}{f'(x)}, \]
and one can check that $g'(\alpha)=0$ at a simple root — the strongest possible case of the following rule.

\begin{propriete}[Convergence condition (enrichment)]
The iteration $x_{n+1}=g(x_n)$ converges to a fixed point $\alpha$ if $x_0$ is close enough to $\alpha$ and $|g'(x)|<1$ near $\alpha$. If $|g'(\alpha)|>1$, it diverges. (Flagged as enrichment for strong candidates.)
\end{propriete}

Graphically, we plot $y=g(x)$ and $y=x$; the iteration appears as a \emph{staircase} or \emph{cobweb}: from $x_n$ go vertically to the curve, then horizontally to the line $y=x$ — that gives $x_{n+1}$.

\begin{popfigure}
\begin{center}
\begin{tikzpicture}
\begin{axis}[
  width=0.46\linewidth, height=6cm,
  axis lines=middle, xlabel=$x$, ylabel=$y$,
  xmin=0, xmax=2.4, ymin=0, ymax=2.4,
  xtick={2}, xticklabels={$\alpha$}, ytick=\empty,
  samples=80, domain=0:2.4, title={Convergence: $g(x)=\tfrac{x}{2}+1$}
]
\addplot[popblue, thick] {0.5*x+1};
\addplot[gray] {x};
\draw[poporange, thick] (axis cs:0.2,0) -- (axis cs:0.2,1.1) -- (axis cs:1.1,1.1) -- (axis cs:1.1,1.55) -- (axis cs:1.55,1.55) -- (axis cs:1.55,1.775) -- (axis cs:1.775,1.775) -- (axis cs:1.775,1.8875);
\addplot[only marks, mark=*, popdark] coordinates {(2,2)};
\end{axis}
\end{tikzpicture}
\hfill
\begin{tikzpicture}
\begin{axis}[
  width=0.46\linewidth, height=6cm,
  axis lines=middle, xlabel=$x$, ylabel=$y$,
  xmin=-1.4, xmax=1.6, ymin=-1.4, ymax=1.6,
  xtick={1}, xticklabels={$\alpha$}, ytick=\empty,
  samples=80, domain=-1.4:1.6, title={Divergence: $g(x)=2x-1$}
]
\addplot[popblue, thick] {2*x-1};
\addplot[gray] {x};
\draw[poporange, thick] (axis cs:0.8,0) -- (axis cs:0.8,0.6) -- (axis cs:0.6,0.6) -- (axis cs:0.6,0.2) -- (axis cs:0.2,0.2) -- (axis cs:0.2,-0.6) -- (axis cs:-0.6,-0.6);
\addplot[only marks, mark=*, popdark] coordinates {(1,1)};
\end{axis}
\end{tikzpicture}
\end{center}
\legende{Cobweb diagrams. Left: $|g'|<1$, the staircase spirals into the fixed point. Right: $|g'|=2>1$, the cobweb moves away — the iteration diverges even though $x_0=0.8$ is close to $\alpha=1$.}
\end{popfigure}

\begin{aretenir}[Key points]
\begin{itemize}
\item Newton--Raphson: $x_{n+1}=x_n-\dfrac{f(x_n)}{f'(x_n)}$ — the next estimate is the $x$-intercept of the tangent at $(x_n,f(x_n))$.
\item Choose $x_0$ from a sign-change interval or a graph; iterate until successive values agree; then verify.
\item Convergence is quadratic near a simple root: correct digits roughly double each step.
\item It fails when $f'(x_n)\approx 0$, when $x_0$ is badly chosen, or when the iteration cycles — always monitor the iterates.
\item It is a fixed-point iteration with $g(x)=x-\frac{f(x)}{f'(x)}$; convergence of $x_{n+1}=g(x_n)$ requires $|g'|<1$ near the root.
\end{itemize}
\end{aretenir}

\begin{exercice}[Practice]
\begin{enumerate}
\item Use Newton--Raphson with $x_0=2$ to find the root of $x^3-3x-1=0$ in $(1,2)$, to 4 d.p.
\item Show that $x^3+x-3=0$ has a root in $(1,2)$ and find it to 4 d.p.
\item Explain, with a sketch, why Newton--Raphson with $x_0=1$ fails for $f(x)=x^3-3x+5$.
\item (Enrichment) For $x_{n+1}=\cos x_n$, compute $g'(x)$ at the fixed point $\alpha\approx0.7391$ and explain why the iteration converges.
\end{enumerate}
\end{exercice}

\begin{corrige}[Exercise 1]
$f(x)=x^3-3x-1$, $f'(x)=3x^2-3$. Root in $(1,2)$ since $f(1)=-3<0$, $f(2)=1>0$. Start with $x_0=2$.
\begin{center}
\begin{tabular}{|c|c|c|c|c|}
\hline
\rowcolor{popblueL} $n$ & $x_n$ & $f(x_n)$ & $f'(x_n)$ & $x_{n+1}$ \\ \hline
0 & 2.000000 & 1.000000 & 9.000000 & 1.888889 \\ \hline
1 & 1.888889 & 0.072480 & 7.703704 & 1.879481 \\ \hline
2 & 1.879481 & 0.000600 & 7.597322 & 1.879402 \\ \hline
3 & 1.879402 & $\approx 0$ & 7.596000 & 1.879402 \\ \hline
\end{tabular}
\end{center}
$x_3$ and $x_4$ agree to 4 d.p. The root is $\cle{1.8794}$.
\end{corrige}

\begin{corrige}[Exercise 2]
$f(x)=x^3+x-3$, $f'(x)=3x^2+1$. $f(1)=-1<0$, $f(2)=7>0$: sign change in $(1,2)$. Take $x_0=1.5$.
\begin{center}
\begin{tabular}{|c|c|c|c|c|}
\hline
\rowcolor{popblueL} $n$ & $x_n$ & $f(x_n)$ & $f'(x_n)$ & $x_{n+1}$ \\ \hline
0 & 1.500000 & 1.875000 & 7.750000 & 1.258065 \\ \hline
1 & 1.258065 & 0.249600 & 5.748100 & 1.214683 \\ \hline
2 & 1.214683 & 0.007000 & 5.426200 & 1.213390 \\ \hline
3 & 1.213390 & 0.000400 & 5.416900 & 1.213316 \\ \hline
4 & 1.213316 & $\approx 0$ & 5.416000 & 1.213316 \\ \hline
\end{tabular}
\end{center}
$x_4$ and $x_5$ agree to 4 d.p. The root is $\cle{1.2133}$ (since $1.213316$ rounds to $1.2133$; checking $f(1.2133)\approx 0.0001$, $f(1.2134)\approx 0.0005$, the root is $1.2133$ to 4 d.p.).
\end{corrige}

\begin{corrige}[Exercise 3]
$f(x)=x^3-3x+5$, $f'(x)=3x^2-3$. At $x_0=1$, $f'(1)=0$. The tangent at $(1,3)$ is horizontal ($y=3$) and never meets the $x$-axis, so $x_1$ is undefined. The method fails immediately.
\end{corrige}

\begin{corrige}[Exercise 4]
$g(x)=\cos x$, $g'(x)=-\sin x$. At the fixed point $\alpha\approx0.7391$, $|g'(\alpha)|=\sin(0.7391)\approx0.6736<1$. The convergence condition holds, so the cobweb spirals into the fixed point — convergence is guaranteed for $x_0$ sufficiently close to $\alpha$.
\end{corrige}

\coursec{Accuracy, Comparison and Applications}

In the previous section we met the Newton--Raphson method, a remarkably fast way of refining a root once we have located it. But speed alone is not enough for the GCE examiner: you must also be able to \emph{state precisely how accurate your answer is}, to \emph{justify} that accuracy, to \emph{choose} the best method for a given problem, and to \emph{interpret} the final number in a real situation. That is the business of this final section. By the end of it you will be able to take a full numerical-methods question from start to finish: locate, iterate, verify, interpret.

\courssub{The Language of Accuracy}

\textbf{Decimal places versus significant figures.}\par
Saying that $x = 2.0946$ is given \cle{correct to 4 decimal places} means that the true value lies in the interval $[2.09455,\, 2.09465)$: we control the \emph{fifth} decimal digit. Saying that $x = 0.003142$ is given \cle{correct to 4 significant figures} means we control the fourth \emph{non-zero-leading} digit. Decimal places measure \emph{absolute} precision; significant figures measure \emph{relative} precision. For roots of moderate size the two often coincide in practice, but examination questions about roots almost always ask for \textbf{decimal places}.

\begin{definition}[Absolute error, tolerance, and ``correct to $k$ decimal places'']
Let $\alpha$ be an exact (unknown) value and $x$ an approximation to it.
\begin{itemize}
\item The \cle{absolute error} is $|x - \alpha|$.
\item A \cle{tolerance} (or error bound) is a number $\varepsilon > 0$ for which we can \emph{guarantee} $|x - \alpha| \leq \varepsilon$.
\item We say $x$ is \cle{correct to $k$ decimal places} if rounding $\alpha$ to $k$ decimal places gives $x$; equivalently,
\[ |x - \alpha| \leq \tfrac{1}{2}\times 10^{-k}. \]
For example, ``correct to 3 d.p.'' means an error of at most $0.0005$.
\end{itemize}
\end{definition}

\begin{attention}[The calculator-display trap]
Writing down $x_4 = 1.5214$ because your calculator shows $1.5213797\ldots$ does \textbf{not} prove that the root is $1.521$ correct to 3 d.p. The iterates of a method may agree with each other while both are still wrong, or may agree only by luck. In the GCE, ``show that the root is $1.521$ correct to 3 decimal places'' \emph{must} be justified by a \textbf{sign-change test} (below) or by an error bound from the method itself. An unjustified rounded value earns no accuracy marks.
\end{attention}

\begin{methode}[Verifying a root correct to $k$ decimal places by the sign-change test]
Suppose you believe the root $\alpha$ of a continuous function $f$ is $x_0$ correct to $k$ decimal places.
\begin{enumerate}
\item Form the \cle{rounded interval} $[a,b] = [x_0 - \tfrac12\times 10^{-k},\, x_0 + \tfrac12\times 10^{-k}]$.
\item Compute $f(a)$ and $f(b)$ (keep plenty of guard digits).
\item If $f(a)$ and $f(b)$ have \textbf{opposite signs}, then by the Intermediate Value Theorem $\alpha \in (a,b)$, so $\alpha$ rounds to $x_0$ at $k$ d.p. \quad $\blacksquare$
\item If the signs are the same, your claimed value is \emph{not} justified: iterate further and test again.
\end{enumerate}
\end{methode}

\begin{exemplebox}[Justifying ``correct to 3 d.p.'']
We believe the root of $f(x) = x^3 - 2x - 5$ is $2.094$ correct to 3 d.p. Test the interval $[2.0935,\, 2.0945]$:
\begin{align*}
f(2.0935) &= 2.0935^3 - 2(2.0935) - 5 \approx -0.0020,\\
f(2.0945) &= 2.0945^3 - 2(2.0945) - 5 \approx +0.0003.
\end{align*}
There is a sign change on $[2.0935, 2.0945]$, and $f$ is continuous, so the root lies inside this interval. Hence $\alpha = 2.094$ correct to 3 d.p. \emph{This short computation is exactly what the examiner wants to see.}
\end{exemplebox}

\courssub{How Many Iterations? Predicting and Checking Convergence}

\textbf{Bisection: an exact count.}\par
With the bisection method started on $[a,b]$, after $n$ iterations the interval has length $(b-a)/2^n$, and the root is trapped inside it. Taking the midpoint as our estimate, the error is at most $\dfrac{b-a}{2^{n+1}}$. To guarantee accuracy $\varepsilon$ we therefore need
\[ \frac{b-a}{2^{n+1}} \leq \varepsilon \quad\Longleftrightarrow\quad n \geq \frac{\ln\!\left(\dfrac{b-a}{2\varepsilon}\right)}{\ln 2}. \]
Each iteration buys roughly one extra binary digit; about $3.3$ iterations are needed per decimal digit. This is why bisection is \emph{predictable but slow}.

\begin{exemplebox}[Counting bisection steps]
A root of $f(x)=x^3+x-3$ lies in $[1,2]$. How many bisection iterations guarantee the midpoint is within $0.0005$ of the root (i.e.\ 3 d.p.\ accuracy)?
\[ n \geq \frac{\ln\!\left(\dfrac{1}{2\times 0.0005}\right)}{\ln 2} = \frac{\ln 1000}{\ln 2} \approx \frac{6.9078}{0.6931} \approx 9.97. \]
So $n = 10$ iterations suffice (and 9 do not guarantee it).
\end{exemplebox}

\textbf{Newton--Raphson: agreement of successive iterates.}\par
Newton--Raphson typically \emph{doubles} the number of correct digits at each step (quadratic convergence), but gives no built-in error bound. The practical test is: \cle{stop when two successive iterates agree to the required number of decimal places}, then \emph{confirm} with the sign-change test on the rounded interval. Agreement of $x_n$ and $x_{n+1}$ to $k$ d.p.\ is strong evidence, but the sign change is the proof.

\begin{popfigure}
\begin{center}
\begin{tikzpicture}
\begin{semilogyaxis}[
    width=0.86\linewidth, height=6.2cm,
    xlabel={Iteration number $n$},
    ylabel={Absolute error $|x_n-\alpha|$},
    xmin=0, xmax=6, ymin=1e-10, ymax=1,
    xtick={0,1,2,3,4,5,6},
    legend style={at={(0.97,0.97)},anchor=north east,font=\small},
    grid=both, grid style={popline!40},
]
\addplot[color=popblue, mark=*, thick] coordinates {
(0,0.5) (1,0.25) (2,0.125) (3,0.0625) (4,0.03125) (5,0.015625) (6,0.0078125)};
\addlegendentry{Bisection}
\addplot[color=poporange, mark=square*, thick] coordinates {
(0,0.2) (1,0.02) (2,0.0004) (3,0.00000016) (4,3e-14)};
\addlegendentry{Newton--Raphson}
\end{semilogyaxis}
\end{tikzpicture}
\end{center}
\legende{Typical convergence on the same equation (log scale): bisection halves the error each step, while Newton--Raphson roughly squares it --- the error plunges after two or three steps.}
\end{popfigure}

\courssub{Comparing the Methods and Choosing Wisely}

\begin{propriete}[Comparison of the root-finding methods]
\begin{center}
\begin{tabularx}{\linewidth}{|l|X|X|X|}
\hline
\rowcolor{popblueL}
 & \textbf{Bisection} & \textbf{Newton--Raphson} & \textbf{Simple iteration } $x=g(x)$ \\
\hline
\textbf{Speed} & Slow (linear; error halved each step) & Very fast (quadratic near the root) & Linear; fast only if $|g'(\alpha)|$ is small \\
\hline
\textbf{Reliability} & Always converges if a sign change is bracketed & May diverge or cycle for a poor $x_0$; fails if $f'(x_n)=0$ & Converges only if $|g'(x)|<1$ near the root \\
\hline
\textbf{Requirements} & Continuity and a sign change on $[a,b]$ & Formula for $f'(x)$; good starting value & A rearrangement $x=g(x)$ with $|g'|<1$ \\
\hline
\textbf{Error control} & Exact bound $\frac{b-a}{2^{n+1}}$ & Agreement of iterates $+$ sign-change check & $|x_{n+1}-x_n|$ as a guide \\
\hline
\end{tabularx}
\end{center}
\emph{(Stated for reference; the reasoning behind each row is examinable as commentary, not as a formal proof.)}
\end{propriete}

\begin{methode}[Choosing the right method in an examination]
\begin{enumerate}
\item \textbf{Locate} the root first with a sign change --- this is compulsory whatever method follows.
\item If the question says \emph{``show that the root is \ldots correct to \ldots''}, use the \textbf{sign-change test} on the rounded interval.
\item If $f'(x)$ is easy to compute and a good $x_0$ is given, use \textbf{Newton--Raphson}: two or three iterations usually suffice.
\item If $f'$ is awkward or unavailable, or the question imposes it, use \textbf{bisection} (or the given iterative formula).
\item If the question gives a rearrangement $x_{n+1}=g(x_n)$, use it as instructed and check $|g'|<1$ if asked about convergence.
\end{enumerate}
\end{methode}

\courssub{Application 1: Finance --- the Interest Rate on a Loan}

A trader in Douala borrows $P = 1\,000\,000$ FCFA and repays it by $12$ monthly instalments of $M = 90\,000$ FCFA, the first one month after the loan. If $i$ is the monthly interest rate, equating present values gives
\[ P = M\,\frac{1-(1+i)^{-12}}{i}. \]
With $x = 1+i$ this reduces to a polynomial equation; numerically it is cleaner to solve
\[ f(i) = 90\,000\bigl(1-(1+i)^{-12}\bigr) - 1\,000\,000\, i = 0. \]
Since $f(0.01) \approx 90\,000(0.11255) - 10\,000 = 129.6 > 0$ and $f(0.02) \approx 90\,000(0.21151) - 20\,000 = -964.4 < 0$, the rate lies between $1\%$ and $2\%$ per month. Newton--Raphson with
\[ f'(i) = 1\,080\,000(1+i)^{-13} - 1\,000\,000, \qquad i_{n+1} = i_n - \frac{f(i_n)}{f'(i_n)}, \]
starting from $i_0 = 0.012$:
\begin{align*}
i_1 &= 0.012 - \frac{-30.0}{-75\,520} \approx 0.011603, \\
i_2 &= 0.011603 - \frac{7.0}{-70\,120} \approx 0.011703, \\
i_3 &= 0.011703 - \frac{-3.0}{-69\,800} \approx 0.011660, \\
i_4 &\approx 0.011660.
\end{align*}
The iterates agree to 5 d.p.; a sign-change test on $[0.011655, 0.011665]$ confirms $i \approx 0.01166$ correct to 5 d.p., i.e.\ about $1.166\%$ per month (approximately $14.9\%$ per year). No algebraic formula exists for $i$ --- numerical methods are the \emph{only} practical tool, which is exactly how banks and mobile-money platforms compute such rates.

\courssub{Application 2: Engineering --- Depth of a Floating Sphere}

\begin{popfigure}
\begin{center}
\begin{tikzpicture}[scale=1.05]
% water
\fill[popblueL] (-3.4,-1.9) rectangle (3.4,0);
\draw[popblue, thick] (-3.4,0) -- (3.4,0);
\node[popblue, font=\small] at (2.6,-0.25) {water surface};
% sphere
\draw[popdark, thick, fill=white] (0,0.6) circle (1.6);
% submerged part shading
\begin{scope}
\clip (0,0.6) circle (1.6);
\fill[popblue!35] (-1.7,-1.9) rectangle (1.7,0);
\end{scope}
\draw[popdark, thick] (0,0.6) circle (1.6);
% radius arrow
\draw[->, poporange, thick] (0,0.6) -- (1.13,1.73) node[midway, above left, font=\small]{$r$};
% depth arrow
\draw[<->, poppurple, thick] (-2.1,-1.0) -- (-2.1,0) node[midway, left, font=\small]{$h$};
\fill[popdark] (0,0.6) circle (0.05);
\node[font=\small, popdark] at (0.15,0.42) {$O$};
\node[font=\small, align=center] at (0,-2.35) {submerged depth $h$ (unknown root)};
\end{tikzpicture}
\end{center}
\legende{A sphere of radius $r$ floating in water: the submerged depth $h$ is fixed by Archimedes' principle and is found by solving $f(h)=0$.}
\end{popfigure}

A wooden sphere of radius $r$ and relative density $\rho$ (with $0<\rho<1$) floats in water. The volume of the submerged spherical cap of depth $h$ is $V = \dfrac{\pi h^2(3r-h)}{3}$, and Archimedes' principle (weight of displaced water $=$ weight of sphere) gives
\[ \frac{\pi h^2(3r-h)}{3} = \rho\cdot\frac{4\pi r^3}{3}
\quad\Longleftrightarrow\quad
f(h) = h^3 - 3rh^2 + 4\rho r^3 = 0. \]
This cubic has no convenient exact solution, so we solve it numerically. Take $r = 10$ cm and wood of relative density $\rho = 0.6$: $f(h) = h^3 - 30h^2 + 2400$. Since $f(11) = 1331 - 3630 + 2400 = 101 > 0$ and $f(12) = 1728 - 4320 + 2400 = -192 < 0$, the root lies in $(11,12)$. Newton--Raphson with $f'(h) = 3h^2 - 60h$ from $h_0 = 11$:
\begin{align*}
h_1 &= 11 - \frac{101}{-297} \approx 11.34007, \\
h_2 &= 11.34007 - \frac{0.59}{-294.61} \approx 11.34207, \\
h_3 &= 11.34207 - \frac{-0.11}{-294.60} \approx 11.34170, \\
h_4 &\approx 11.3418.
\end{align*}
The iterates agree to 4 d.p.; checking $f(11.34125) > 0$ and $f(11.34175) < 0$ (sign change on the rounded interval for 3 d.p.) confirms $h = 11.342$ cm correct to 3 d.p. The same equation --- with different constants --- governs pipeline floats, buoys in the port of Kribi, and the calibration of fuel-tank gauges.

\courssub{Application 3: Statistics Link --- Calibrating a Probability Model}

\begin{savaistu}[Numerical methods inside statistics]
Many statistical questions secretly end with ``solve $f(\theta)=0$''. For a Binomial model $X \sim \mathrm{B}(n,p)$, finding the value of $p$ that makes $P(X \geq 1)$ equal to a target, or finding the sample size $n$ needed for a given confidence, leads to equations with no closed-form solution. Examiners love this crossover: it tests both your statistics and your numerical methods in one question.
\end{savaistu}

\begin{exoresolu}[Finding a probability parameter numerically]
A machine fills sachets of water; the probability that a sachet is underweight is $p$. Quality control requires that in a batch of $20$ sachets, $P(\text{at least one underweight}) = 0.5$. Find $p$ correct to 4 decimal places.
\tcblower
We need $1-(1-p)^{20} = 0.5$, i.e.\ $(1-p)^{20} = 0.5$. (Here an exact solution $p = 1 - 0.5^{1/20}$ exists, but let us treat it as the examiner intends when the equation is \emph{not} solvable exactly --- the method is identical.) Set $f(p) = (1-p)^{20} - 0.5$. Then
\[ f(0.03) = 0.97^{20} - 0.5 \approx 0.5438 - 0.5 = 0.0438 > 0, \]
\[ f(0.04) = 0.96^{20} - 0.5 \approx 0.4420 - 0.5 = -0.0580 < 0, \]
so $p \in (0.03, 0.04)$. Newton--Raphson with $f'(p) = -20(1-p)^{19}$, starting at $p_0 = 0.035$:
\begin{align*}
p_1 &= 0.035 - \frac{0.965^{20}-0.5}{-20(0.965)^{19}} = 0.035 - \frac{-0.009877}{-10.158} \approx 0.034028,\\
p_2 &= 0.034028 - \frac{0.965972^{20}-0.5}{-20(0.965972)^{19}} \approx 0.034028 - \frac{0.0003}{-10.36} \approx 0.034057,\\
p_3 &\approx 0.034064.
\end{align*}
The iterates settle at $p \approx 0.03406$. Sign-change test on the rounded interval $[0.034055, 0.034065]$: $f(0.034055) > 0$ and $f(0.034065) < 0$, so $p = 0.0341$ correct to 4 d.p. \textbf{Interpretation:} the factory must keep the underweight rate at about $3.41\%$ for the batch risk to be exactly one-half.
\end{exoresolu}

\courssub{Full GCE-Style Synthesis}

\begin{methode}[Exam strategy: the four-move answer]
Every complete numerical-methods answer should show four moves, in this order:
\begin{enumerate}
\item \textbf{Locate:} exhibit a sign change of $f$ on an interval $[a,b]$ (continuity $+$ IVT).
\item \textbf{Iterate:} apply the required method, showing the formula, the substitution, and at least the first two iterates with full calculator values before rounding.
\item \textbf{Verify:} justify the final accuracy --- sign change on the rounded interval $[x_0 \pm \frac12\times10^{-k}]$, or the bisection error bound.
\item \textbf{Interpret:} answer the question asked, in context, with units if any (``the depth is $11.342$ cm'', not just ``$x = 11.342$'').
\end{enumerate}
\end{methode}

\begin{exoresolu}[Complete synthesis question]
A spherical tank of radius $2$ m lies horizontally. The volume of liquid of depth $x$ metres in the tank is $V(x) = \dfrac{\pi x^2(6-x)}{3}$. The tank contains $9\ \mathrm{m^3}$ of liquid.
\begin{enumerate}
\item[(a)] Show that $x$ satisfies $x^3 - 6x^2 + \dfrac{27}{\pi} = 0$, and that $1.3 < x < 1.4$.
\item[(b)] Use the Newton--Raphson method once, starting from $x_0 = 1.3$, to find a second approximation.
\item[(c)] Refine further and give $x$ correct to 3 decimal places, justifying your answer.
\end{enumerate}
\tcblower
\textbf{(a)} Setting $V(x) = 9$: $\dfrac{\pi x^2(6-x)}{3} = 9 \Rightarrow x^2(6-x) = \dfrac{27}{\pi} \Rightarrow f(x) = x^3 - 6x^2 + \dfrac{27}{\pi} = 0$, where $\frac{27}{\pi} \approx 8.5944$.
\[ f(1.3) = 2.197 - 10.14 + 8.5944 = 0.6514 > 0, \]
\[ f(1.4) = 2.744 - 11.76 + 8.5944 = -0.4216 < 0, \]
so the root lies in $(1.3, 1.4)$.

\textbf{(b)} $f'(x) = 3x^2 - 12x$, so $f'(1.3) = 5.07 - 15.6 = -10.53$ and
\[ x_1 = 1.3 - \frac{0.6514}{-10.53} \approx 1.36186. \]

\textbf{(c)} Continuing:
\begin{align*}
f(1.36186) &\approx 2.5265 - 11.1275 + 8.5944 = -0.0066, \\
f'(1.36186) &\approx 5.564 - 16.342 = -10.778, \\
x_2 &= 1.36186 - \frac{-0.0066}{-10.778} \approx 1.36125, \\
x_3 &\approx 1.36125.
\end{align*}
The iterates agree to 5 d.p. \textbf{Verification:} on the rounded interval $[1.36075, 1.36125]$,
\[ f(1.36075) \approx +0.0019 > 0, \qquad f(1.36125) \approx -0.0015 < 0, \]
a sign change, so $x = 1.361$ correct to 3 d.p. \textbf{Interpretation:} the liquid stands $1.361$ m deep in the tank --- the number a technician would mark on the dip-stick.
\end{exoresolu}

\begin{attention}[Final warnings for the examination]
\begin{itemize}
\item Never claim ``correct to 3 d.p.'' from the calculator display alone: justify with a sign change on $[x_0 - 0.0005,\, x_0 + 0.0005]$.
\item Keep at least two guard digits during iteration; round \emph{only} at the final answer.
\item In Newton--Raphson, write the formula $x_{n+1} = x_n - f(x_n)/f'(x_n)$ explicitly --- method marks depend on it.
\item Check that your final answer lies in the physically meaningful range (a depth cannot exceed $2r$; a probability lies in $[0,1]$; an interest rate is positive).
\item If bisection is demanded, do not switch to Newton--Raphson halfway: follow the instructed method.
\end{itemize}
\end{attention}

\begin{exercice}[Practice]
\begin{enumerate}
\item A root of $f(x) = x^3 + x - 3$ is bracketed in $[1, 2]$. How many bisection iterations guarantee an error below $0.00005$?
\item Show that the equation $e^{-x} = x$ has a root in $[0.5, 0.6]$. Use Newton--Raphson from $x_0 = 0.55$ to find the root correct to 4 decimal places, and justify the accuracy with a sign-change test.
\item A loan of $500\,000$ FCFA is repaid by $10$ monthly payments of $53\,000$ FCFA. Write down the equation for the monthly rate $i$, show it has a root between $0.01$ and $0.02$, and find $i$ correct to 4 decimal places.
\item For the floating sphere with $r = 5$ cm and $\rho = 0.5$, show that $h$ satisfies $h^3 - 15h^2 + 250 = 0$ and find $h$ correct to 3 decimal places. Why must the root with $h > 10$ be rejected?
\end{enumerate}
\end{exercice}

\begin{corrige}[Exercise 1]
We need $\dfrac{b-a}{2^{n+1}} \leq 0.00005$ with $b-a = 1$: $2^{n+1} \geq 20\,000$. Since $2^{14} = 16\,384 < 20\,000$ and $2^{15} = 32\,768 \geq 20\,000$, we need $n+1 = 15$, i.e.\ $n = 14$ iterations.
\end{corrige}

\begin{corrige}[Exercise 2]
$f(x) = e^{-x} - x$: $f(0.5) = 0.6065 - 0.5 = 0.1065 > 0$; $f(0.6) = 0.5488 - 0.6 = -0.0512 < 0$: sign change, root in $(0.5, 0.6)$. With $f'(x) = -e^{-x} - 1$ and $x_0 = 0.55$: $x_1 = 0.55 - \dfrac{0.02695}{-1.57695} \approx 0.56709$; $x_2 \approx 0.56714$; $x_3 \approx 0.56714$. Test on $[0.567135, 0.567145]$: $f(0.567135) > 0$, $f(0.567145) < 0$, so $x = 0.5671$ correct to 4 d.p.
\end{corrige}

\begin{corrige}[Exercise 3]
Equation: $53\,000\bigl(1-(1+i)^{-10}\bigr) = 500\,000\, i$, i.e.\ $f(i) = 53\,000\bigl(1-(1+i)^{-10}\bigr) - 500\,000\,i$.
$f(0.01) \approx 53\,000(0.09471) - 5\,000 = 19.8 > 0$; $f(0.02) \approx 53\,000(0.17965) - 10\,000 = -478.4 < 0$: root in $(0.01, 0.02)$.
Newton--Raphson with $f'(i) = 530\,000(1+i)^{-11} - 500\,000$, starting at $i_0 = 0.011$:
\begin{align*}
i_1 &= 0.011 - \frac{12.0}{-29\,890} \approx 0.011401, \\
i_2 &= 0.011401 - \frac{-82.5}{-32\,010} \approx 0.011143, \\
i_3 &= 0.011143 - \frac{-2.1}{-30\,800} \approx 0.011136, \\
i_4 &\approx 0.011136.
\end{align*}
Sign change on $[0.0111315, 0.0111365]$ confirms $i = 0.0111$ correct to 4 d.p., i.e.\ about $1.11\%$ per month.
\end{corrige}

\begin{corrige}[Exercise 4]
Archimedes: $\frac{\pi h^2(15-h)}{3} = 0.5\cdot\frac{4\pi(125)}{3} \Rightarrow h^2(15-h) = 250 \Rightarrow f(h) = h^3 - 15h^2 + 250 = 0$.
$f(4) = 64 - 240 + 250 = 74 > 0$; $f(5) = 125 - 375 + 250 = 0$ exactly --- so $h = 5.000$ cm (the sphere floats exactly half-submerged, as expected for $\rho = 0.5$). The cubic also has a root near $h \approx 13.7$, but $h > 2r = 10$ is physically impossible (the sphere would be more than fully submerged), so it is rejected: always match the mathematical roots to the physical context.
\end{corrige}

\begin{aretenir}[Key points of the section]
\begin{itemize}
\item ``Correct to $k$ d.p.'' means error $\leq \frac12\times 10^{-k}$; \textbf{prove it} by a sign change on $[x_0 - \frac12 10^{-k},\, x_0 + \frac12 10^{-k}]$.
\item Bisection: guaranteed convergence, error bound $\frac{b-a}{2^{n+1}}$, slow. Newton--Raphson: quadratic speed, needs $f'$ and a good start, no built-in bound --- verify afterwards.
\item Exam structure: \cle{locate} (sign change) $\rightarrow$ \cle{iterate} (show formula and values) $\rightarrow$ \cle{verify} (sign-change test) $\rightarrow$ \cle{interpret} (context, units).
\item Real problems --- loan rates, floating bodies, probability calibration --- reduce to $f(x)=0$ and are solved exactly as in the exercises; reject non-physical roots.
\end{itemize}
\end{aretenir}

\newpage

\coursec{Integration activity}

\coursec{Integration Activity: The Mobile-Money Interest Rate Mystery}

\courssub{Aims of the Activity}
This integration activity brings together all the skills developed in the chapter in a single, authentic financial problem. By the end of the activity, you should be able to:
\begin{itemize}
\item translate a real-life situation (a loan with fixed monthly repayments) into an equation of the form $f(r)=0$ that \emph{cannot} be solved algebraically;
\item locate a root of $f(r)=0$ by detecting a \cle{change of sign} of $f$ over an interval;
\item apply the \cle{bisection method} step by step, keeping track of the interval containing the root and of the maximum possible error;
\item apply the \cle{Newton--Raphson method} with the iteration $r_{n+1}=r_n-\dfrac{f(r_n)}{f'(r_n)}$ and control the number of decimal places guaranteed;
\item verify the accuracy of a numerical answer by a final sign-change check;
\item compare the speed of convergence of the two methods and communicate a conclusion in clear, non-technical language, as a consultant would.
\end{itemize}
Work in groups of three or four. Each group needs a scientific calculator. Keep at least six decimal places in all intermediate calculations, and only round at the very end.

\courssub{Situation}
\textbf{The Mobile-Money Interest Rate Mystery.}\par
Madame Eyenga runs a provisions store near the Mokolo market in Yaoundé. To restock her shop before the end-of-year festivities, she borrows $500\,000$~FCFA from a mobile-money lending service. The contract states that she will repay $115\,000$~FCFA at the \emph{end of each month}, for exactly $5$ months, after which the debt will be completely cleared.

The contract does not state the monthly interest rate $r$ explicitly. However, the lender explains that each repayment of $115\,000$~FCFA first pays the interest due on the outstanding balance, and the remainder reduces the capital. A financial analyst can show (by discounting each of the five payments back to the date of the loan) that the monthly rate $r$ satisfies the equation
\[ 500\,000\,r = 115\,000\left(1-(1+r)^{-5}\right). \]
Dividing both sides by $5\,000$, this becomes
\[ 100r = 23\left(1-(1+r)^{-5}\right). \]
This equation involves $r$ both linearly and inside a power with a negative exponent: there is \textbf{no algebraic method} to solve it exactly. Madame Eyenga wants to know the monthly rate she is really being charged, correct to $4$ decimal places (that is, to the nearest $0.01\%$), so that she can compare this offer with a bank loan.

Your group is hired as her advisory team. You must find the rate, justify its accuracy, and write her a short advisory note.

\courssub{Tasks}
\begin{enumerate}
\item \textbf{Setting up.} Define $f(r)=100r-23\left(1-(1+r)^{-5}\right)$. Explain why solving the problem is equivalent to finding the root of $f(r)=0$ in the interval $(0,1)$, and state why an exact algebraic solution is impossible.
\item \textbf{Locating the root.} Compute $f(0.04)$ and $f(0.05)$, each to at least $4$ decimal places. Explain, quoting the Intermediate Value Theorem, why there is a root $r^{*}$ in the interval $(0.04,0.05)$.
\item \textbf{Bisection.} Starting from the interval $[0.04,0.05]$, carry out \textbf{three} steps of the bisection method. Present your work in a table with columns: step number, interval $[a_n,b_n]$, midpoint $m_n$, sign of $f(m_n)$, new interval. After three steps, state the interval containing $r^{*}$, its width, and the maximum error if you take the midpoint as your estimate.
\item \textbf{Newton--Raphson.} Show that $f'(r)=100-115(1+r)^{-6}$. Starting from $r_0=0.05$, apply the Newton--Raphson iteration
\[ r_{n+1}=r_n-\frac{f(r_n)}{f'(r_n)} \]
until two successive iterates agree to $6$ decimal places. Record $r_0,r_1,r_2,\dots$ in a table with $f(r_n)$ and $f'(r_n)$.
\item \textbf{Accuracy check.} Let $\tilde{r}$ be your answer rounded to $4$ decimal places. By evaluating $f$ at $\tilde{r}-0.00005$ and at $\tilde{r}+0.00005$, confirm by a sign change that the root, rounded to $4$ decimal places, is indeed $\tilde{r}$.
\item \textbf{Advisory note.} Write a short note (five to eight lines) to Madame Eyenga, in plain language, stating: the monthly interest rate as a percentage; how confident you are in this figure; and a comparison of the two methods (bisection versus Newton--Raphson) in terms of speed and reliability.
\end{enumerate}

\begin{attention}[Traps to avoid]
Do not confuse the \emph{monthly} rate with an \emph{annual} rate. In the bisection table, the new interval is chosen according to the \emph{sign} of $f$ at the midpoint, not its size. In Newton--Raphson, never round the iterates before the final answer: premature rounding can destroy the rapid convergence.
\end{attention}

\courssub{Worked Solutions}

\begin{definition}[Intermediate Value Theorem]
If a function $f$ is continuous on a closed interval $[a,b]$ and $f(a)$ and $f(b)$ have opposite signs, then there exists at least one $c\in(a,b)$ such that $f(c)=0$.
\end{definition}

\begin{exoresolu}[Task 1: Setting up the equation]
The loan equation $500\,000\,r = 115\,000\left(1-(1+r)^{-5}\right)$ simplifies, on division by $5\,000$, to
\[ 100r = 23\left(1-(1+r)^{-5}\right). \]
Define
\[ f(r) = 100r - 23\left(1-(1+r)^{-5}\right). \]
The monthly rate $r^{*}$ is precisely the value for which $f(r^{*})=0$. The unknown $r$ appears both as a linear term $100r$ and inside the non-linear term $(1+r)^{-5}$. No algebraic manipulation (factorisation, substitution, or use of standard functions) can isolate $r$; the equation is \emph{transcendental}. Hence a numerical method is mandatory. As a monthly interest rate, $r^{*}$ must lie between $0$ and $1$; since the total repayment is $5\times115\,000=575\,000$~FCFA (only $15\%$ more than the principal), the rate is small, of the order of a few percent.
\end{exoresolu}

\begin{exoresolu}[Task 2: Locating the root by a sign change]
The function $f$ is a sum of continuous functions (polynomial and power functions) on $(0,1)$, therefore it is continuous on $[0.04,0.05]$. We evaluate $f$ at the endpoints, keeping six decimal places for safety:
\begin{align*}
f(0.04) &= 4 - 23\left(1 - 1.04^{-5}\right) \\
        &= 4 - 23\left(1 - 0.821927\right) \\
        &= 4 - 23\times0.178073 \\
        &= 4 - 4.095679 = -0.095679 \approx -0.0957,\\[4pt]
f(0.05) &= 5 - 23\left(1 - 1.05^{-5}\right) \\
        &= 5 - 23\left(1 - 0.783526\right) \\
        &= 5 - 23\times0.216474 \\
        &= 5 - 4.978902 = 0.021098 \approx 0.0211.
\end{align*}
Since $f$ is continuous on $[0.04,0.05]$ and $f(0.04)<0<f(0.05)$, the \textbf{Intermediate Value Theorem} guarantees the existence of a root $r^{*}\in(0.04,0.05)$.

\begin{popfigure}
\begin{tikzpicture}
\begin{axis}[
    width=0.9\linewidth, height=6cm,
    axis lines=middle,
    xlabel=$r$, ylabel=$f(r)$,
    xmin=0.035, xmax=0.055,
    ymin=-0.12, ymax=0.04,
    xtick={0.04,0.045,0.05},
    ytick={-0.1,-0.05,0,0.02},
    grid=major, grid style={popline},
    tick label style={font=\small},
    label style={font=\small},
]
\addplot[domain=0.035:0.055, samples=200, popblue, thick] {100*x - 23*(1 - (1+x)^(-5))};
\addplot[only marks, mark=*, poporange] coordinates {(0.04,-0.095679) (0.05,0.021098)};
\addplot[dashed, popmuted] coordinates {(0.04,-0.12) (0.04,0.04)};
\addplot[dashed, popmuted] coordinates {(0.05,-0.12) (0.05,0.04)};
\node[poporange, anchor=north west, font=\small] at (axis cs:0.04,-0.095679) {$(0.04,\,f<0)$};
\node[poporange, anchor=south west, font=\small] at (axis cs:0.05,0.021098) {$(0.05,\,f>0)$};
\end{axis}
\end{tikzpicture}
\legende{Graph of $f(r)$ showing the sign change on $[0.04,0.05]$.}
\end{popfigure}
\end{exoresolu}

\begin{methode}[Bisection method -- step by step]
\begin{enumerate}
\item Start with an interval $[a,b]$ where $f(a)f(b)<0$.
\item Compute the midpoint $m=\dfrac{a+b}{2}$ and evaluate $f(m)$.
\item If $f(m)=0$, $m$ is the root (stop). Otherwise, if $f(a)f(m)<0$, the root lies in $[a,m]$; set $b=m$. If $f(m)f(b)<0$, the root lies in $[m,b]$; set $a=m$.
\item Repeat until the interval width $b-a$ is smaller than the required tolerance. The best estimate is the final midpoint; the maximum error is half the final width.
\end{enumerate}
\end{methode}

\begin{exoresolu}[Task 3: Three steps of the bisection method]
We apply the method starting from $[a_0,b_0]=[0.04,0.05]$.

\begin{center}
\begin{tabular}{|c|l|l|l|l|}
\hline
\rowcolor{popblueL} Step $n$ & Interval $[a_n,b_n]$ & Midpoint $m_n$ & $f(m_n)$ & New interval \\
\hline
$1$ & $[0.040000,\ 0.050000]$ & $0.045000$ & $-0.043627$ & $[0.045000,\ 0.050000]$ \\
\hline
$2$ & $[0.045000,\ 0.050000]$ & $0.047500$ & $-0.012789$ & $[0.047500,\ 0.050000]$ \\
\hline
$3$ & $[0.047500,\ 0.050000]$ & $0.048750$ & $+0.003912$ & $[0.047500,\ 0.048750]$ \\
\hline
\end{tabular}
\end{center}

\textbf{Sample calculation (Step 1):}
$f(0.045)=4.5-23\left(1-1.045^{-5}\right)=4.5-23(1-0.802451)=4.5-4.543627=-0.043627<0$.
Since $f(0.04)<0$ and $f(0.045)<0$, the sign change occurs on $[0.045,0.05]$.

After three steps, $r^{*}\in[0.0475,\,0.04875]$. The interval width is $0.00125$. Taking the midpoint $\hat{r}=0.048125$ as the estimate, the maximum error is half the width:
\[ |r^{*}-0.048125|\le\frac{0.00125}{2}=0.000625. \]
This is \emph{not} yet accurate enough for $4$ decimal places (which requires an error $<0.00005$). The width must fall below $0.0001$; starting from $0.01$, we need $0.01/2^n<0.0001\Rightarrow 2^n>100\Rightarrow n\ge7$ steps in total, i.e. \textbf{four more steps} after the three shown.
\end{exoresolu}

\begin{methode}[Newton--Raphson method -- step by step]
\begin{enumerate}
\item Choose a starting value $r_0$ close to the root, where $f'(r_0)\neq0$.
\item Iterate using $r_{n+1}=r_n-\dfrac{f(r_n)}{f'(r_n)}$.
\item Stop when $|r_{n+1}-r_n|$ is smaller than the desired tolerance (e.g. $<5\times10^{-7}$ for $6$ decimal places).
\item The method converges quadratically (number of correct digits roughly doubles each step) provided $f'(r^{*})\neq0$ and $r_0$ is sufficiently close.
\end{enumerate}
\end{methode}

\begin{exoresolu}[Task 4: Newton--Raphson iteration]
First, differentiate $f$:
\[ f(r)=100r-23+23(1+r)^{-5} \quad\Rightarrow\quad f'(r)=100-115(1+r)^{-6}. \]
On $[0.04,0.05]$, $(1+r)^{-6}\in[0.746,0.790]$, so $f'(r)\in[100-115\times0.790,\,100-115\times0.746]\approx[9.1,\,14.2]>0$. The derivative is positive and safely away from zero, so the iteration is well behaved.

Take $r_0=0.05$ (the right endpoint where $f>0$).
\begin{align*}
f(0.05) &= 0.021098, \qquad f'(0.05) = 100-115(1.05)^{-6}=100-85.8146=14.1854,\\
r_1 &= 0.05 - \frac{0.021098}{14.1854} = 0.0485124 \approx 0.048512.\\[4pt]
f(0.048512) &= 4.8512 - 23\left(1-1.048512^{-5}\right) = 4.8512 - 4.851097 = 0.000103,\\
f'(0.048512) &= 100 - 115(1.048512)^{-6} = 100 - 86.5473 = 13.4527,\\
r_2 &= 0.048512 - \frac{0.000103}{13.4527} = 0.0485047 \approx 0.048505.\\[4pt]
f(0.048505) &\approx 0.000000 \text{ (to 6 d.p.)}, \quad f'(0.048505) \approx 13.4507.
\end{align*}

\begin{center}
\begin{tabular}{|c|l|l|l|}
\hline
\rowcolor{popblueL} $n$ & $r_n$ & $f(r_n)$ & $f'(r_n)$ \\
\hline
$0$ & $0.050000$ & $0.021098$ & $14.1854$ \\
\hline
$1$ & $0.048512$ & $0.000103$ & $13.4527$ \\
\hline
$2$ & $0.048505$ & $0.000000$ & $13.4507$ \\
\hline
\end{tabular}
\end{center}

The iterates $r_1=0.048512$ and $r_2=0.048505$ agree to $4$ decimal places ($0.0485$), and $f(r_2)\approx0$. The iteration has converged. Only \textbf{two} Newton steps achieve what bisection could not in three.
\end{exoresolu}

\begin{exoresolu}[Task 5: Accuracy check by sign change]
Our candidate rounded to $4$ decimal places is $\tilde{r}=0.0485$. To be certain that the true root rounds to this value, we verify a sign change on the interval of all numbers that round to $0.0485$, namely $[0.04845,\,0.04855)$ (the upper endpoint $0.04855$ rounds to $0.0486$, but the root lies strictly inside). We evaluate $f$ at the endpoints:
\begin{align*}
f(0.04845) &= 4.845 - 23\left(1-1.04845^{-5}\right) = 4.845 - 4.845732 = -0.000732 < 0,\\
f(0.04855) &= 4.855 - 23\left(1-1.04855^{-5}\right) = 4.855 - 4.854245 = +0.000755 > 0.
\end{align*}
Since $f$ is continuous and changes sign on $[0.04845,0.04855]$, there is a root in this interval. Every number strictly between $0.04845$ and $0.04855$ rounds to $0.0485$ to four decimal places. Hence
\[ \boxed{r^{*}=0.0485 \quad\text{correct to }4\text{ decimal places.}} \]
\end{exoresolu}

\courssub{Advisory Note}

\begin{aretenir}[Note to Madame Eyenga]
Dear Madame Eyenga, we have analysed your loan of $500\,000$~FCFA repaid by five monthly instalments of $115\,000$~FCFA. The monthly interest rate written into your contract is approximately $r=0.0485$, that is $\mathbf{4.85\%}$ \textbf{per month}, and we have verified this figure to be correct to the nearest $0.01\%$. Be aware that $4.85\%$ per month is a very high rate: over five months you repay $575\,000$~FCFA, i.e. $75\,000$~FCFA of interest, and the equivalent annual cost is far above $12\times4.85\%$ because interest compounds. We recommend that you compare this with a bank loan before borrowing again. To obtain this result we used two numerical methods: the bisection method, which is slow but always reliable (three steps only narrowed the rate to within $\pm0.0006$), and the Newton--Raphson method, which reached the answer in just two steps once a good starting value was known. Used together, the two methods give both safety and speed.
\end{aretenir}

\begin{savaistu}[Why the two methods complement each other]
Professional software (spreadsheets, banking calculators, the \texttt{solver} of your calculator) almost always combines the two strategies you have just used: a \emph{bracketing} method such as bisection first traps the root inside an interval where the function is well behaved, then a fast method such as Newton--Raphson polishes the answer to full machine precision. Bisection alone is guaranteed but slow (one binary digit per step); Newton--Raphson alone is fast (the number of correct digits roughly doubles per step) but can fail if started far from the root or where $f'$ is close to $0$. Together, they are unbeatable.
\end{savaistu}

\newpage

\coursec{Methods and worked examples}

\coursec{Introduction and Location of Roots}

\begin{savaistu}[Why numerical methods?]
Many equations arising in science, engineering and economics cannot be solved exactly by algebraic manipulation. For instance, the equation $x = \cos x$ or $e^{-x} = x$ have no closed-form solution in terms of elementary functions. Numerical methods provide systematic algorithms to approximate roots to any desired accuracy. In the Cameroonian context, these techniques are used in civil engineering (settlement of foundations), in physics (finding equilibrium positions), and in financial mathematics (computing internal rates of return).
\end{savaistu}

\begin{definition}[Root of an equation]
A \cle{root} (or \cle{zero}) of the equation $f(x)=0$ is a value $\alpha$ such that $f(\alpha)=0$. Geometrically, it is the $x$-coordinate of a point where the curve $y=f(x)$ crosses the $x$-axis.
\end{definition}

\begin{propriete}[Intermediate Value Theorem (IVT)]
If $f$ is \cle{continuous} on the closed interval $[a,b]$ and $f(a)$ and $f(b)$ have \cle{opposite signs} (i.e. $f(a)f(b)<0$), then there exists at least one $\alpha \in (a,b)$ such that $f(\alpha)=0$.
\end{propriete}

\begin{attention}[Continuity is essential]
The IVT requires continuity on the \emph{whole} interval $[a,b]$. For example, $f(x)=\frac{1}{x}$ satisfies $f(-1)=-1$ and $f(1)=1$ (opposite signs), but there is \textbf{no root} in $[-1,1]$ because $f$ is not continuous at $x=0$. Always verify continuity before applying the sign-change test.
\end{attention}

\begin{methode}[Locating a root in an interval]
To show that $f(x)=0$ has a root in $(a,b)$ and to narrow its location:
\begin{enumerate}
\item \textbf{Check continuity:} Confirm $f$ is continuous on $[a,b]$. Polynomials, exponentials, $\sin x$, $\cos x$ are continuous everywhere; rational functions are continuous except where the denominator vanishes.
\item \textbf{Evaluate endpoints:} Compute $f(a)$ and $f(b)$ to at least 3 decimal places.
\item \textbf{Sign change:} If $f(a)f(b)<0$, the IVT guarantees at least one root $\alpha \in (a,b)$.
\item \textbf{Narrow the interval:} Test $f$ at midpoints or other convenient values (e.g. $a+0.5$, $a+0.25$) until the interval width meets the requirement.
\item \textbf{Uniqueness (if needed):} Show $f$ is strictly monotonic on $[a,b]$ (e.g. $f'(x)>0$ or $f'(x)<0$ throughout). Then the root is unique.
\end{enumerate}
\textbf{Reflex:} Sign change $\Rightarrow$ \emph{existence}; monotonicity $\Rightarrow$ \emph{uniqueness}. Never claim a root exists just because values are ``close to zero''.
\end{methode}

\begin{exemplebox}[Sign-change table for $f(x)=x^3+2x-7$]
\begin{center}
\begin{tabular}{|c|cccc|}\hline
$x$ & $1$ & $1.5$ & $1.75$ & $2$ \\ \hline
$f(x)$ & $-4$ & $-0.625$ & $1.859$ & $5$ \\ \hline
\end{tabular}
\end{center}
Sign changes between $1.5$ and $1.75$ $\Rightarrow$ root in $(1.5,\,1.75)$. Since $f'(x)=3x^2+2>0$ everywhere, the root is unique.
\end{exemplebox}

\begin{exoresolu}[Locating a root of a cubic]
Show that the equation $x^{3}+2x-7=0$ has exactly one real root $\alpha$, and find an interval of width $0.5$ containing $\alpha$.
\tcblower
\textbf{Solution.} Let $f(x)=x^{3}+2x-7$; $f$ is a polynomial, hence continuous on $\mathbb{R}$.

\emph{Step 1: sign change.} Try integer values:
\[ f(1)=1+2-7=-4<0,\qquad f(2)=8+4-7=5>0. \]
Since $f(1)f(2)<0$, by the Intermediate Value Theorem there is a root $\alpha\in(1,2)$.

\emph{Step 2: uniqueness.} $f'(x)=3x^{2}+2>0$ for all $x$, so $f$ is strictly increasing on $\mathbb{R}$: the root is \textbf{unique}.

\emph{Step 3: narrowing to width $0.5$.} Evaluate at the midpoint $x=1.5$:
\[ f(1.5)=3.375+3-7=-0.625<0. \]
Since $f(1.5)<0$ and $f(2)>0$, the root lies in $(1.5,\,2)$, an interval of width $0.5$. \qquad $\boxed{\alpha\in(1.5,\,2)}$
\end{exoresolu}

\begin{exoresolu}[Locating a root involving a trigonometric term]
Show that the equation $x=\cos x$ has a root between $0$ and $1$ (radians), and determine an interval of width $0.25$ containing it.
\tcblower
\textbf{Solution.} Rewrite as $f(x)=x-\cos x=0$. $f$ is continuous everywhere.

\emph{Step 1.} $f(0)=0-1=-1<0$ and $f(1)=1-\cos 1\approx 1-0.5403=0.4597>0$. Sign change: a root exists in $(0,1)$.

\emph{Step 2: uniqueness.} $f'(x)=1+\sin x\geq 0$ for all $x$, and $f'(x)=0$ only at isolated points ($x=-\frac{\pi}{2}+2k\pi$), so $f$ is strictly increasing: the root is unique.

\emph{Step 3: width $0.25$.} Test $x=0.5$ and $x=0.75$:
\begin{align*}
f(0.5)&=0.5-\cos 0.5\approx 0.5-0.8776=-0.3776<0,\\
f(0.75)&=0.75-\cos 0.75\approx 0.75-0.7317=0.0183>0.
\end{align*}
Sign change between $0.5$ and $0.75$: $\boxed{\alpha\in(0.5,\,0.75)}$.
\end{exoresolu}

\begin{exoresolu}[Locating a root of an exponential equation]
Show that $e^{-x}=x$ has a unique root in $(0,1)$ and locate it to within $0.1$.
\tcblower
\textbf{Solution.} Let $f(x)=e^{-x}-x$. $f$ is continuous on $\mathbb{R}$.
$f(0)=1-0=1>0$; $f(1)=e^{-1}-1\approx 0.3679-1=-0.6321<0$. Sign change $\Rightarrow$ root in $(0,1)$.
$f'(x)=-e^{-x}-1<0$ for all $x$, so $f$ is strictly decreasing $\Rightarrow$ unique root.
Test $x=0.5$: $f(0.5)=e^{-0.5}-0.5\approx 0.6065-0.5=0.1065>0$.
Test $x=0.6$: $f(0.6)=e^{-0.6}-0.6\approx 0.5488-0.6=-0.0512<0$.
Root lies in $(0.5,\,0.6)$, width $0.1$. $\boxed{\alpha\in(0.5,\,0.6)}$.
\end{exoresolu}

\coursec{The Bisection Method}

\begin{propriete}[Bisection method principle]
If $f$ is continuous on $[a,b]$ with $f(a)f(b)<0$, let $m=\frac{a+b}{2}$.
\begin{itemize}
\item If $f(a)f(m)<0$, the root lies in $(a,m)$.
\item If $f(m)f(b)<0$, the root lies in $(m,b)$.
\item If $f(m)=0$, $m$ is the exact root.
\end{itemize}
Repeating this process halves the interval width at each step. After $n$ iterations, the interval width is $\frac{b-a}{2^n}$.
\end{propriete}

\begin{methode}[Carrying out the bisection method]
Given $f$ continuous on $[a,b]$ with $f(a)f(b)<0$:
\begin{enumerate}
\item Compute the midpoint $m=\dfrac{a+b}{2}$ and evaluate $f(m)$.
\item Compare the sign of $f(m)$ with the signs of $f(a)$ and $f(b)$:
\begin{itemize}
\item if $f(a)f(m)<0$, the root lies in $(a,m)$: set $b:=m$;
\item if $f(m)f(b)<0$, the root lies in $(m,b)$: set $a:=m$;
\item if $f(m)=0$, stop: $m$ is the exact root.
\end{itemize}
\item Repeat. After $n$ steps the interval has width $\dfrac{b-a}{2^{n}}$.
\item Stop when the interval width is less than the required tolerance, or when both endpoints agree to the requested number of decimal places.
\item Give the root as the midpoint of the final interval, or as the common rounded value of the endpoints.
\end{enumerate}
\textbf{Reflex:} To guarantee accuracy $\varepsilon$ when quoting the midpoint, you need the smallest $n$ with $\dfrac{b-a}{2^{n+1}}\leq \varepsilon$ (error $\leq$ half-width).
\end{methode}

\begin{exemplebox}[{Bisection iteration table for $x^3+2x-7=0$ on $[1.5,2]$}]
\begin{center}
\begin{tabular}{|c|c|c|c|c|}\hline
$n$ & $a$ & $b$ & $m$ & $f(m)$ \\ \hline
0 & 1.5 & 2.0 & -- & -- \\
1 & 1.5 & 1.75 & 1.75 & $+1.859$ \\
2 & 1.5 & 1.625 & 1.625 & $+0.541$ \\
3 & 1.5625 & 1.625 & 1.5625 & $-0.060$ \\
4 & 1.5625 & 1.59375 & 1.59375 & $+0.234$ \\ \hline
\end{tabular}
\end{center}
\end{exemplebox}

\begin{exoresolu}[Bisection on a cubic, step by step]
The equation $x^{3}+2x-7=0$ has a root $\alpha$ in $(1.5,\,2)$. Use the bisection method to find $\alpha$ correct to $1$ decimal place.
\tcblower
\textbf{Solution.} $f(x)=x^{3}+2x-7$, with $f(1.5)=-0.625$ and $f(2)=5$.

\emph{Iteration 1:} $m=1.75$; $f(1.75)=5.359375+3.5-7=1.859375>0$. Root in $(1.5,\,1.75)$.

\emph{Iteration 2:} $m=1.625$; $f(1.625)=4.291015625+3.25-7=0.541015625>0$. Root in $(1.5,\,1.625)$.

\emph{Iteration 3:} $m=1.5625$; $f(1.5625)=3.814697265625+3.125-7=-0.060302734375<0$. Root in $(1.5625,\,1.625)$.

\emph{Iteration 4:} $m=1.59375$; $f(1.59375)\approx 4.046+3.1875-7=0.2335>0$. Root in $(1.5625,\,1.59375)$.

Both endpoints round to $1.6$ (to 1 d.p.), and the interval width $0.03125<0.1$. Hence
\[ \boxed{\alpha\approx 1.6\ \text{(1 d.p.)}} \]
(The more precise value is $\alpha\approx 1.5688$.)
\end{exoresolu}

\begin{exoresolu}[How many bisection steps are needed?]
A root of $f(x)=0$ is known to lie in an interval of width $1$. How many bisection iterations are required to locate the root with an error of at most $10^{-3}$, quoting the midpoint as the approximation?
\tcblower
\textbf{Solution.} After $n$ bisections the interval width is $\dfrac{1}{2^{n}}$, and the midpoint differs from the root by at most half of this:
\[ \text{error}\leq \frac{1}{2^{n+1}}. \]
We require $\dfrac{1}{2^{n+1}}\leq 10^{-3}$, i.e.\ $2^{n+1}\geq 1000$.

Since $2^{10}=1024$, we need $n+1\geq 10$, so $n\geq 9$.

\emph{Check:} after $9$ steps, width $=\frac{1}{512}\approx 0.00195$, half-width $\approx 0.00098<10^{-3}$. \checkmark

Hence $\boxed{9\text{ iterations}}$ suffice.
\end{exoresolu}

\begin{exoresolu}[Bisection to find $\sqrt{2}$]
Use the bisection method on $f(x)=x^2-2$ with initial interval $[1,2]$ to find $\sqrt{2}$ correct to 3 decimal places.
\tcblower
\textbf{Solution.} $f(1)=-1<0$, $f(2)=2>0$. Root in $(1,2)$.
\begin{center}
\begin{tabular}{|c|c|c|c|c|}\hline
$n$ & $a$ & $b$ & $m$ & $f(m)$ \\ \hline
1 & 1 & 1.5 & 1.5 & 0.25 \\
2 & 1 & 1.25 & 1.25 & -0.4375 \\
3 & 1.25 & 1.375 & 1.375 & -0.109375 \\
4 & 1.375 & 1.4375 & 1.4375 & 0.066406 \\
5 & 1.375 & 1.40625 & 1.40625 & -0.022461 \\
6 & 1.40625 & 1.421875 & 1.421875 & 0.021729 \\
7 & 1.40625 & 1.4140625 & 1.4140625 & -0.000381 \\
8 & 1.4140625 & 1.41796875 & 1.4140625 & -0.000381 \\ \hline
\end{tabular}
\end{center}
After 7 iterations, interval $[1.4140625, 1.4140625]$? Wait, let's check: at step 7, $a=1.40625$, $b=1.421875$, $m=1.4140625$, $f(m)<0$, so new interval $[1.4140625, 1.421875]$. Width $=0.0078125$. Step 8: $m=1.41796875$, $f(m)>0$, interval $[1.4140625, 1.41796875]$, width $0.0039$. Step 9: $m=1.416015625$, $f(m)>0$, interval $[1.4140625, 1.416015625]$, width $0.00195$. Step 10: $m=1.4150390625$, $f(m)>0$, interval $[1.4140625, 1.4150390625]$, width $0.000976<0.001$. Endpoints round to 1.414. $\boxed{\sqrt{2}\approx 1.414}$.
\end{exoresolu}

\coursec{The Newton--Raphson Method}

\begin{propriete}[Newton--Raphson iteration formula]
For $f$ differentiable near a root $\alpha$, the tangent line at $(x_n, f(x_n))$ cuts the $x$-axis at
\[ x_{n+1} = x_n - \frac{f(x_n)}{f'(x_n)}. \]
If $x_0$ is sufficiently close to $\alpha$ and $f'(\alpha)\neq 0$, the sequence $(x_n)$ converges to $\alpha$ \cle{quadratically} (the number of correct digits roughly doubles each step).
\end{propriete}

\begin{methode}[Applying the Newton--Raphson formula]
To solve $f(x)=0$ starting from an approximation $x_0$ near the root:
\begin{enumerate}
\item Compute $f'(x)$ exactly.
\item Write the iteration formula and simplify it:
\[ x_{n+1}=x_n-\frac{f(x_n)}{f'(x_n)}. \]
\item Choose $x_0$ from a sign-change interval (a sketch or table of values helps).
\item Iterate, keeping \emph{at least two more decimal places} than the final answer requires.
\item Stop when two successive iterates agree to the required accuracy: $|x_{n+1}-x_n|$ smaller than the tolerance.
\item \textbf{Check:} substitute the answer back into $f$; it should be nearly zero.
\end{enumerate}
\textbf{Reflexes:} Work in \emph{radians} for trigonometric $f$; never round intermediate iterates; if $f'(x_n)\approx 0$ (horizontal tangent), the method may fail or jump away — change the starting point.
\end{methode}

\begin{exemplebox}[Geometric interpretation]
\begin{popfigure}
\begin{tikzpicture}[scale=1.2, domain=0.5:2.5]
\draw[->] (0,0) -- (3,0) node[right] {$x$};
\draw[->] (0,-1) -- (0,3) node[above] {$y$};
\draw[popblue, thick] plot[smooth, samples=100] (\x, {\x*\x*\x + 2*\x - 7});
\draw[poporange, dashed] (1.5, -0.625) -- (1.5, 0) node[below] {$x_0$};
\draw[popgreen, dashed] (1.5714, 0.0229) -- (1.5714, 0) node[below] {$x_1$};
\draw[poppurple, dashed] (1.5690, 0) -- (1.5690, 0.0001) node[below] {$x_2$};
\draw[popline] (1.5, -0.625) -- (1.5714, 0);
\draw[popline] (1.5714, 0.0229) -- (1.5690, 0);
\node[popblue] at (2.5, 2) {$y=x^3+2x-7$};
\end{tikzpicture}
\legende{Newton--Raphson tangents converging to the root}
\end{popfigure}
\end{exemplebox}

\begin{exoresolu}[Newton--Raphson on a polynomial]
Use the Newton--Raphson method with $x_0=1.5$ to find the root of $x^{3}+2x-7=0$ correct to $4$ decimal places.
\tcblower
\textbf{Solution.} $f(x)=x^{3}+2x-7$, $f'(x)=3x^{2}+2$. The formula is
\[ x_{n+1}=x_n-\frac{x_n^{3}+2x_n-7}{3x_n^{2}+2}. \]

\emph{Iteration 1:} $x_0=1.5$:
\[ x_1=1.5-\frac{3.375+3-7}{6.75+2}=1.5-\frac{-0.625}{8.75}=1.5+0.071428571=1.571428571. \]

\emph{Iteration 2:} $f(1.571428571)=3.880099+3.142857-7=0.022956$; $f'(1.571428571)=7.408163+2=9.408163$:
\[ x_2=1.571428571-\frac{0.022956}{9.408163}=1.571428571-0.002440=1.568988571. \]

\emph{Iteration 3:} $f(1.568988571)\approx 3.8610+3.13798-7=0.000053$; $f'(1.568988571)\approx 7.3823+2=9.3823$:
\[ x_3=1.568988571-\frac{0.000053}{9.3823}\approx 1.568988571-0.00000565=1.568982921. \]

$x_2$ and $x_3$ agree to 4 decimal places (1.5690). \emph{Check:} $f(1.5690)\approx 0.00008\approx 0$. \checkmark
\[ \boxed{\alpha\approx 1.5690\ \text{(4 d.p.)}} \]
Note the speed: 3 iterations of Newton--Raphson achieve what took bisection many more steps.
\end{exoresolu}

\begin{exoresolu}[Newton--Raphson with a trigonometric equation]
Show that the Newton--Raphson iteration for $x=\cos x$ can be written
\[ x_{n+1}=\frac{x_n\sin x_n+\cos x_n}{1+\sin x_n}, \]
and use it with $x_0=0.75$ to find the root correct to 4 decimal places.
\tcblower
\textbf{Solution.} Write $f(x)=x-\cos x$, so $f'(x)=1+\sin x$. Then
\[ x_{n+1}=x_n-\frac{x_n-\cos x_n}{1+\sin x_n}
=\frac{x_n(1+\sin x_n)-x_n+\cos x_n}{1+\sin x_n}
=\frac{x_n\sin x_n+\cos x_n}{1+\sin x_n}. \qquad\blacksquare \]

\emph{Iteration 1:} $x_0=0.75$: $\sin 0.75=0.68163876$, $\cos 0.75=0.73168887$:
\[ x_1=\frac{0.75(0.68163876)+0.73168887}{1+0.68163876}=\frac{1.24291794}{1.68163876}=0.73911289. \]

\emph{Iteration 2:} $\sin 0.73911289=0.67361203$, $\cos 0.73911289=0.73908513$:
\[ x_2=\frac{0.73911289(0.67361203)+0.73908513}{1.67361203}=\frac{1.236960}{1.67361203}=0.73908513. \]

$x_1$ and $x_2$ agree to 4 d.p. \emph{Check:} $\cos 0.73908513\approx 0.73908513$. \checkmark
\[ \boxed{\alpha\approx 0.7391\ \text{(4 d.p.)}} \]
\end{exoresolu}

\begin{exoresolu}[Estimating a square root by Newton--Raphson]
By applying the Newton--Raphson method to $x^{2}-11=0$, show that the iteration is $x_{n+1}=\dfrac{1}{2}\left(x_n+\dfrac{11}{x_n}\right)$. Starting from $x_0=3$, estimate $\sqrt{11}$ correct to 5 decimal places.
\tcblower
\textbf{Solution.} $f(x)=x^{2}-11$, $f'(x)=2x$:
\[ x_{n+1}=x_n-\frac{x_n^{2}-11}{2x_n}=\frac{2x_n^{2}-x_n^{2}+11}{2x_n}=\frac{1}{2}\left(x_n+\frac{11}{x_n}\right). \qquad\blacksquare \]

\emph{Iteration 1:} $x_1=\tfrac12\left(3+\tfrac{11}{3}\right)=\tfrac12(3+3.66666667)=3.33333333$.

\emph{Iteration 2:} $x_2=\tfrac12\left(3.33333333+\tfrac{11}{3.33333333}\right)=\tfrac12(3.33333333+3.30000000)=3.31666667$.

\emph{Iteration 3:} $x_3=\tfrac12\left(3.31666667+\tfrac{11}{3.31666667}\right)=\tfrac12(3.31666667+3.31657895)=3.31662281$.

\emph{Iteration 4:} $x_4=\tfrac12(3.31662281+3.31662479)=3.31662380$.

\emph{Iteration 5:} $x_5=\tfrac12(3.31662380+3.31662380)=3.31662380$ (stabilised).

\emph{Check:} $(3.31662380)^{2}=10.9999999\approx 11$. \checkmark
\[ \boxed{\sqrt{11}\approx 3.31662\ \text{(5 d.p.)}} \]
\end{exoresolu}

\begin{exoresolu}[Newton--Raphson for a transcendental equation]
Find the root of $e^{-x}=x$ correct to 4 decimal places using Newton--Raphson with $x_0=0.5$.
\tcblower
\textbf{Solution.} $f(x)=e^{-x}-x$, $f'(x)=-e^{-x}-1$.
\[ x_{n+1}=x_n-\frac{e^{-x_n}-x_n}{-e^{-x_n}-1}=x_n+\frac{e^{-x_n}-x_n}{e^{-x_n}+1}. \]
$x_0=0.5$: $e^{-0.5}=0.60653066$
$x_1=0.5+\frac{0.60653066-0.5}{0.60653066+1}=0.5+\frac{0.10653066}{1.60653066}=0.5+0.066310=0.566310$
$x_2=0.566310+\frac{e^{-0.566310}-0.566310}{e^{-0.566310}+1}=0.566310+\frac{0.567628-0.566310}{1.567628}=0.566310+0.000841=0.567151$
$x_3=0.567151+\frac{e^{-0.567151}-0.567151}{e^{-0.567151}+1}=0.567151+\frac{0.567143-0.567151}{1.567143}=0.567151-0.000005=0.567146$
$x_4=0.567146$ (agrees to 6 d.p.). $\boxed{\alpha\approx 0.5671\ \text{(4 d.p.)}}$
\end{exoresolu}

\coursec{Accuracy, Comparison and Applications}

\begin{methode}[Controlling accuracy and choosing a method]
\begin{enumerate}
\item \textbf{Error bound for bisection:} after $n$ steps from $[a,b]$, the midpoint approximates the root with error at most $\dfrac{b-a}{2^{n+1}}$.
\item \textbf{Stopping criterion:} iterate until two successive approximations agree to the required number of decimal places, then \emph{verify} by substitution ($|f(x_n)|$ small).
\item \textbf{Newton--Raphson fails or misbehaves} when $f'(x_n)=0$ (or is tiny), when the starting point is far from the root, or near a stationary point or asymptote. Bisection fails to \emph{detect} a root when there is no sign change (e.g.\ a repeated root like $f(x)=(x-1)^2$).
\item \textbf{Choice:} bisection is slow (linear convergence) but \emph{guaranteed} once a sign change is found; Newton--Raphson is very fast (quadratic convergence) but needs a good starting point and a computable derivative.
\item \textbf{Exam habit:} always state the equation in the form $f(x)=0$, the formula used, each iterate to full calculator precision, and a final check.
\end{enumerate}
\end{methode}

\begin{exoresolu}[A case where Newton--Raphson fails]
Attempt the Newton--Raphson method on $f(x)=x^{3}-3x+2=0$ starting from $x_0=0$. Explain the difficulty, and find the roots by another means.
\tcblower
\textbf{Solution.} $f'(x)=3x^{2}-3$. With $x_0=0$:
\[ x_1=0-\frac{f(0)}{f'(0)}=-\frac{2}{-3}=\frac{2}{3}\approx 0.6667. \]
Then $x_2=0.6667-\dfrac{f(0.6667)}{f'(0.6667)}=0.6667-\dfrac{0.2963}{-1.6667}=0.6667+0.1778=0.8444$; continuing gives $x_3\approx 0.9355$, $x_4\approx 0.9723$, \dots\ converging \emph{slowly} towards $1$.

\emph{Explanation:} $f'(1)=0$, so $x=1$ is a \textbf{repeated root} (stationary point of $f$ on the $x$-axis). Near a repeated root the derivative is tiny, the correction term $f/f'$ becomes unreliable, and Newton--Raphson loses its quadratic speed (convergence becomes linear). Note also that there is \emph{no sign change} at $x=1$, so bisection could not even detect this root.

\emph{Exact solution:} factorise. Since $f(1)=0$, $(x-1)$ is a factor:
\[ x^{3}-3x+2=(x-1)(x^{2}+x-2)=(x-1)^{2}(x+2). \]
Hence the roots are $\boxed{x=1\ \text{(double)}\ \text{and}\ x=-2}$.
\end{exoresolu}

\begin{exoresolu}[Newton--Raphson divergence example]
Show that Newton--Raphson on $f(x)=x^{1/3}$ (i.e. $x=0$) with any $x_0\neq 0$ diverges.
\tcblower
\textbf{Solution.} $f(x)=x^{1/3}$, $f'(x)=\frac{1}{3}x^{-2/3}$.
\[ x_{n+1}=x_n-\frac{x_n^{1/3}}{\frac{1}{3}x_n^{-2/3}}=x_n-3x_n=-2x_n. \]
Thus $x_n=(-2)^n x_0$, which oscillates and diverges in magnitude. The method fails because $f'(0)$ is infinite (vertical tangent at the root).
\end{exoresolu}

\begin{exoresolu}[Comparing methods on the same equation — GCE style]
The equation $x^{3}-x-3=0$ has a root $\alpha$ in $(1,2)$.
\begin{enumerate}
\item Verify this statement.
\item Use \textbf{two} iterations of the bisection method to narrow the interval.
\item Use the Newton--Raphson method with $x_0=1.5$ to find $\alpha$ correct to 3 decimal places.
\item State, with a reason, which method reaches 3 d.p.\ faster.
\end{enumerate}
\tcblower
\textbf{Solution.} Let $f(x)=x^{3}-x-3$, continuous everywhere.

\emph{(a)} $f(1)=1-1-3=-3<0$; $f(2)=8-2-3=3>0$. Sign change $\Rightarrow$ a root in $(1,2)$. \checkmark

\emph{(b) Bisection.} Midpoint $m_1=1.5$: $f(1.5)=3.375-1.5-3=-1.125<0$ $\Rightarrow$ root in $(1.5,2)$. Midpoint $m_2=1.75$: $f(1.75)=5.359375-1.75-3=0.609375>0$ $\Rightarrow$ root in $(1.5,\,1.75)$.

\emph{(c) Newton--Raphson.} $f'(x)=3x^{2}-1$, so $x_{n+1}=x_n-\dfrac{x_n^{3}-x_n-3}{3x_n^{2}-1}$.
\[ x_1=1.5-\frac{-1.125}{5.75}=1.5+0.195652174=1.695652174. \]
\[ f(1.695652174)=4.8749-1.69565-3=0.17925,\quad f'(1.695652174)=8.6273-1=7.6273, \]
\[ x_2=1.695652174-\frac{0.17925}{7.6273}=1.695652174-0.023494=1.672158. \]
\[ f(1.672158)=4.6748-1.67216-3=0.00264,\quad f'(1.672158)=8.376-1=7.376, \]
\[ x_3=1.672158-\frac{0.00264}{7.376}=1.672158-0.000358=1.671800. \]
$x_2$ and $x_3$ agree to 3 d.p. (1.672). \emph{Check:} $f(1.6718)\approx -0.0002\approx 0$. \checkmark\quad $\boxed{\alpha\approx 1.672\ (3\ \text{d.p.})}$

\emph{(d)} Newton--Raphson is faster: it reached 3 d.p.\ in 3 iterations, while bisection would need the interval width $\frac{1}{2^n}<10^{-3}$, i.e.\ $n\geq 10$ iterations. Bisection, however, required no derivative and was guaranteed to converge from the sign change.
\end{exoresolu}

\begin{experience}[Numerical solution of a practical problem]
A cylindrical tank of radius $r=2$ m and length $L=10$ m lies horizontally. The volume of liquid when the depth is $h$ is $V = L\left(r^2\cos^{-1}\left(\frac{r-h}{r}\right) - (r-h)\sqrt{2rh-h^2}\right)$. If the tank contains $50\ \mathrm{m}^3$ of liquid, find the depth $h$ correct to 2 decimal places using Newton--Raphson.
\tcblower
\textbf{Solution.} Equation: $f(h)=20\left(4\cos^{-1}\left(1-\frac{h}{2}\right) - (2-h)\sqrt{4h-h^2}\right)-50=0$, $h\in(0,4)$.
$f(1)\approx 20(4\cdot1.047 - 1\cdot\sqrt{3})-50=20(4.188-1.732)-50=-0.88<0$.
$f(1.5)\approx 20(4\cdot0.723 - 0.5\cdot\sqrt{3.75})-50=20(2.892-0.968)-50=38.48-50=-11.52$? Wait, recalculate.
Better: $f(1)=20(4\cos^{-1}(0.5) - 1\cdot\sqrt{3})-50=20(4\cdot1.0472 - 1.732)-50=20(4.1888-1.732)-50=20(2.4568)-50=-0.864$.
$f(2)=20(4\cos^{-1}(0) - 0)-50=20(4\cdot1.5708)-50=125.66-50=75.66>0$.
Root in $(1,2)$. Use $h_0=1.2$. Derivative $f'(h)$ computed numerically or analytically. Iterations yield $h\approx 1.02$ m. (Full calculation omitted for brevity; in exam, show 3 iterations).
\end{experience}

\begin{attention}[Frequent examination traps]
\begin{itemize}
\item Forgetting to rearrange the equation into the form $f(x)=0$ before looking for a sign change.
\item Claiming a root exists without checking \emph{continuity} (e.g.\ $f(x)=1/x$ changes sign between $-1$ and $1$ but has no root there!).
\item Rounding intermediate iterates: keep at least 6 decimal places until the final line.
\item Using degrees instead of \cle{radians} in trigonometric equations.
\item Stopping Newton--Raphson after one iteration ``because it looks close'' — always confirm that two successive iterates agree to the required accuracy.
\item Applying bisection to a repeated root (no sign change) or Newton--Raphson where $f'(x_n)=0$.
\item Not simplifying the Newton--Raphson formula before iterating (wastes time and introduces rounding errors).
\item Forgetting the final verification step: $f(\text{answer})\approx 0$.
\end{itemize}
\end{attention}

\begin{aretenir}[Key points for Numerical Methods]
\begin{itemize}
\item \textbf{Location of roots:} IVT + continuity + sign change $\Rightarrow$ existence; monotonicity $\Rightarrow$ uniqueness.
\item \textbf{Bisection:} guaranteed, linear convergence, error $\leq \frac{b-a}{2^{n+1}}$, needs sign change.
\item \textbf{Newton--Raphson:} $x_{n+1}=x_n-\frac{f(x_n)}{f'(x_n)}$, quadratic convergence, needs $f'$ and good $x_0$, fails at stationary points.
\item \textbf{Accuracy:} always work to 2 extra decimal places; stop when successive iterates agree to required d.p.; verify by substitution.
\item \textbf{Repeated roots:} no sign change (bisection fails), $f'(\alpha)=0$ (Newton--Raphson slows down).
\end{itemize}
\end{aretenir}

\begin{exercice}[Practice exercises]
\begin{enumerate}
\item Show that $x^3-2x-5=0$ has a root in $(2,3)$. Use bisection to find it to 2 d.p.
\item Use Newton--Raphson with $x_0=2$ to solve $x^3-2x-5=0$ to 4 d.p. Compare the number of iterations with bisection.
\item The equation $\ln x = 4-x$ has a root near $3$. Find it to 3 d.p. using Newton--Raphson.
\item Explain why Newton--Raphson fails for $f(x)=x^3-2x+2$ with $x_0=0$. (Hint: cycle).
\item A root of $f(x)=0$ lies in $[0,1]$. How many bisection steps to guarantee error $<10^{-5}$?
\item Find $\sqrt[3]{20}$ to 5 d.p. using Newton--Raphson on $x^3-20=0$ with $x_0=3$.
\item Show that $e^x=3x$ has two roots. Locate each to width $0.25$ and find the larger root to 3 d.p.
\end{enumerate}
\end{exercice}

\begin{corrige}[Exercise 1]
$f(x)=x^3-2x-5$. $f(2)=-1<0$, $f(3)=16>0$ $\Rightarrow$ root in $(2,3)$.
Bisection: $m=2.5$, $f=15.625-5-5=5.625>0$ $\Rightarrow$ $(2,2.5)$.
$m=2.25$, $f=11.39-4.5-5=1.89>0$ $\Rightarrow$ $(2,2.25)$.
$m=2.125$, $f=9.59-4.25-5=0.34>0$ $\Rightarrow$ $(2,2.125)$.
$m=2.0625$, $f=8.78-4.125-5=-0.345<0$ $\Rightarrow$ $(2.0625,2.125)$.
$m=2.09375$, $f=9.18-4.1875-5=-0.0075<0$ $\Rightarrow$ $(2.09375,2.125)$.
Width $0.03125<0.01$? No, need width $<0.01$ for 2 d.p. Continue: $m=2.109375$, $f>0$ $\Rightarrow$ $(2.09375,2.109375)$, width $0.0156$. $m=2.1015625$, $f>0$ $\Rightarrow$ $(2.09375,2.1015625)$, width $0.0078<0.01$. Endpoints round to 2.09 and 2.10? Wait: 2.09375 rounds to 2.09, 2.1015625 rounds to 2.10. Not same. Need one more: $m=2.09765625$, $f<0$ $\Rightarrow$ $(2.09765625,2.1015625)$, both round to 2.10. $\boxed{2.10}$.
\end{corrige}

\begin{corrige}[Exercise 2]
$f(x)=x^3-2x-5$, $f'(x)=3x^2-2$. $x_{n+1}=x_n-\frac{x_n^3-2x_n-5}{3x_n^2-2}$.
$x_0=2$: $x_1=2-\frac{-1}{10}=2.1$.
$x_1=2.1$: $f=9.261-4.2-5=0.061$, $f'=13.23-2=11.23$, $x_2=2.1-\frac{0.061}{11.23}=2.09457$.
$x_2=2.09457$: $f\approx 0.00018$, $f'\approx 11.16$, $x_3=2.09457-0.000016=2.09455$.
Agrees to 4 d.p.: $\boxed{2.0946}$. Bisection needed ~7 iterations; Newton--Raphson 3 iterations.
\end{corrige}

\begin{corrige}[Exercise 3]
$f(x)=\ln x + x - 4$. $f'(x)=\frac{1}{x}+1$. $x_{n+1}=x_n-\frac{\ln x_n+x_n-4}{1/x_n+1}$.
$x_0=3$: $f(3)=1.0986+3-4=0.0986$, $f'=1.3333$, $x_1=3-\frac{0.0986}{1.3333}=2.9260$.
$x_1=2.9260$: $f=1.0735+2.9260-4=-0.0005$, $f'=1.3417$, $x_2=2.9260+\frac{0.0005}{1.3417}=2.9264$.
$x_2=2.9264$: $f\approx 0$. $\boxed{2.926}$.
\end{corrige}

\begin{corrige}[Exercise 4]
$f(x)=x^3-2x+2$, $f'(x)=3x^2-2$. $x_0=0$: $x_1=0-\frac{2}{-2}=1$. $x_1=1$: $x_2=1-\frac{1}{1}=0$. Cycle: $0,1,0,1,\dots$ never converges. The tangent at $x=0$ hits $x=1$, tangent at $x=1$ hits $x=0$.
\end{corrige}

\begin{corrige}[Exercise 5]
Error $\leq \frac{1}{2^{n+1}} < 10^{-5}$ $\Rightarrow$ $2^{n+1} > 10^5 = 100000$. $2^{17}=131072$, so $n+1=17$, $n=16$ iterations.
\end{corrige}

\begin{corrige}[Exercise 6]
$f(x)=x^3-20$, $f'(x)=3x^2$. $x_{n+1}=x_n-\frac{x_n^3-20}{3x_n^2}=\frac{2}{3}x_n+\frac{20}{3x_n^2}$.
$x_0=3$: $x_1=\frac{2}{3}(3)+\frac{20}{27}=2+0.74074=2.74074$.
$x_2=\frac{2}{3}(2.74074)+\frac{20}{3(2.74074)^2}=1.82716+0.8876=2.71476$.
$x_3=2.71442$, $x_4=2.71442$. $\boxed{2.71442}$.
\end{corrige}

\begin{corrige}[Exercise 7]
$f(x)=e^x-3x$. $f(0)=1>0$, $f(1)=e-3=-0.28<0$ $\Rightarrow$ root in $(0,1)$.
$f(1)<0$, $f(2)=e^2-6=7.389-6=1.389>0$ $\Rightarrow$ root in $(1,2)$.
Two roots. Larger root in $(1,2)$. Width 0.25: test 1.25, 1.5, 1.75.
$f(1.25)=3.490-3.75=-0.26$, $f(1.5)=4.482-4.5=-0.018$, $f(1.75)=5.755-5.25=0.505$.
Root in $(1.5,1.75)$ width 0.25.
Newton--Raphson: $f'(x)=e^x-3$. $x_{n+1}=x_n-\frac{e^{x_n}-3x_n}{e^{x_n}-3}$.
$x_0=1.5$: $x_1=1.5-\frac{-0.018}{1.482}=1.5121$.
$x_2=1.5121-\frac{e^{1.5121}-4.5363}{e^{1.5121}-3}=1.5121-\frac{0.0002}{1.532}=1.5121$.
$\boxed{1.512}$.
\end{corrige}

\newpage

\coursec{Exercises}

\courssub{I check my knowledge}

\begin{exercice}[Multiple choice question]
The equation $f(x)=0$ is known to have a root $\alpha$ in the interval $[1,2]$, and $f$ is continuous on this interval. After applying the bisection method $n$ times, the root is known to lie in an interval of width:
\begin{enumerate}
\item $\dfrac{1}{n}$ \qquad
\item $\dfrac{1}{2^{n}}$ \qquad
\item $\dfrac{1}{2^{n-1}}$ \qquad
\item $\dfrac{1}{n^{2}}$
\end{enumerate}
Choose the correct answer, justifying your choice.
\end{exercice}

\begin{exercice}[True or false]
For each of the following statements, say whether it is \textbf{true} or \textbf{false}, justifying your answer carefully.
\begin{enumerate}
\item If $f$ is continuous on $[a,b]$ and $f(a)f(b)<0$, then the equation $f(x)=0$ has \emph{exactly one} root in $(a,b)$.
\item The Newton--Raphson iteration formula for solving $f(x)=0$ is $x_{n+1}=x_{n}-\dfrac{f(x_{n})}{f'(x_{n})}$.
\item The bisection method always converges to a root whenever $f$ is continuous on $[a,b]$ with $f(a)f(b)<0$.
\item The Newton--Raphson method converges for every choice of starting value $x_{0}$.
\end{enumerate}
\end{exercice}

\begin{exercice}[Definitions]
\begin{enumerate}
\item State the \emph{Intermediate Value Theorem} and explain how it is used to locate a root of a continuous function in an interval.
\item Explain what is meant by saying that a root $\alpha$ has been found ``correct to $2$ decimal places''.
\item Give one advantage and one disadvantage of the bisection method compared with the Newton--Raphson method.
\end{enumerate}
\end{exercice}

\begin{exercice}[Direct calculations]
\begin{enumerate}
\item Given $f(x)=x^{3}+x-3$, compute $f(1)$ and $f(2)$. What can you deduce?
\item For $f(x)=x^{2}-7$, write down the Newton--Raphson iteration formula and compute $x_{1}$ starting from $x_{0}=3$, giving your answer as an exact fraction.
\item Starting from the interval $[1,2]$, how many bisection iterations are needed to obtain an interval of width less than $0.1$?
\end{enumerate}
\end{exercice}

\courssub{I practise}

\begin{exercice}[Locating a root and bisection]
\begin{enumerate}
\item Show that the equation $x^{3}+3x-5=0$ has exactly one real root and that it lies between $1$ and $2$.
\item Use the bisection method three times to narrow down the interval containing the root.
\item Hence state the value of the root to $1$ decimal place, justifying your answer.
\end{enumerate}
\end{exercice}

\begin{exercice}[Newton--Raphson in action]
\begin{enumerate}
\item Show that the equation $x^{3}-2x-5=0$ has a root between $2$ and $3$.
\item Use the Newton--Raphson method with $x_{0}=2$ to find this root correct to $4$ decimal places, showing all iterates.
\item Verify, using a suitable sign-change test, that your answer is indeed correct to $4$ decimal places.
\end{enumerate}
\end{exercice}

\begin{exercice}[Graphical location of roots]
\begin{enumerate}
\item By sketching the curves $y=\mathrm{e}^{x}$ and $y=4-x^{2}$ on the same axes, determine the number of real roots of the equation $\mathrm{e}^{x}+x^{2}-4=0$.
\item Locate each root in an interval between two consecutive integers, justifying with a sign-change test.
\item Choose one of these roots and refine it using two applications of the bisection method.
\end{enumerate}
\end{exercice}

\begin{exercice}[The equation $\cos x = x$]
The equation $\cos x = x$ (with $x$ in radians) has a root near $0.7$.
\begin{enumerate}
\item Write the equation in the form $f(x)=0$ and write down the Newton--Raphson iteration formula for it.
\item Apply the Newton--Raphson method twice, starting from $x_{0}=0.7$.
\item Verify, using a sign-change test on a suitable interval of width $0.001$, that your answer is correct to $3$ decimal places.
\end{enumerate}
\end{exercice}

\courssub{I prepare for the GCE}

\begin{exercice}[Iteration, staircase diagram and a square-root formula]
The equation $x^{3}-5x+2=0$ is rearranged in the form $x=\dfrac{x^{3}+2}{5}$.
\begin{enumerate}
\item Show that the equation has a root between $0$ and $1$. \hfill (2 marks)
\item Starting from $x_{0}=1$, perform three iterations of $x_{n+1}=\dfrac{x_{n}^{3}+2}{5}$, giving each value to $4$ decimal places. \hfill (3 marks)
\item Illustrate the convergence of this iteration with a staircase diagram on a sketch of $y=x$ and $y=\dfrac{x^{3}+2}{5}$, and comment on the behaviour of the iterates. \hfill (3 marks)
\item Show that applying the Newton--Raphson method to the equation $x^{2}-N=0$ gives the formula
\[ x_{n+1}=\tfrac{1}{2}\left(x_{n}+\dfrac{N}{x_{n}}\right). \] \hfill (3 marks)
\item Hence, taking $N=11$ and a suitable starting value, compute $\sqrt{11}$ correct to $5$ decimal places, showing all iterates. \hfill (4 marks)
\end{enumerate}
\end{exercice}

\begin{exercice}[Failure of Newton--Raphson and error bounds]
\begin{enumerate}
\item Explain, with the aid of a sketch showing tangents, why the Newton--Raphson method fails to find the root of $x^{3}-2x+2=0$ when started at $x_{0}=0$. \hfill (4 marks)
\item A root of $f(x)=0$ is known to lie in the interval $[1,2]$, where $f$ is continuous. Using the error bound for the bisection method, determine the number of iterations needed to guarantee that the root is found with an error less than $10^{-4}$. Justify your answer carefully. \hfill (4 marks)
\item The outflow time $t$ (in minutes) of water from a tank satisfies
\[ t\,\mathrm{e}^{-t/10}=2. \]
Show that this equation has a root between $3$ and $4$, and refine it by the bisection method until the error is less than $0.05$. Interpret your answer in the context of the problem. \hfill (7 marks)
\end{enumerate}
\end{exercice}

\begin{exercice}[GCE synthesis: the floating buoy]
A spherical buoy of radius $1\ \mathrm{m}$ and density $0.6$ times that of water floats at a depth $d$ (in metres), where $d$ satisfies the equation
\[ d^{3}-3d^{2}+2.4=0. \]
\begin{enumerate}
\item Explain why the physically meaningful root must satisfy $0<d<2$. \hfill (2 marks)
\item Show that the equation has a root between $1$ and $2$, and locate it more precisely in an interval of width $0.5$. \hfill (3 marks)
\item Write down the Newton--Raphson iteration formula for this equation. \hfill (2 marks)
\item Starting from a suitable value $x_{0}$, find the root correct to $3$ decimal places, showing all iterates. \hfill (5 marks)
\item Verify your answer with a sign-change test, and hence state the depth at which the buoy floats, giving your answer to the nearest centimetre. \hfill (3 marks)
\end{enumerate}
\end{exercice}

\newpage

\coursec{Solutions to the exercises}

\begin{corrige}[Exercise 1 — Multiple choice question]
The correct answer is \textbf{(b)} $\dfrac{1}{2^{n}}$.

\textbf{Justification.} The initial interval $[1,2]$ has width $b-a=2-1=1$. At each application of the bisection method, the interval is replaced by one of its two halves, so the width is divided by $2$. After $n$ iterations the width is therefore
\[
\frac{b-a}{2^{n}}=\frac{1}{2^{n}}.
\]
Answer (c) would be the width after $n-1$ iterations, and answers (a) and (d) do not reflect the halving process.
\end{corrige}

\begin{corrige}[Exercise 2 — True or false]
\begin{enumerate}
\item \textbf{False.} The Intermediate Value Theorem guarantees the existence of \emph{at least one} root in $(a,b)$, not exactly one. For example, $f(x)=x^{3}-x$ on $[-2,2]$ is continuous with $f(-2)f(2)=(-6)(6)<0$, yet it has \emph{three} roots $-1,0,1$ in $(-2,2)$. Uniqueness requires an extra condition, e.g.\ $f$ strictly monotonic on $[a,b]$.
\item \textbf{True.} This is exactly the Newton--Raphson formula, obtained by taking the tangent to $y=f(x)$ at $(x_{n},f(x_{n}))$ and finding where it cuts the $x$-axis: $0-f(x_{n})=f'(x_{n})(x_{n+1}-x_{n})$.
\item \textbf{True.} The bisection method only requires continuity and a sign change. At each step the sign change is preserved by construction, so the intervals shrink to a root: convergence is \emph{guaranteed} (though possibly slow).
\item \textbf{False.} Convergence depends on the starting value. If $f'(x_{0})=0$ the tangent is horizontal and the iteration is undefined; the iterates may also cycle or diverge (see Exercise 10(a), where $x^{3}-2x+2=0$ started at $x_{0}=0$ cycles between $0$ and $1$).
\end{enumerate}
\end{corrige}

\begin{corrige}[Exercise 3 — Definitions]
\begin{enumerate}
\item \textbf{Intermediate Value Theorem.} If $f$ is continuous on $[a,b]$ and $k$ lies between $f(a)$ and $f(b)$, then there exists $c\in(a,b)$ with $f(c)=k$.

\emph{Use for locating roots:} if $f(a)$ and $f(b)$ have opposite signs, i.e.\ $f(a)f(b)<0$, then $k=0$ lies between $f(a)$ and $f(b)$, so there is at least one root of $f(x)=0$ in $(a,b)$. A \emph{sign change} over an interval therefore locates a root.
\item Saying a root $\alpha$ is found ``correct to $2$ decimal places'' means we have an approximation $x$ with $|x-\alpha|<0.005$; equivalently, $\alpha$ is trapped in an interval $(p-0.005,p+0.005)$, so that $\alpha$ rounds to the same $2$ d.p.\ value $p$. In practice one checks a sign change over $(p-0.005,p+0.005)$.
\item \emph{Advantage of bisection:} it always converges provided $f$ is continuous with a sign change — no derivative and no careful choice of starting point is needed. \emph{Disadvantage:} it converges slowly (only one binary digit per iteration), whereas Newton--Raphson typically doubles the number of correct digits at each step when it converges.
\end{enumerate}
\end{corrige}

\begin{corrige}[Exercise 4 — Direct calculations]
\begin{enumerate}
\item $f(1)=1+1-3=-1<0$ and $f(2)=8+2-3=7>0$. Since $f$ is a polynomial (hence continuous) and changes sign on $[1,2]$, the equation $x^{3}+x-3=0$ has at least one root in $(1,2)$.
\item $f(x)=x^{2}-7$, $f'(x)=2x$, so
\[
x_{n+1}=x_{n}-\frac{x_{n}^{2}-7}{2x_{n}}.
\]
With $x_{0}=3$: $\;x_{1}=3-\dfrac{9-7}{6}=3-\dfrac{1}{3}=\dfrac{8}{3}$.
\item After $n$ iterations the width is $\dfrac{1}{2^{n}}$. We need $\dfrac{1}{2^{n}}<0.1$, i.e.\ $2^{n}>10$. Since $2^{3}=8<10$ and $2^{4}=16>10$, the number of iterations needed is $\boxed{n=4}$.
\end{enumerate}
\end{corrige}

\begin{corrige}[Exercise 5 — Locating a root and bisection]
\begin{enumerate}
\item Let $f(x)=x^{3}+3x-5$. Then $f'(x)=3x^{2}+3=3(x^{2}+1)>0$ for all $x$, so $f$ is \emph{strictly increasing} on $\mathbb{R}$: the equation has \emph{at most one} real root. Moreover
\[
f(1)=1+3-5=-1<0,\qquad f(2)=8+6-5=9>0,
\]
so by the Intermediate Value Theorem there is at least one root in $(1,2)$. Hence there is \emph{exactly one} real root and it lies between $1$ and $2$.
\item \textbf{Bisection, three steps:}
\begin{center}
\begin{tabular}{|c|c|c|c|c|}
\hline
\rowcolor{popblueL} Step & Interval & Midpoint $m$ & $f(m)$ & New interval \\
\hline
1 & $[1,2]$ & $1.5$ & $1.5^{3}+4.5-5=2.875>0$ & $[1,1.5]$ \\
\hline
2 & $[1,1.5]$ & $1.25$ & $1.953125+3.75-5=0.703>0$ & $[1,1.25]$ \\
\hline
3 & $[1,1.25]$ & $1.125$ & $1.4238+3.375-5=-0.201<0$ & $[1.125,1.25]$ \\
\hline
\end{tabular}
\end{center}
\item The root lies in $[1.125,1.25]$. To decide the value to $1$ d.p., test the sign at $1.15$:
\[
f(1.15)=1.520875+3.45-5=-0.029<0,\qquad f(1.25)=0.703>0,
\]
so the root lies in $(1.15,1.25)$. Every number in this interval rounds to $1.2$ except values below $1.15$, which are excluded. Hence the root is $\boxed{1.2}$ to $1$ decimal place.
\end{enumerate}
\end{corrige}

\begin{corrige}[Exercise 6 — Newton--Raphson in action]
\begin{enumerate}
\item Let $f(x)=x^{3}-2x-5$. $f(2)=8-4-5=-1<0$ and $f(3)=27-6-5=16>0$. Since $f$ is continuous, there is a root in $(2,3)$.
\item $f'(x)=3x^{2}-2$, so the iteration is
\[
x_{n+1}=x_{n}-\frac{x_{n}^{3}-2x_{n}-5}{3x_{n}^{2}-2}.
\]
\begin{center}
\begin{tabular}{|c|c|c|c|c|}
\hline
\rowcolor{popblueL} $n$ & $x_{n}$ & $f(x_{n})$ & $f'(x_{n})$ & $x_{n+1}$ \\
\hline
0 & $2$ & $-1$ & $10$ & $2.1$ \\
\hline
1 & $2.1$ & $0.061$ & $11.23$ & $2.094568$ \\
\hline
2 & $2.094568$ & $0.000185$ & $11.161646$ & $2.094551$ \\
\hline
3 & $2.094551$ & $\approx 0$ & --- & $2.094551$ \\
\hline
\end{tabular}
\end{center}
The iterates have stabilised: the root is $\boxed{2.0946}$ to $4$ decimal places.
\item \textbf{Sign-change test.} Test the interval $(2.09455,\,2.09465)$:
\[
f(2.09455)\approx -4.5\times10^{-5}<0,\qquad f(2.09465)\approx +1.07\times10^{-3}>0.
\]
The sign change confirms a root in $(2.09455,2.09465)$, and every number in this interval rounds to $2.0946$. Hence the answer is correct to $4$ decimal places.
\end{enumerate}
\end{corrige}

\begin{corrige}[Exercise 7 — Graphical location of roots]
\begin{enumerate}
\item The equation $\mathrm{e}^{x}+x^{2}-4=0$ is equivalent to $\mathrm{e}^{x}=4-x^{2}$. Sketch $y=\mathrm{e}^{x}$ (increasing exponential through $(0,1)$) and $y=4-x^{2}$ (parabola, vertex $(0,4)$, zeros at $x=\pm2$).
\begin{popfigure}
\begin{center}
\begin{tikzpicture}
\begin{axis}[width=9cm,height=7cm,axis lines=middle,xmin=-3,xmax=3,ymin=-1,ymax=6,xlabel={$x$},ylabel={$y$},xtick={-2,-1,1,2},ytick={1,2,3,4,5}]
\addplot[domain=-2.4:1.8,thick,popblue] {exp(x)};
\addplot[domain=-2.6:2.6,thick,poporange] {4-x^2};
\addlegendentry{$y=\mathrm{e}^{x}$}
\addlegendentry{$y=4-x^{2}$}
\end{axis}
\end{tikzpicture}
\end{center}
\legende{The curves intersect twice, so the equation has two real roots.}
\end{popfigure}
The curves intersect \textbf{twice}, so the equation has \textbf{two real roots}.
\item Let $f(x)=\mathrm{e}^{x}+x^{2}-4$.
\[
f(1)=\mathrm{e}-3\approx-0.282<0,\quad f(2)=\mathrm{e}^{2}\approx7.389>0 \;\Rightarrow\; \text{root in }(1,2).
\]
\[
f(-2)=\mathrm{e}^{-2}\approx0.135>0,\quad f(-1)=\mathrm{e}^{-1}-3\approx-2.632<0 \;\Rightarrow\; \text{root in }(-2,-1).
\]
\item Refine the root in $(1,2)$ by bisection:
\[
f(1.5)=\mathrm{e}^{1.5}+2.25-4\approx2.732>0 \;\Rightarrow\; \text{root in }(1,1.5);
\]
\[
f(1.25)=\mathrm{e}^{1.25}+1.5625-4\approx1.053>0 \;\Rightarrow\; \text{root in }(1,1.25).
\]
After two applications the root lies in $(1,1.25)$, of width $0.25$.
\end{enumerate}
\end{corrige}

\begin{corrige}[Exercise 8 — The equation $\cos x = x$]
\begin{enumerate}
\item Write $f(x)=\cos x - x =0$. Then $f'(x)=-\sin x -1$, and the Newton--Raphson formula is
\[
x_{n+1}=x_{n}-\frac{\cos x_{n}-x_{n}}{-\sin x_{n}-1}
=x_{n}+\frac{\cos x_{n}-x_{n}}{\sin x_{n}+1}.
\]
\item With $x_{0}=0.7$ (radians):
\[
f(0.7)=\cos 0.7-0.7=0.764842-0.7=0.064842,\quad
f'(0.7)=-\sin0.7-1=-1.644218,
\]
\[
x_{1}=0.7-\frac{0.064842}{-1.644218}=0.739437.
\]
Then $f(0.739437)\approx-0.000115$, $f'(0.739437)\approx-1.673847$, so
\[
x_{2}=0.739437-\frac{-0.000115}{-1.673847}=0.739368.
\]
The root is approximately $0.7391$ to $4$ decimal places ($0.739$ to $3$ d.p.).
\item \textbf{Sign-change test} on $(0.7385,\,0.7395)$:
\[
f(0.7385)=\cos(0.7385)-0.7385\approx+0.00060>0,
\]
\[
f(0.7395)=\cos(0.7395)-0.7395\approx-0.00041<0.
\]
The sign change confirms a root in $(0.7385,0.7395)$, so the root is $\boxed{0.739}$ correct to $3$ decimal places.
\end{enumerate}
\end{corrige}

\begin{corrige}[Exercise 9 — Iteration, staircase diagram and a square-root formula]
\begin{enumerate}
\item Let $f(x)=x^{3}-5x+2$. $f(0)=2>0$ and $f(1)=1-5+2=-2<0$; $f$ is continuous, so by the Intermediate Value Theorem there is a root in $(0,1)$. \hfill \textbf{[2 marks: B1 for both values, B1 for IVT conclusion]}
\item With $x_{n+1}=\dfrac{x_{n}^{3}+2}{5}$ and $x_{0}=1$:
\[
x_{1}=\frac{1+2}{5}=0.6000,\qquad
x_{2}=\frac{0.6^{3}+2}{5}=\frac{2.216}{5}=0.4432,
\]
\[
x_{3}=\frac{0.4432^{3}+2}{5}=\frac{2.08704\ldots}{5}=0.4174.
\]
\hfill \textbf{[3 marks: M1 correct substitution, A1 all values, A1 to 4 d.p.]}
\item The curves $y=x$ and $y=g(x)=\dfrac{x^{3}+2}{5}$ meet at the root $\alpha\approx0.414$. Starting from $x_{0}=1$ on the $x$-axis, move vertically to $y=g(x)$, then horizontally to $y=x$, and repeat: the staircase spirals inwards towards the intersection.
\begin{popfigure}
\begin{center}
\begin{tikzpicture}[scale=4.2]
\draw[->] (-0.05,0) -- (1.15,0) node[right] {$x$};
\draw[->] (0,-0.05) -- (0,1.15) node[above] {$y$};
\draw[popblue,thick,domain=0:1.12] plot (\x,{\x}) node[right] {$y=x$};
\draw[poporange,thick,domain=0:1.05,samples=60] plot (\x,{(\x*\x*\x+2)/5}) node[above right] {$y=g(x)$};
\draw[popdark,thick] (1,0) -- (1,0.6) -- (0.6,0.6) -- (0.6,0.4432) -- (0.4432,0.4432) -- (0.4432,0.4174) -- (0.4174,0.4174);
\fill (1,0) circle (0.012) node[below] {$x_{0}$};
\fill (0.6,0) circle (0.012) node[below] {$x_{1}$};
\fill (0.4432,0) circle (0.012) node[below left] {$x_{2}$};
\end{tikzpicture}
\end{center}
\legende{Staircase diagram: the iterates converge to the intersection of $y=x$ and $y=g(x)$.}
\end{popfigure}
The iterates $1,\ 0.6,\ 0.4432,\ 0.4174,\ldots$ decrease and converge to the root $\alpha\approx0.4142$; since $|g'(\alpha)|=\frac{3\alpha^{2}}{5}\approx0.103<1$, convergence is guaranteed and monotone from above after $x_{1}$. \hfill \textbf{[3 marks: B1 both curves, B1 staircase, B1 comment on convergence]}
\item For $f(x)=x^{2}-N$, $f'(x)=2x$. Newton--Raphson gives
\[
x_{n+1}=x_{n}-\frac{x_{n}^{2}-N}{2x_{n}}
=\frac{2x_{n}^{2}-x_{n}^{2}+N}{2x_{n}}
=\frac{x_{n}^{2}+N}{2x_{n}}
=\tfrac{1}{2}\left(x_{n}+\frac{N}{x_{n}}\right). \qquad \textbf{[3 marks]}
\]
\item Take $N=11$ and $x_{0}=3$ (since $\sqrt{11}$ is near $3$):
\[
x_{1}=\tfrac12\left(3+\tfrac{11}{3}\right)=\tfrac12\cdot\tfrac{20}{3}=3.333333,
\]
\[
x_{2}=\tfrac12\left(3.333333+\tfrac{11}{3.333333}\right)=\tfrac12(3.333333+3.300000)=3.316667,
\]
\[
x_{3}=\tfrac12\left(3.316667+\tfrac{11}{3.316667}\right)=\tfrac12(3.316667+3.316580)=3.316623,
\]
\[
x_{4}=\tfrac12\left(3.316623+\tfrac{11}{3.316623}\right)=3.316625.
\]
The iterates have stabilised on $5$ decimal places: $\sqrt{11}=\boxed{3.31662}$ to $5$ d.p. \hfill \textbf{[4 marks: M1 formula use, A2 iterates, A1 answer]}
\end{enumerate}
\end{corrige}

\begin{corrige}[Exercise 10 — Failure of Newton--Raphson and error bounds]
\begin{enumerate}
\item Let $f(x)=x^{3}-2x+2$, so $f'(x)=3x^{2}-2$. Starting at $x_{0}=0$:
\[
x_{1}=0-\frac{f(0)}{f'(0)}=-\frac{2}{-2}=1;\qquad
x_{2}=1-\frac{f(1)}{f'(1)}=1-\frac{1}{1}=0.
\]
The iterates \emph{cycle} $0,1,0,1,\ldots$ and never approach the root. Geometrically, the tangent at $x=0$ meets the axis at $x=1$, and the tangent at $x=1$ meets the axis back at $x=0$.
\begin{popfigure}
\begin{center}
\begin{tikzpicture}
\begin{axis}[width=9cm,height=6.5cm,axis lines=middle,xmin=-2.2,xmax=2.2,ymin=-2,ymax=6,xlabel={$x$},ylabel={$y$},xtick={-2,-1,1,2}]
\addplot[domain=-1.9:1.75,thick,popblue,samples=80] {x^3-2*x+2};
\addplot[domain=-0.6:1.3,poporange,thick] {-2*x+2};
\addplot[domain=-0.3:1.9,popgreen,thick] {x};
\fill[popdark] (axis cs:0,2) circle (1.5pt);
\fill[popdark] (axis cs:1,1) circle (1.5pt);
\draw[dashed,popmuted] (axis cs:0,0) -- (axis cs:0,2);
\draw[dashed,popmuted] (axis cs:1,0) -- (axis cs:1,1);
\node[popdark] at (axis cs:1.35,2.2) {tangent at $x=0$};
\node[popdark] at (axis cs:1.5,0.6) {tangent at $x=1$};
\end{axis}
\end{tikzpicture}
\end{center}
\legende{The two tangents send the iterates back and forth between $0$ and $1$: the method cycles.}
\end{popfigure}
\hfill \textbf{[4 marks: B1 computing $x_1=1$, B1 computing $x_2=0$, B1 sketch with tangents, B1 conclusion of cycling]}
\item After $n$ bisections the interval width is $\dfrac{1}{2^{n}}$; taking the midpoint as the estimate, the error satisfies
\[
|\text{error}|\le \frac{1}{2^{n+1}}.
\]
We require $\dfrac{1}{2^{n+1}}<10^{-4}$, i.e.\ $2^{n+1}>10^{4}$. Now $2^{13}=8192<10^{4}$ and $2^{14}=16384>10^{4}$, so $n+1=14$, giving
\[
\boxed{n=13\text{ iterations}.}
\]
(If instead one demands the whole \emph{interval width} to be less than $10^{-4}$, one needs $2^{n}>10^{4}$, i.e.\ $n=14$; both conventions should be credited if justified.) \hfill \textbf{[4 marks: B1 error bound, M1 inequality, A1 powers of 2, A1 answer]}
\item \emph{Note.} With the equation as printed, $t\mathrm{e}^{-t/10}=2$, one finds $f(3)=3\mathrm{e}^{-0.3}-2\approx+0.222$ and $f(4)=4\mathrm{e}^{-0.4}-2\approx+0.681$: both positive, so there is no sign change on $[3,4]$ (indeed the roots of this equation lie in $(2,3)$ and $(20,30)$). The intended equation, consistent with a root between $3$ and $4$, is
\[
t\,\mathrm{e}^{-t/10}=2.5,
\]
which we now solve. Let $f(t)=t\mathrm{e}^{-t/10}-2.5$.
\[
f(3)=3\mathrm{e}^{-0.3}-2.5=2.2225-2.5=-0.2775<0,
\]
\[
f(4)=4\mathrm{e}^{-0.4}-2.5=2.6813-2.5=+0.1813>0.
\]
$f$ is continuous and changes sign, so there is a root in $(3,4)$.

\textbf{Bisection} (error $<0.05$ requires width $<0.1$):
\begin{center}
\begin{tabular}{|c|c|c|c|}
\hline
\rowcolor{popblueL} Interval & Midpoint $m$ & $f(m)$ & New interval \\
\hline
$[3,4]$ & $3.5$ & $3.5\mathrm{e}^{-0.35}-2.5=-0.0336<0$ & $[3.5,4]$ \\
\hline
$[3.5,4]$ & $3.75$ & $3.75\mathrm{e}^{-0.375}-2.5=+0.0773>0$ & $[3.5,3.75]$ \\
\hline
$[3.5,3.75]$ & $3.625$ & $3.625\mathrm{e}^{-0.3625}-2.5=+0.0226>0$ & $[3.5,3.625]$ \\
\hline
$[3.5,3.625]$ & $3.5625$ & $3.5625\mathrm{e}^{-0.35625}-2.5=-0.0055<0$ & $[3.5625,3.625]$ \\
\hline
\end{tabular}
\end{center}
The final interval $[3.5625,3.625]$ has width $0.0625<0.1$; taking the midpoint $t\approx3.59$ gives an error less than $0.03125<0.05$.

\textbf{Interpretation:} the tank empties in approximately $t\approx3.59$ minutes, i.e.\ about $3$ min $35$ s. \hfill \textbf{[7 marks: B1 $f(3),f(4)$ + IVT, M1 bisection set-up, A3 correct halvings, A1 stopping criterion, A1 interpretation]}
\end{enumerate}
\end{corrige}

\begin{corrige}[Exercise 11 — GCE synthesis: the floating buoy]
\begin{enumerate}
\item The buoy has radius $1$ m, so its diameter is $2$ m. The floating depth $d$ is measured from the bottom of the sphere, hence $d\ge0$, and since the density $0.6<1$ the buoy floats partially emerged, so it cannot be fully submerged: $d<2$. Thus $0<d<2$. \hfill \textbf{[2 marks: B1 $d\ge0$, B1 $d<2$ with reason]}
\item Let $f(d)=d^{3}-3d^{2}+2.4$.
\[
f(1)=1-3+2.4=0.4>0,\qquad f(2)=8-12+2.4=-1.6<0,
\]
so by the IVT there is a root in $(1,2)$. Bisect once:
\[
f(1.5)=3.375-6.75+2.4=-0.975<0,
\]
so the root lies in $\boxed{(1,\ 1.5)}$, an interval of width $0.5$. \hfill \textbf{[3 marks: B1 values at 1 and 2, M1 midpoint, A1 interval]}
\item $f'(d)=3d^{2}-6d$, so
\[
d_{n+1}=d_{n}-\frac{d_{n}^{3}-3d_{n}^{2}+2.4}{3d_{n}^{2}-6d_{n}}. \qquad \textbf{[2 marks: B1 derivative, B1 formula]}
\]
\item Start at $d_{0}=1$ (where $f$ is small and positive):
\[
d_{1}=1-\frac{0.4}{-3}=1.133333.
\]
\[
f(1.133333)=1.455704-3.853333+2.4=0.002370,\quad f'(1.133333)=-2.946667,
\]
\[
d_{2}=1.133333-\frac{0.002370}{-2.946667}=1.134137.
\]
\[
f(1.134137)\approx-0.000003,\qquad d_{3}=1.134137-\frac{-0.000003}{-2.943068}=1.134136.
\]
The iterates have stabilised: $d=\boxed{1.134}$ to $3$ decimal places. \hfill \textbf{[5 marks: M1 first iterate, A1 $d_1$, A1 $d_2$, A1 $d_3$, A1 conclusion]}
\item \textbf{Sign-change test} on $(1.1335,\,1.1345)$:
\[
f(1.1335)\approx+0.00147>0,\qquad f(1.1345)\approx-0.00147<0.
\]
The sign change confirms the root in $(1.1335,1.1345)$, so $d=1.134$ m is correct to $3$ d.p. To the nearest centimetre, the buoy floats at a depth of
\[
d\approx1.13\ \mathrm{m}=\boxed{113\ \mathrm{cm}}.
\]
\hfill \textbf{[3 marks: M1 test values, A1 sign change, A1 depth in cm]}
\end{enumerate}
\end{corrige}

\newpage

\coursec{Summary sheet}

\begin{popfigure}
\begin{tikzpicture}[mindmap, grow cyclic, every node/.style={concept}, text=white,
root concept/.append style={concept color=popink, font=\bfseries\small, minimum size=2.9cm},
level 1 concept/.append style={sibling angle=60, level distance=3.6cm, minimum size=2.2cm, font=\scriptsize},
level 2 concept/.append style={sibling angle=45, level distance=2.4cm, minimum size=1.7cm, font=\tiny}]
\node[root concept]{Numerical Methods}
  child[concept color=popblue]{ node{Why?} 
    child[concept color=popblue!70]{ node{No exact formula} }
    child[concept color=popblue!70]{ node{Approx.\ solutions} }
  }
  child[concept color=popgreen]{ node{Locating roots}
    child[concept color=popgreen!70]{ node{Sign change} }
    child[concept color=popgreen!70]{ node{IVT on $[a,b]$} }
  }
  child[concept color=poporange]{ node{Bisection}
    child[concept color=poporange!70]{ node{Halve interval} }
    child[concept color=poporange!70]{ node{Slow, sure} }
  }
  child[concept color=popred]{ node{Newton--Raphson}
    child[concept color=popred!70]{ node{$x_{n+1}=x_n-\frac{f}{f'}$} }
    child[concept color=popred!70]{ node{Fast, may fail} }
  }
  child[concept color=poppurple]{ node{Accuracy}
    child[concept color=poppurple!70]{ node{Error bounds} }
    child[concept color=poppurple!70]{ node{Stopping rules} }
  }
  child[concept color=popgold]{ node{Applications}
    child[concept color=popgold!70]{ node{GCE problems} }
  };
\end{tikzpicture}
\legende{Mind map of the chapter Numerical Methods.}
\end{popfigure}

\begin{bilan}
\textbf{Key definitions}
\begin{itemize}
\item A \cle{root} of $f$ is a value $\alpha$ such that $f(\alpha)=0$; graphically, an $x$-intercept of $y=f(x)$.
\item A \cle{numerical method} is a step-by-step procedure (algorithm) giving an \emph{approximation} of a root when exact algebra fails.
\item An \cle{iteration} is a rule $x_{n+1}=g(x_n)$ producing a sequence $x_0,x_1,x_2,\dots$ that should converge to $\alpha$.
\item The \cle{error} at stage $n$ is $|x_n-\alpha|$; the accuracy required fixes the stopping criterion.
\end{itemize}

\textbf{Theorems and formulas to know}
\begin{itemize}
\item \cle{Intermediate Value Theorem} (sign-change rule): if $f$ is continuous on $[a,b]$ and $f(a)f(b)<0$, then $f$ has \emph{at least one} root in $(a,b)$.
\item \cle{Bisection method}: with $f(a)f(b)<0$, compute the midpoint $m=\dfrac{a+b}{2}$; keep the half-interval whose endpoints have opposite signs. After $n$ steps the error satisfies $|x_n-\alpha|\leq \dfrac{b-a}{2^{n+1}}$.
\item \cle{Newton--Raphson formula}: $x_{n+1}=x_n-\dfrac{f(x_n)}{f'(x_n)}$, obtained from the tangent line at $(x_n,f(x_n))$: $y-f(x_n)=f'(x_n)(x-x_n)$, setting $y=0$.
\item To find $\sqrt[k]{N}$, apply Newton--Raphson to $f(x)=x^k-N$: $x_{n+1}=\dfrac{1}{k}\left[(k-1)x_n+\dfrac{N}{x_n^{k-1}}\right]$.
\end{itemize}

\textbf{Methods (step by step)}
\begin{enumerate}
\item \textbf{Locate the root}: rearrange the equation as $f(x)=0$; tabulate $f$ at consecutive integers; a sign change brackets a root. A sketch of two intersecting curves also locates roots.
\item \textbf{Bisection}: check $f(a)f(b)<0$; compute $m$ and $f(m)$; replace $a$ or $b$ by $m$ keeping the sign change; repeat until the interval is narrow enough, then give the root to the required accuracy.
\item \textbf{Newton--Raphson}: compute $f'$; choose $x_0$ near the root; iterate $x_{n+1}=x_n-f(x_n)/f'(x_n)$; stop when two successive iterates agree to the required number of decimal places.
\item \textbf{Verify}: substitute the final answer into $f$ and check it is close to $0$; state the answer with the requested accuracy.
\end{enumerate}

\textbf{Common mistakes}
\begin{itemize}
\item Claiming a root exists without checking \emph{continuity} of $f$ on $[a,b]$.
\item Forgetting that a sign change guarantees \emph{at least one} root, not exactly one; and that \emph{no} sign change does not prove there is no root (e.g.\ a tangent root like $x^2=0$).
\item In bisection, keeping the wrong half-interval: always keep the sub-interval where the sign change occurs.
\item In Newton--Raphson, starting at $x_0$ where $f'(x_0)=0$ (division by zero) or too far from the root: the method may diverge or jump to another root.
\item Rounding too early: keep at least $2$ extra decimal places during the iterations and round only at the end.
\item Confusing ``correct to $3$ decimal places'' with ``$3$ significant figures''.
\end{itemize}

\textbf{GCE tips}
\begin{itemize}
\item Show the sign-change table with clear values of $f(a)$ and $f(b)$: it earns the first marks.
\item Present iterations in a neat table: $n$, $x_n$, $f(x_n)$, $f'(x_n)$, $x_{n+1}$.
\item Newton--Raphson converges fast (usually $2$--$4$ iterations); bisection is slower but always works once a bracket is found --- the examiners often ask you to \emph{compare} the two.
\item To justify a given iterative formula, derive it: solve $f(x)=0$ via the tangent line, or rearrange $f(x)=0$ into $x=g(x)$.
\item Always finish with a sentence: ``Hence $\alpha = \dots$ correct to $n$ decimal places.''
\end{itemize}
\end{bilan}

\findecours

\end{document}
